%\documentclass[overheadsbeamer.tex]{subfiles}
\documentclass[handoutsbeamer.tex]{subfiles}

\begin{document}

\thispagestyle{empty}

\ona{ \setcounter{chapter}{4} } % ONE LESS
\lecture{Variational Quantum Algorithms}{Lect05}
\chapter{Variational Quantum Algorithms}\label{C.VQE}
\ona{\sloppy}

\ona{This \chap{} is concerned with variational quantum algorithms. We introduce them through the first of them to be proposed, the variational quantum eigensolver (VQE) of \citet{peruzzo2014variational}, drawing on the theory of \citet{mcclean2016theory} and the review of \citet{cerezo2021variational}. After describing the variational loop and its use for molecules, we show how a suitable ansatz arises from the discretisation of an adiabatic evolution; we bound the error of the discretisation and compare, for small hydrogen molecules, the discretised schedule with VQE started from it. We then introduce a second variational algorithm, the quantum approximate optimization algorithm (QAOA) for combinatorial optimization, as far as is needed here, leaving its proper treatment to \Chap~\chref{C.QAOA}, and follow a single run of it in detail, from the circuit to the measured copies and the gradients obtained from shifted circuits. The greater part of the \chap{} is devoted to the question of how hard it is to \emph{train} such an algorithm. Here we follow \citet{bittel2021training}, who formulate the training as a family of computational problems, with a measure of error that distinguishes what the circuit cannot represent from what the optimizer fails to find. They show that the training is $\mathbf{NP}$-hard even when the quantum part of the computation is trivial, namely for circuits without entanglement, on logarithmically many qubits, with a single parameter, for QAOA and for free fermions, and that in the worst case no polynomial-time algorithm can guarantee a non-trivial error. We set out their formulation first and then give the proofs, all but one of them in full. Parts of the \chap{} are adapted from our course on quantum computing \citep{aspman2024quantum}.}

\section{The variational principle}

\begin{frame}\ft{What we need from quantum mechanics}
\begin{itemize}\setlength\itemsep{0.3em}
    \item A \emph{state} of $n$ qubits is a unit vector $\ket\psi \in \C^{2^n}$, written in the computational basis $\{\ket{x} : x \in \{0,1\}^n\}$ as $\ket\psi = \sum_x \psi_x \ket x$; a measurement in this basis returns $x$ with probability $|\psi_x|^2$.
    \item An \emph{observable} is a Hermitian matrix $O$, whose expectation in the state $\ket\psi$ is $\braket{\psi|O|\psi}$. On $n$ qubits every observable is a real combination of Pauli strings, and the expectation of a Pauli string can be estimated by measuring each qubit in a suitable basis (\ona{Section~\ref{sec:ch05:vqa}}\onp{below}).
    \item The evolution is unitary; a gate $e^{-i\theta H}$, with $H$ Hermitian, is a unitary depending on the parameter $\theta$.
\end{itemize}
\begin{thm}[Variational principle]\label{thm:ch05:variational}
Let $H \in \Herm(\C^d)$ have smallest eigenvalue $E_0$. For every non-zero $\ket\psi$, $\braket{\psi|H|\psi}/\braket{\psi|\psi} \ge E_0$, with equality if and only if $\ket\psi$ is a ground state.
\end{thm}
\ona{\emph{Proof.} Expand $\ket\psi = \sum_k C_k\ket{\phi_k}$ in an orthonormal eigenbasis, $H\ket{\phi_k} = E_k\ket{\phi_k}$. Then $\braket{\psi|H|\psi} = \sum_k\big(|C_k|^2E_k\big) \ge E_0\sum_k|C_k|^2 = E_0\braket{\psi|\psi}$, and equality holds if and only if $C_k = 0$ whenever $E_k > E_0$. \qed}
\end{frame}

\section{Variational quantum algorithms}\label{sec:ch05:vqa}

\begin{frame}\ft{The Pauli basis}
\onp{\footnotesize}
\ona{The Hamiltonians and observables of this \chap{} are written in the Pauli basis, whose properties we recall. The Pauli matrices are}
\begin{equation}\label{eq:ch05:pauli}
    I = \begin{pmatrix}1&0\\0&1\end{pmatrix},\qquad X = \begin{pmatrix}0&1\\1&0\end{pmatrix},\qquad Y = \begin{pmatrix}0&-i\\i&0\end{pmatrix},\qquad Z = \begin{pmatrix}1&0\\0&-1\end{pmatrix}.
\end{equation}
\ona{They are Hermitian and unitary, they square to the identity, any two of $X$, $Y$ and $Z$ anticommute, and $XY = iZ$, $YZ = iX$, $ZX = iY$. On $n$ qubits, a \emph{Pauli string} is a tensor product $P = P_1\otimes\cdots\otimes P_n$ with every $P_j\in\{I,X,Y,Z\}$, and its \emph{weight} is the number of factors different from $I$. The $4^n$ Pauli strings are orthogonal with respect to the trace inner product,}
\begin{equation}\label{eq:ch05:paulibasis}
    \Tr(PQ) = 2^n\,\delta_{PQ}, \qquad\text{so that}\qquad M = 2^{-n}\sum_P\Tr(PM)\,P \quad\text{for every } M\in\C^{2^n\times2^n},
\end{equation}
\ona{and they form a basis of the space of all $2^n\times2^n$ complex matrices, the \emph{generalised Pauli basis}. The coefficients $\Tr(PM)$ are real when $M$ is Hermitian. Three further facts will be used. Any two Pauli strings either commute or anticommute. Since $P^2 = \one$, the rotation generated by a Pauli string is $e^{-i\theta P/2} = \cos\frac\theta2\,\one - i\sin\frac\theta2\,P$. And the expectation of $P$ is estimated by rotating each factor $X$ to $Z$ by a Hadamard gate, and each factor $Y$ by $S^\dagger$ followed by a Hadamard gate, measuring in the computational basis, and averaging the parity $(-1)^{\sum_jx_j}$ of the bits on the support of $P$. The size of an instance is then the number of Pauli strings with non-zero coefficients, which is $|E|+1$ for the MaxCut Hamiltonian of a graph and $\mathcal O(N^4)$ for a molecular Hamiltonian on $N$ orbitals. For a space whose dimension $n$ is not a power of two, the clock and shift matrices $Z_n\ket j = \omega^j\ket j$ and $X_n\ket j = \ket{j+1 \bmod n}$, with $\omega = e^{2\pi i/n}$, give $n^2$ unitaries $X_n^aZ_n^b$ that are orthogonal in the same sense, and for $n = 2$ they are the Pauli matrices up to phases; we shall instead embed such a space in the smallest number of qubits that contains it.}
\onp{$4^n$ Pauli strings $P = P_1\otimes\cdots\otimes P_n$: Hermitian, unitary, $P^2 = \one$, commuting or anticommuting, and $e^{-i\theta P/2} = \cos\frac\theta2\one - i\sin\frac\theta2P$. A Hermitian matrix has real coefficients $\Tr(PM)$. To estimate $\braket P$, rotate $X$ and $Y$ factors to $Z$ and average the parity of the measured bits.}
\end{frame}

\begin{frame}\ft{The variational loop}
\onp{\footnotesize}
\onp{\begin{columns}[T]\begin{column}{0.6\textwidth}}
\ona{A \emph{variational quantum algorithm} (VQA) restricts the state to those prepared by a parametrised circuit, $\ket{\Psi(\vartheta)} = U(\vartheta)\ket{\Psi_0}$, and minimises the upper bound $C(\vartheta) = \braket{\Psi(\vartheta)|H|\Psi(\vartheta)}$ with a classical optimizer (Fig.~\ref{fig:ch05:vqa}). Such an algorithm is specified by the following ingredients.}
\begin{itemize}\setlength\itemsep{0.15em}
    \item a dataset $\mathcal D$ and its quantum embedding, here the Hamiltonian $H = \sum_k h_k P_k$ in the Pauli basis;
    \item a parametrised circuit $U(\vartheta) = U_L(\vartheta_L)\cdots U_1(\vartheta_1)$, $U_i = e^{-i\vartheta_iH_i}$, and an initial state $\ket{\Psi_0}$, which together form the \emph{ansatz};
    \item the cost $C(\vartheta) = \sum_k h_k\braket{\Psi(\vartheta)|P_k|\Psi(\vartheta)}$, estimated from repeated measurements up to an additive error of order $1/\sqrt{\text{shots}}$;
    \item a classical optimizer $\mathcal A$ for $\min_\vartheta C(\vartheta)$, either derivative-free (COBYLA, SPSA) or using gradients, which the parameter-shift rule of \Chap~\chref{C.PSR} provides.
\end{itemize}
\onp{\end{column}\begin{column}{0.38\textwidth}\centering
\includegraphics[height=0.74\textheight,width=\linewidth,keepaspectratio]{figures/vqa/vqa2024.pdf}
\end{column}\end{columns}}
\ona{\begin{figure}[htb]\centering\includegraphics[width=0.92\linewidth,height=0.42\textheight,keepaspectratio]{figures/vqa/vqa2024.pdf}\caption{The variational loop. The quantum computer estimates the cost of the parametrised circuit by sampling, and a classical optimizer proposes new parameters.}\label{fig:ch05:vqa}\end{figure}}
\end{frame}

\mode<all>{\gdef\fittab#1#2#3#4#5#6{{\setlength\tabcolsep{4pt}\begin{tabular}{@{}>{\raggedright\arraybackslash}p{#1\linewidth}>{\raggedright\arraybackslash}p{#2\linewidth}>{\raggedright\arraybackslash}p{#3\linewidth}>{\raggedright\arraybackslash}p{#4\linewidth}>{\raggedright\arraybackslash}p{#5\linewidth}>{\raggedright\arraybackslash}p{#6\linewidth}@{}}\toprule
VQA & initial state & generators $H_i$ & observable $O$ & layers $L$ & covered in \\\midrule
hardware-efficient VQE & $\ket0^{\otimes N}$ & $\frac12Y_k$, $\frac12Z_k$ on one qubit, $\frac12Z_kZ_l$ on two & molecular $H = \sum_kh_kP_k$ & number of gates & \ona{Secs.~\ref{sec:ch05:vqa}, \ref{sec:ch05:barren}}\onp{slides \pageref{fr:ch05:vqemol}, \pageref{fr:ch05:barren}} \\
UCCSD, Trotterised & $\ket{\rm HF}$, a basis state & $i(\tau_k - \tau_k^\dagger)$ for an excitation $\tau_k$: a sum of Pauli strings & molecular $H$ & excitations $\times$ Trotter steps & \ona{Secs.~\ref{sec:ch05:vqa}, \ref{sec:ch05:adiabatic}}\onp{slides \pageref{fr:ch05:vqemol}, \pageref{fr:ch05:gates}} \\
QAOA of depth $p$ & $\ket+^{\otimes n}$ & $H_C$ and the mixer $H_M = \sum_jX_j$, alternately & cost $H_C$ & $2p$ & \ona{Secs.~\ref{sec:ch05:qaoa}--\ref{sec:ch05:hardness}}\onp{slides \pageref{fr:ch05:qaoa}--\pageref{fr:ch05:workedlast}}; \Chap~\chref{C.QAOA} \\\bottomrule
\end{tabular}}}}
\begin{frame}\ft{Three variational algorithms in a common form}
\onp{\footnotesize}
\ona{Every ansatz in this \chap{} has the form of the loop of Fig.~\ref{fig:ch05:vqa}, namely a product of layers $U_i(\vartheta_i) = e^{-i\vartheta_iH_i}$, each with a generator $H_i$ and a single real parameter, applied to an initial state, together with an observable whose expectation is to be minimised. Table~\ref{tab:ch05:fit} writes in this form the two ans\"atze for molecules described on the next slide, and the quantum approximate optimization algorithm (QAOA) for combinatorial optimization, which is introduced in Section~\ref{sec:ch05:qaoa}. Two conventions are needed. First, a layer may act on all the qubits at once, and a single parameter may control many gates, as $\gamma$ controls all the $ZZ$ rotations of a layer of QAOA; the generator is then the sum of theirs. Second, the initial state may always be taken to be $\ket0^{\otimes N}$, since a fixed initial state $V\ket0^{\otimes N}$, such as the Hartree--Fock determinant $X^{x}\ket0^{\otimes N}$ or the uniform superposition $\ket+^{\otimes n}$, can be absorbed by conjugation,}
\begin{equation}
    U_L(\vartheta_L)\cdots U_1(\vartheta_1)V\ket0 = V\,e^{-i\vartheta_LV^\dagger H_LV}\cdots e^{-i\vartheta_1V^\dagger H_1V}\ket0,
\end{equation}
\ona{that is, by replacing every $H_i$ by $V^\dagger H_iV$ and $O$ by $V^\dagger OV$. When $V$ is a product of single-qubit Clifford gates, such as $X$ and the Hadamard gate, Pauli strings are mapped to Pauli strings, and the number of terms is unchanged. The training problems of Section~\ref{sec:ch05:formal} adopt this convention.}
\onp{\begin{center}\scriptsize\fittab{0.13}{0.11}{0.24}{0.13}{0.1}{0.14}\end{center}
\footnotesize A fixed initial state $V\ket0$ is absorbed by $H_i\mapsto V^\dagger H_iV$ and $O\mapsto V^\dagger OV$.}
\ona{\begin{table}[htb]\centering\caption{Three variational algorithms in the form of the variational loop, with the sections of this \chap{} in which they are treated. A fixed initial state can be absorbed into the generators and the observable; $\tau_k$ denotes a fermionic excitation operator mapped to qubits.}\label{tab:ch05:fit}\footnotesize\fittab{0.13}{0.11}{0.24}{0.13}{0.1}{0.14}\end{table}}
\ona{Recall that second quantization describes the electrons of a molecule by the occupation numbers of a finite set of spin orbitals, each of which is either empty or holds one electron, so that a configuration is a bit string and a state is a unit vector in $\C^{2^N}$, with one qubit for each of the $N$ spin orbitals; this is the representation introduced in the section \emph{Quantum chemistry for non-chemists} of \Chap~\chref{C.EQA}. An electron in spin orbital $p$ is created by $a_p^\dagger$ and annihilated by $a_p$, operators that satisfy $\{a_p,a_q^\dagger\} = \delta_{pq}$ and $\{a_p,a_q\} = 0$, and the molecular Hamiltonian takes the form $H = \sum_{pq}h_{pq}a_p^\dagger a_q + \frac12\sum_{pqrs}h_{pqrs}a_p^\dagger a_q^\dagger a_ra_s$, with one- and two-electron integrals $h_{pq}$ and $h_{pqrs}$ that are computed classically. The Jordan--Wigner transformation, $a_p = Z_1\cdots Z_{p-1}\,\frac12(X_p + iY_p)$, maps these operators to qubits while preserving the anticommutation relations. It writes $H$, and the excitation operators $\tau_k$ of the table, such as $a_a^\dagger a_i$ and $a_a^\dagger a_b^\dagger a_ja_i$, as sums of Pauli strings, and it gives $a_p^\dagger a_p = \frac12(\one - Z_p)$.}
\end{frame}

\begin{frame}\ft{VQE for molecules}\label{fr:ch05:vqemol}
\ona{The variational quantum eigensolver of \citet{peruzzo2014variational} is the variational algorithm for the electronic-structure problem of \Chap~\chref{C.Motivation}. The Hamiltonian $H$ is the second-quantised molecular Hamiltonian mapped to qubits, with $\mathcal O(N^4)$ Pauli terms for $N$ orbitals, and the aim is its ground-state energy. The principal choices are the following.}
\begin{description}\setlength\itemsep{0.2em}
    \item[Initial state] the Hartree--Fock determinant $\ket{\Psi_0}$, a computational basis state whose energy is already the mean-field energy; the circuit is to add the correlation.
    \item[Chemically inspired ansatz] unitary coupled cluster with single and double excitations (UCCSD), $U(\vartheta) = e^{T(\vartheta) - T(\vartheta)^\dagger}$ with excitation operators $T$, which is exact in a limit but deep in practice.
    \item[Hardware-efficient ansatz] layers of single-qubit rotations and native entangling gates \citep{kandala2017hardware}, which are shallow and expressive but take no account of the problem.
    \item[Measurement] commuting Pauli terms are measured together; the number of shots needed for precision $\epsilon$ grows as $(\sum_k|h_k|)^2/\epsilon^2$ in the worst case.
\end{description}
\ona{On hardware, VQE has given the ground-state energies of small molecules, such as H$_2$, LiH and BeH$_2$, to chemical accuracy after error mitigation \citep{kandala2017hardware}. Its scaling to larger molecules is discussed in \Chap~\chref{C.EQA}.}
\end{frame}

\begin{frame}\ft{VQE as state preparation for early fault-tolerant algorithms}
\ona{VQE carries no guarantee on the quality of the state it returns, and there is evidence that estimating energies by preparing and measuring states, as VQE does, is not competitive for molecules of industrial interest, since the number of shots needed for chemical accuracy is prohibitive. \citet{zhang2022computing} propose a division of labour that retains what VQE does well.}
\begin{enumerate}\setlength\itemsep{0.2em}
    \item VQE, or Hartree--Fock, or a short adiabatic evolution, prepares a state $\ket{\phi_0}$ whose overlap with the ground state is $\eta = |\braket{\phi_0|\psi_0}|^2$.
    \item The state $\ket{\phi_0}$ is passed to a routine for estimating ground-state properties, which uses Hadamard tests of $e^{-ijH}$ and $O$ with classical Fourier post-processing. It returns $\braket{\psi_0|O|\psi_0}$ to accuracy $\epsilon$ with maximal evolution time $\widetilde{\mathcal O}(\gamma^{-1})$, whatever the quality of $\ket{\phi_0}$, at a total cost of $\widetilde{\mathcal O}(\gamma^{-1}\epsilon^{-2}\eta^{-2})$ (\Chap~\chref{C.Motivation}).
\end{enumerate}
\ona{The variational state need not be accurate, provided that its overlap with the ground state is not small. The aim of finding the ground state is thus replaced by the weaker aim of finding a state with non-negligible overlap $\eta$, which relates VQE to the warm starts of \Chaps~\chref{C.WSRounded}--\chref{C.WSContinuous}, where the initial state, too, needs only a good overlap with the optimum. Among the properties that can be estimated in this way are the entries $D_{pq} = \braket{\psi_0|a_p^\dagger a_q|\psi_0}$ of the one-particle reduced density matrix, and hence charge densities and dipole moments; each entry is a combination of four expectations of Pauli strings.}
\onp{VQE supplies the overlap $\eta$, and Fourier post-processing of Hadamard tests supplies the guarantee \citep{zhang2022computing}.}
\end{frame}

\section{VQE as discretised adiabatic computation}\label{sec:ch05:adiabatic}

\begin{frame}\ft{From an adiabatic path to an ansatz}
\onp{\footnotesize}
\ona{It remains to say where good ans\"atze come from. The construction that carries a guarantee is taken from adiabatic computation (\Chap~\chref{C.Annealing}). Let $H_0$ be a Hamiltonian whose ground state $\ket{\Psi_0}$ is easily prepared, let $H_1$ be the Hamiltonian whose ground state is sought, and consider the slow evolution along the path}
\begin{equation}\label{eq:ch05:path}
    H(s) = (1-s)\,H_0 + s\,H_1,\qquad i\frac{d}{dt}\ket{\psi(t)} = H(t/T)\ket{\psi(t)},\quad \ket{\psi(0)} = \ket{\Psi_0},\quad 0\le t\le T.
\end{equation}
\ona{By the adiabatic theorem of \Chap~\chref{C.Annealing}, a total time $T$ of the order of $1/\Delta_{\min}^2$, where $\Delta_{\min}$ is the smallest gap along the path, leaves $\ket{\psi(T)}$ close to the ground state $\ket g$ of $H_1$; we write $\varepsilon_{\rm ad}(T) = 1 - |\braket{g|\psi(T)}|^2$ for the remaining infidelity. For the electronic-structure problem, $H_1 = H$ is the molecular Hamiltonian, and the natural choice of $H_0$ is the \emph{Fock operator} $F = \sum_p\varepsilon_p\,a_p^\dagger a_p$ of the Hartree--Fock method, with the orbital energies $\varepsilon_p$. Among states with the given numbers of electrons, its ground state is the Hartree--Fock determinant $\ket{\rm HF}$, the usual initial state of VQE. A gate-based computer cannot carry out the evolution \eqref{eq:ch05:path} directly, and it is turned into a circuit, and then into VQE, in three steps.}
\begin{enumerate}\setlength\itemsep{0.15em}
    \item The Hamiltonian is made piecewise constant: on the $k$-th of $L$ intervals of length $\Delta t = T/L$, $H(t/T)$ is replaced by $H(s_k)$, with $s_k = (k-\frac12)/L$.
    \item Each factor is split as in the Trotter product formula, $e^{-i\Delta t H(s_k)}$ being replaced by $e^{-i\beta_kH_0}\,e^{-i\gamma_kH_1}$, with $\gamma_k = s_k\Delta t$ and $\beta_k = (1-s_k)\Delta t$.
    \item The angles are left free: the circuit is kept, and an optimizer chooses all $2L$ angles $(\boldsymbol\beta,\boldsymbol\gamma)$ of
\end{enumerate}
\begin{equation}\label{eq:ch05:hva}
    \ket{\Psi(\boldsymbol\beta,\boldsymbol\gamma)} = e^{-i\beta_LH_0}e^{-i\gamma_LH_1}\cdots e^{-i\beta_1H_0}e^{-i\gamma_1H_1}\ket{\Psi_0}.
\end{equation}
\ona{The first factor to act is $e^{-i\gamma_1H_1}$, since $e^{-i\beta_1H_0}$ would merely multiply $\ket{\Psi_0}$ by a phase. The first two steps introduce an error, which we bound below; the third cannot make matters worse, since the schedule is itself a point of the parameter space.}
\end{frame}

\begin{frame}\ft{The discretisation}
\ona{Figure~\ref{fig:ch05:schedule} illustrates the three steps for $L = 6$. The linear schedule $s = t/T$ is replaced by a staircase, and each step of the staircase by two pulses, one of $H_1$ for a time $\gamma_k$ and one of $H_0$ for a time $\beta_k$. The share of $H_1$ grows from layer to layer, as the schedule does; in VQE the lengths of the pulses become free parameters.}
\ona{\begin{figure}[htb]}\begin{center}
\begin{tikzpicture}[x=1.1cm,y=1cm,every node/.style={font=\scriptsize}]
  \draw[-{Latex}] (0,0) -- (6.4,0) node[right] {$t$};
  \draw[-{Latex}] (0,0) -- (0,2.6) node[above] {$s$};
  \draw (0.07,2.2) -- (-0.07,2.2) node[left] {$1$};
  \draw (6,0.07) -- (6,-0.07) node[below] {$T$};
  \draw[thick,gray] (0,0) -- (6,2.2);
  \foreach \k in {1,...,6} {\draw[very thick,blue!70!black] ({\k-1},{2.2*(\k-0.5)/6}) -- (\k,{2.2*(\k-0.5)/6});}
  \draw[densely dotted] (4,0) -- (4,1.65); \draw[densely dotted] (5,0) -- (5,2.017);
  \draw[{Latex}-{Latex}] (4,-0.25) -- node[below] {$\Delta t$} (5,-0.25);
  \begin{scope}[yshift=-1.75cm]
    \foreach \k in {1,...,6} {
      \pgfmathsetmacro{\sk}{(\k-0.5)/6}
      \fill[orange!70] ({\k-1},0) rectangle ({\k-1+\sk},0.5);
      \fill[blue!30] ({\k-1+\sk},0) rectangle (\k,0.5);
      \draw ({\k-1},0) rectangle (\k,0.5);
      \node at ({\k-0.5},-0.22) {$k = \k$};
    }
    \node[anchor=east] at (-0.1,0.25) {pulses};
  \end{scope}
  \node[note,anchor=west] (n1) at (6.6,2.45) {the schedule $s = t/T$};
  \draw[notearrow] (n1.west) to[bend right=15] (4.5,1.68);
  \node[note,anchor=west] (n2) at (6.6,1.55) {step 1: $H(s_k)$ for a time $\Delta t$,\\with $s_k = (k-\frac12)/L$};
  \draw[notearrow] (n2.west) to[bend left=10] (5.6,2.03);
  \node[note,anchor=west] (n3) at (6.6,-1.3) {step 2: $H_1$ for $\gamma_k = s_k\Delta t$ (orange),\\then $H_0$ for $\beta_k = (1-s_k)\Delta t$ (blue)};
  \draw[notearrow] (n3.west) -- (5.4,-1.5);
  \node[note,anchor=north] at (3,-2.25) {in step 3 the $2L$ lengths of the pulses become free parameters,\\and the adiabatic schedule is one point of the parameter space};
\end{tikzpicture}
\end{center}
\ona{\caption{Discretising the adiabatic path with $L = 6$ layers. Top: the linear schedule (grey) and its piecewise-constant approximation at the midpoints (blue). Bottom: each layer as a pulse of $H_1$ followed by a pulse of $H_0$, of lengths $\gamma_k$ and $\beta_k$ with $\gamma_k + \beta_k = \Delta t$.}\label{fig:ch05:schedule}\end{figure}}
\onp{\footnotesize A staircase, then two pulses for each step; VQE leaves the lengths of the pulses free.}
\end{frame}

\begin{frame}\ft{The error of the discretisation}
\onp{\footnotesize}
\begin{lem}\label{lem:ch05:duhamel}
If $H(t)$ and $\tilde H(t)$ generate the unitaries $U(t)$ and $\tilde U(t)$ on $[0,T]$, then $\|U(T) - \tilde U(T)\|\le\int_0^T\|H(t) - \tilde H(t)\|\,dt$.
\end{lem}
\ona{\emph{Proof.} From $U' = -iHU$ and $\tilde U' = -i\tilde H\tilde U$, $\frac d{dt}(U^\dagger\tilde U) = iU^\dagger(H - \tilde H)\tilde U$. Integrate from $0$ to $T$ and use that unitaries preserve norms. \qed}
\begin{lem}[First-order Trotter error]\label{lem:ch05:trotter}
For Hermitian $A$ and $B$, $\|e^{-i(A+B)} - e^{-iA}e^{-iB}\|\le\frac12\|[A,B]\|$.
\end{lem}
\ona{\emph{Proof.} Let $V(t) = e^{iBt}e^{iAt}e^{-i(A+B)t}$, so that the left-hand side is $\|V(1) - \one\|$. Then $V'(t) = ie^{iBt}e^{iAt}\big(e^{-iAt}Be^{iAt} - B\big)e^{-i(A+B)t}$, and $\|e^{-iAt}Be^{iAt} - B\|\le t\|[A,B]\|$, since the derivative of $e^{-iA\tau}Be^{iA\tau}$ is $-ie^{-iA\tau}[A,B]e^{iA\tau}$. Integrating, $\|V(1) - \one\|\le\int_0^1t\,\|[A,B]\|\,dt = \frac12\|[A,B]\|$. \qed}
\ona{We apply the two lemmas to the path \eqref{eq:ch05:path}. In step 1, $\|H(t/T) - H(s_k)\| = |t/T - s_k|\,\|H_1 - H_0\|\le\|H_1-H_0\|/(2L)$ on the $k$-th interval, so that by Lemma~\ref{lem:ch05:duhamel} the error is at most $T\|H_1-H_0\|/(2L)$. In step 2, Lemma~\ref{lem:ch05:trotter} with $A = \gamma_kH_1$ and $B = \beta_kH_0$ bounds the error of the $k$-th layer by $\frac12\beta_k\gamma_k\|[H_0,H_1]\|\le\frac18\Delta t^2\|[H_0,H_1]\|$, since $s_k(1-s_k)\le\frac14$. The errors of a product of unitaries add, and the $L$ layers therefore contribute at most $T^2\|[H_0,H_1]\|/(8L)$. Altogether,}
\begin{equation}\label{eq:ch05:disc}
    \big\|\ket{\psi(T)} - \ket{\Psi(\boldsymbol\beta^{\rm ad},\boldsymbol\gamma^{\rm ad})}\big\| \;\le\; \eta_L \coloneqq \frac{T\,\|H_1 - H_0\|}{2L} + \frac{T^2\,\|[H_0,H_1]\|}{8L},
\end{equation}
\ona{where $(\boldsymbol\beta^{\rm ad},\boldsymbol\gamma^{\rm ad})$ denote the angles of the schedule. Both terms decrease as $1/L$; a splitting of second order would make the second decrease as $1/L^2$.}
\onp{\footnotesize Both terms decrease as $1/L$; if $H_0$ and $H_1$ commute, step 2 is exact.}
\end{frame}

\begin{frame}\ft{The ansatz contains the discretised evolution}
\onp{\scriptsize}
\begin{propn}\label{propn:ch05:adiabatic}
At the angles of the schedule, the state \eqref{eq:ch05:hva} has approximation error, relative to the spectral width of $H_1$, at most
\begin{equation*}
    \delta \;\le\; \varepsilon_{\rm ad}(T) + \eta_L.
\end{equation*}
Hence the model mismatch $\delta_m$ of the ansatz with $L$ layers obeys the same bound, and so does the result of any optimizer started at these angles that never increases the energy.
\end{propn}
\ona{\emph{Proof.} For any state $\ket\phi$ with $|\braket{g|\phi}|^2 = 1-\epsilon$, expanding in eigenvectors gives $\braket{\phi|H_1|\phi} - E_0\le\epsilon\,w(H_1)$, so $\ket{\psi(T)}$ has $\delta\le\varepsilon_{\rm ad}(T)$. For unit vectors $\phi,\phi'$ and any $c\in\R$, $|\braket{\phi|H_1|\phi} - \braket{\phi'|H_1|\phi'}|\le2\|H_1 - c\one\|\,\|\phi-\phi'\|$; with $c$ the midpoint of the spectrum, $2\|H_1 - c\one\| = w(H_1)$, and \eqref{eq:ch05:disc} adds at most $\eta_L$ to $\delta$. \qed}
\begin{itemize}\setlength\itemsep{0.2em}
    \item If the gap along the path is at least inverse-polynomial and the norms are polynomial, then $T$ and $L = \mathcal O\big((T\|H_1-H_0\| + T^2\|[H_0,H_1]\|)/\epsilon\big)$ are polynomial, and the ansatz contains a state with $\delta\le\epsilon$. Since adiabatic computation is universal \citep{aharonov2007adiabatic}, so is this family of ans\"atze.
    \item The schedule supplies angles with a known error, from which an optimizer can only improve. This is a warm start in the sense of \Chaps~\chref{C.WSRounded}--\chref{C.WSContinuous}; its guarantee is not affected by the persistent local minima of \citet{bittel2021training}, although it does not promise the global optimum.
    \item Combinatorial optimization corresponds to a diagonal $H_1$, whose ground states are the optimal solutions, with $H_0 = -\sum_jX_j$; it is taken up in \ona{Section~\ref{sec:ch05:qaoa}}\onp{a later section}.
\end{itemize}
\end{frame}

\begin{frame}\ft{The gates for a molecule}\label{fr:ch05:gates}
\onp{\footnotesize}
\ona{Under the Jordan--Wigner transformation of Section~\ref{sec:ch05:vqa}, $a_p^\dagger a_p = \frac12(\one - Z_p)$, and the two kinds of layer in \eqref{eq:ch05:hva} differ considerably in cost.}
\begin{itemize}\setlength\itemsep{0.2em}
    \item The Fock operator is $F = \mathrm{const} - \frac12\sum_p\varepsilon_pZ_p$, so that $e^{-i\beta F}$ is, up to a phase, a product of single-qubit rotations $R_Z(-\beta\varepsilon_p)$, and the Fock layers are cheap.
    \item The molecular Hamiltonian $H = \sum_\alpha h_\alpha P_\alpha$ has $\mathcal O(N^4)$ Pauli strings for $N$ orbitals, which do not commute. A further Trotter splitting, $e^{-i\gamma H}\approx\prod_\alpha e^{-i\gamma h_\alpha P_\alpha}$, turns each factor into a rotation generated by a Pauli string, which is implemented as a ladder of CNOTs around a single $R_Z$.
    \item If each group of commuting terms is given its own angle in each layer, one obtains the \emph{Hamiltonian variational ansatz} of \citet{Wecker2015}, which has more parameters and is more expressive, and is correspondingly harder to optimize.
    \item For H$_2$ in the minimal basis, the states reachable from $\ket{\rm HF}$ span only two dimensions, and a single layer with free angles, $e^{-i\beta F}e^{-i\gamma H}\ket{\rm HF}$, already attains the exact ground energy, whereas the schedule itself, with $T = 16$, needs $L = 9$ layers for chemical accuracy, $1.6$ mHa.
\end{itemize}
\ona{All the numbers in this section were obtained with \texttt{code/ch05\_adiabatic.py}, which constructs the qubit Hamiltonians from the integrals computed by PySCF, as in \Chap~\chref{C.Motivation}, checks their ground energies against the full configuration interaction of PySCF, and works in the sector with the given numbers of spin-up and spin-down electrons, which $F$ and $H$ both preserve.}
\end{frame}

\begin{frame}\ft{The duration of the evolution}
\ona{Figure~\ref{fig:ch05:adiabaticT} shows the error of the continuous evolution \eqref{eq:ch05:path} as a function of $T$, in units of $\hbar$ per Hartree, for H$_2$ at $0.75$~\angstrom{} and for a linear chain of four hydrogen atoms, H$_4$, at spacings of $1.0$ and $1.5$~\angstrom{}, on eight qubits. Along the paths, the gap within the singlet states decreases to $0.52$~Hartree at $1.0$~\angstrom{} and to $0.17$~Hartree at $1.5$~\angstrom{}, in both cases at $s = 1$, and the stretched chain converges correspondingly more slowly. The errors decrease roughly as $1/T^2$, which is the rate for a linear schedule.}
\ona{\gdef\atH{0.4\linewidth}}\onp{\gdef\atH{0.34\linewidth}}
\ona{\begin{figure}[htb]}\begin{center}
\begin{tikzpicture}
\begin{axis}[width=0.5\linewidth,height=\atH,clip=false,xmode=log,ymode=log,log basis x=2,xmin=1,xmax=64,ymin=3e-6,ymax=0.3,xtick={1,2,4,8,16,32,64},xticklabels={1,2,4,8,16,32,64},xlabel={total time $T$},ylabel={energy error (Hartree)},thick,mark size=1.6pt]
  \addplot[green!45!black,mark=*] table[x=T,y=m0] {figures/ch05/adiabatic_T.dat};
  \addplot[blue!70!black,mark=*] table[x=T,y=m1] {figures/ch05/adiabatic_T.dat};
  \addplot[red!70!black,mark=*] table[x=T,y=m2] {figures/ch05/adiabatic_T.dat};
  \addplot[gray,densely dotted,very thick,domain=1:64,samples=2] {1.6e-3};
  \node[note,anchor=west] at (axis cs:72,4.0e-4) {H$_4$, $1.5$~\angstrom};
  \node[note,anchor=west] at (axis cs:72,5.8e-5) {H$_4$, $1.0$~\angstrom};
  \node[note,anchor=west] at (axis cs:72,6.8e-6) {H$_2$, $0.75$~\angstrom};
  \node[note,anchor=west] at (axis cs:72,1.6e-3) {chemical accuracy};
  \node[note,anchor=south west] (c2) at (axis cs:1.1,1.2e-5) {a smaller gap requires\\a longer evolution};
  \draw[notearrow] (c2.north) to[bend left=20] (axis cs:8,1.6e-2);
\end{axis}
\end{tikzpicture}
\end{center}
\ona{\caption{The adiabatic evolution \eqref{eq:ch05:path} from the Fock operator to the molecular Hamiltonian, in the minimal basis: the energy error at time $T$ above the exact ground energy. Of the times computed, chemical accuracy, $1.6$~mHa (dotted), is first reached at $T = 4$ for H$_2$, $T = 16$ for H$_4$ at $1.0$~\angstrom{} and $T = 64$ at $1.5$~\angstrom.}\label{fig:ch05:adiabaticT}\end{figure}}
\onp{\footnotesize The errors decrease as $1/T^2$; the stretched chain, which has the smallest gap, needs the longest evolution.}
\end{frame}

\begin{frame}\ft{The schedule and VQE for H$_4$ at $1.0$~\angstrom}
\ona{Figures~\ref{fig:ch05:adiabaticL} and~\ref{fig:ch05:adiabaticL2} fix $T = 16$ and vary the number of layers $L$ for the two chains of four atoms. The blue curves show the discretised schedule, with the angles of steps 1 and 2, and the red curves show VQE, here a quasi-Newton algorithm, such as Broyden--Fletcher--Goldfarb--Shanno (BFGS), which is the one we use, applied to all $2L$ angles and started from the schedule. The schedule is accurate only when $\Delta t = T/L$ is small. VQE reaches chemical accuracy from $L = 6$ onwards, but its error is not monotone in $L$: with five layers it does worse than with four, because BFGS stops at a local minimum.}
\ona{\gdef\alH{0.42\linewidth}\gdef\alW{0.56\linewidth}}\onp{\gdef\alH{0.4\linewidth}\gdef\alW{0.52\linewidth}}
\ona{\begin{figure}[htb]}\begin{center}
\begin{tikzpicture}
\begin{axis}[name=al,width=\alW,height=\alH,clip=false,ymode=log,xmin=0.5,xmax=12.5,ymin=1e-10,ymax=3,xtick={1,2,4,6,8,10,12},ytick={1e-10,1e-8,1e-6,1e-4,1e-2,1},xlabel={layers $L$},ylabel={energy error (Hartree)},thick,mark size=1.6pt,title={\scriptsize H$_4$, $1.0$~\angstrom}]
  \addplot[blue!70!black,mark=o] table[x=L,y=disc1] {figures/ch05/adiabatic_L.dat};
  \addplot[red!70!black,mark=*] table[x=L,y=vqe1] {figures/ch05/adiabatic_L.dat};
  \addplot[blue!70!black,dashed,domain=0.5:12.5,samples=2] {8.1e-4};
  \addplot[gray,densely dotted,very thick,domain=0.5:12.5,samples=2] {1.6e-3};
  \node[note,anchor=west] (a1) at (axis cs:13.3,1.5) {blue: the schedule, accurate only\\once $\Delta t = T/L$ is small};
  \draw[notearrow] (a1.west) to[bend right=10] (axis cs:10.2,1.2e-3);
  \node[note,anchor=west] (a2) at (axis cs:13.3,1e-2) {a local minimum\\at $L = 5$};
  \draw[notearrow] (a2.west) to[bend right=15] (axis cs:5.2,3.4e-2);
  \node[note,anchor=west] (a3) at (axis cs:13.3,1e-5) {red: VQE from the schedule's\\angles, chemically accurate\\from $L = 6$};
  \draw[notearrow] (a3.west) to[bend left=10] (axis cs:6.2,1.0e-4);
  \node[note,anchor=west] (a4) at (axis cs:13.3,1e-8) {dashed: the continuous\\evolution; dotted: chemical\\accuracy, $1.6$ mHa};
  \draw[notearrow] (a4.west) to[bend left=15] (axis cs:12.4,8.1e-4);
\end{axis}
\end{tikzpicture}
\end{center}
\ona{\caption{The discretised adiabatic schedule (blue) and VQE started from it (red), for H$_4$ at spacing $1.0$~\angstrom{} with $T = 16$: energy errors above the exact ground energy as functions of the number of layers $L$. Dashed: the continuous evolution for the same $T$. Dotted: chemical accuracy.}\label{fig:ch05:adiabaticL}\end{figure}}
\onp{\footnotesize The schedule needs a small $\Delta t$; VQE started from its angles needs few layers, but may stop at a local minimum.}
\end{frame}

\begin{frame}\ft{The schedule and VQE for H$_4$ at $1.5$~\angstrom}
\ona{At $1.5$~\angstrom{} the gap is smaller and $T = 16$ is too short, so that the schedule stalls near the error of the continuous evolution, above chemical accuracy, however many layers it is given. VQE started from it reaches chemical accuracy from $L = 4$ onwards, and errors below $10^{-6}$ Hartree from $L = 10$ onwards; the freedom in the angles compensates for the short evolution.}
\ona{\gdef\alH{0.42\linewidth}\gdef\alW{0.56\linewidth}}\onp{\gdef\alH{0.4\linewidth}\gdef\alW{0.52\linewidth}}
\ona{\begin{figure}[htb]}\begin{center}
\begin{tikzpicture}
\begin{axis}[name=al,width=\alW,height=\alH,clip=false,ymode=log,xmin=0.5,xmax=12.5,ymin=1e-10,ymax=3,xtick={1,2,4,6,8,10,12},ytick={1e-10,1e-8,1e-6,1e-4,1e-2,1},xlabel={layers $L$},ylabel={energy error (Hartree)},thick,mark size=1.6pt,title={\scriptsize H$_4$, $1.5$~\angstrom}]
  \addplot[blue!70!black,mark=o] table[x=L,y=disc2] {figures/ch05/adiabatic_L.dat};
  \addplot[red!70!black,mark=*] table[x=L,y=vqe2] {figures/ch05/adiabatic_L.dat};
  \addplot[blue!70!black,dashed,domain=0.5:12.5,samples=2] {5.6e-3};
  \addplot[gray,densely dotted,very thick,domain=0.5:12.5,samples=2] {1.6e-3};
  \node[note,anchor=west] (b1) at (axis cs:13.3,1.5) {blue: the schedule stalls near\\the continuous error (dashed),\\above chemical accuracy (dotted)};
  \draw[notearrow] (b1.west) to[bend right=10] (axis cs:10.2,4.2e-3);
  \node[note,anchor=west] (b2) at (axis cs:13.3,1e-4) {red: VQE from the schedule's\\angles, chemically accurate\\from $L = 4$};
  \draw[notearrow] (b2.west) to[bend left=10] (axis cs:4.2,8.5e-4);
  \node[note,anchor=west] (b3) at (axis cs:13.3,1e-8) {below $10^{-6}$ Hartree\\from $L = 10$};
  \draw[notearrow] (b3.west) to[bend left=10] (axis cs:10.2,9e-7);
\end{axis}
\end{tikzpicture}
\end{center}
\ona{\caption{As Fig.~\ref{fig:ch05:adiabaticL}, for H$_4$ at spacing $1.5$~\angstrom.}\label{fig:ch05:adiabaticL2}\end{figure}}
\onp{\footnotesize For this gap $T$ is too short, and the schedule stalls; VQE does not.}
\end{frame}

\begin{frame}\ft{The scope of the adiabatic argument}
\onp{\footnotesize}
\ona{The adiabatic argument is the clearest justification of VQE available, and its scope should be stated precisely.}
\begin{itemize}\setlength\itemsep{0.25em}
    \item It provides an ansatz whose model mismatch is bounded in terms of the gap, the norms and the commutator (\ona{Proposition~\ref{propn:ch05:adiabatic}}\onp{the proposition}), a starting point with a certified error, and universality.
    \item It does not provide the gap. For molecules the gap along the path from the Fock operator is often large, but it decreases as bonds are stretched\ona{, as in Fig.~\ref{fig:ch05:adiabaticT}}; for hard optimization problems it is exponentially small (\Chap~\chref{C.Annealing}), and so are the guarantees.
    \item It does not make training easy. The freedom in the angles is what allows VQE to improve on the schedule with few layers, but it also exposes the optimizer to the landscapes of \citet{bittel2021training} studied below; \ona{Fig.~\ref{fig:ch05:adiabaticL}}\onp{the example of H$_4$} already exhibits a local minimum.
    \item It does not account for measurements. The bounds concern states, and each energy must still be estimated from measured copies, as the worked example below shows, with the numbers of shots discussed in \Chap~\chref{C.PerIteration}.
\end{itemize}
\ona{The same argument underlies the algorithm for combinatorial optimization of Section~\ref{sec:ch05:qaoa} and its convergence as the depth grows (\Chap~\chref{C.QAOA}), and also the warm starts of \Chaps~\chref{C.WSRounded}--\chref{C.WSContinuous}, which replace the adiabatic starting point by one obtained from a classical algorithm.}
\end{frame}

\section{What kind of algorithm is a VQA?}

\begin{frame}\ft{VQAs as local search algorithms}
A VQA may be regarded as a local search algorithm. Local search starts from an initial solution $s_0$, often chosen at random, and sets $s^* = s_0$. Until a stopping criterion is met, it perturbs the current solution slightly, and if the perturbed solution $s'$ satisfies $f(s') < f(s^*)$, it sets $s^* = s'$. It returns the last stored $s^*$.

A VQA proceeds in the same way. It starts from initial parameters $\vartheta_0$, often chosen at random, and sets $\vartheta^* = \vartheta_0$. Until a stopping criterion, usually a number of iterations, is met, the classical optimizer proposes new parameters $\vartheta'$, and if $C(\vartheta') < C(\vartheta^*)$, it sets $\vartheta^* = \vartheta'$. It returns the last stored $\vartheta^*$.

\ona{Local search has a complexity class of its own. \citet{johnson1988easy} defined the class $\mathbf{PLS}$, for \emph{polynomial local search}, of search problems in which an initial solution, the value of any solution and, for any solution that is not locally optimal, a better neighbour can all be computed in polynomial time; the task is to find a local optimum. Since a local optimum always exists, $\mathbf{PLS}$ lies between $\mathbf{FP}$ and the class $\mathbf{TFNP}$ of total search problems. Finding a cut of a weighted graph that cannot be improved by moving a single vertex is complete for $\mathbf{PLS}$; for unweighted graphs, such as those of Section~\ref{sec:ch05:hardness}, every improving move gains an edge, and greedy search stops after at most $|E|$ moves. In general, following improving moves may take exponentially many steps, and it is not known whether a local optimum can always be found in polynomial time by other means: whether $\mathbf{PLS} = \mathbf{FP}$ is open, and no problem in $\mathbf{PLS}$ is $\mathbf{NP}$-hard unless $\mathbf{NP} = \mathbf{coNP}$ \citep{johnson1988easy}. The complexity of local search, and with it that of finding the local minima at which VQAs stop, is thus largely open.}
\onp{\footnotesize Local search is captured by the class $\mathbf{PLS}$ \citep{johnson1988easy}; whether a local optimum can always be found in polynomial time is open.}
\end{frame}

\begin{frame}\ft{VQAs are not approximation algorithms}
On the other hand, VQAs are not \emph{approximation algorithms} in the usual sense. An approximation algorithm for a problem is a polynomial-time algorithm that returns, on \emph{every} instance, a solution within a guaranteed factor of the optimum. The Goemans--Williamson algorithm for MaxCut, for example, guarantees the ratio $0.878$ on every instance. No such guarantee is known for VQAs in general.

\ona{To turn such statements into theorems, one must specify which computational problem the classical optimizer is asked to solve, and how its failure is to be measured. This is what \citet{bittel2021training} do, and we follow them in Section~\ref{sec:ch05:formal}, after a second variational algorithm and a worked example.}
\end{frame}

\section{The quantum approximate optimization algorithm}\label{sec:ch05:qaoa}

\begin{frame}\ft{QAOA, a second variational algorithm}\label{fr:ch05:qaoa}
\onp{\footnotesize}
\onp{The \emph{quantum approximate optimization algorithm} (QAOA) \citep{farhi2014quantum} is the VQA for combinatorial optimization, treated in full in \Chap~\chref{C.QAOA}. The objective $f$ to be maximised is encoded in a diagonal cost Hamiltonian,}
\ona{VQE is not the only variational algorithm. The \emph{quantum approximate optimization algorithm} (QAOA) of \citet{farhi2014quantum} is the corresponding algorithm for combinatorial optimization. It is treated in full in \Chap~\chref{C.QAOA}, and here we introduce only what the worked example below requires. Let $f:\{0,1\}^n\to\R$ be an objective to be maximised, such as the size of a cut in a graph, and encode it in a diagonal \emph{cost Hamiltonian},}
\begin{equation}\label{eq:ch05:HC}
    H_C\ket x = -f(x)\ket x\quad\text{for all } x\in\{0,1\}^n;\qquad\text{for MaxCut,}\quad H_C = -\tfrac12\sum_{ij\in E}(\one - Z_iZ_j),
\end{equation}
\ona{whose ground states are the optimal bit strings. For MaxCut, $\braket{x|H_C|x} = -\mathrm{cut}(x)$, since $Z_iZ_j$ takes the value $-1$ on $\ket x$ exactly when the edge $ij$ is cut. QAOA of depth $p$ is the VQA with the following ingredients:}
\begin{itemize}\setlength\itemsep{0.15em}
    \item the initial state $\ket+^{\otimes n}$, the uniform superposition of all bit strings, prepared by a Hadamard gate on each qubit;
    \item the ansatz of $p$ alternating layers of the \emph{cost unitary} $e^{-i\gamma_kH_C}$ and the \emph{mixer} $e^{-i\beta_kH_M}$, where $H_M = \sum_jX_j$,
    \begin{equation}\label{eq:ch05:qaoa}
        \ket{\psi_p(\boldsymbol\beta,\boldsymbol\gamma)} = e^{-i\beta_pH_M}e^{-i\gamma_pH_C}\cdots e^{-i\beta_1H_M}e^{-i\gamma_1H_C}\ket+^{\otimes n};
    \end{equation}
    \ona{The name \emph{mixer} describes the purpose of this layer. The cost unitary is diagonal in the computational basis; it multiplies each bit string $\ket x$ by the phase $e^{i\gamma f(x)}$ and leaves the probabilities of measuring the bit strings unchanged, so that layers of cost unitaries alone would never change the uniform distribution. The mixer $e^{-i\beta H_M} = \prod_jR_X(2\beta)$ is not diagonal. Since $H_M$ is the adjacency matrix of the hypercube of bit strings, the mixer moves amplitude between bit strings that differ in a single bit and thereby mixes them, so that the phases introduced by the cost layer interfere and amplitude can accumulate on good bit strings. It plays the role of the transverse field that drives the transitions in quantum annealing (\Chap~\chref{C.Annealing}).}
    \onp{The cost layer only introduces phases $e^{i\gamma f(x)}$; the mixer moves amplitude between bit strings that differ in one bit, so that the phases interfere.}
    \item the cost $F_p(\boldsymbol\beta,\boldsymbol\gamma) = \braket{\psi_p|H_C|\psi_p}$, which is minus the expected objective of a measured bit string, minimised over the $2p$ angles by a classical optimizer\ona{, as in Fig.~\ref{fig:ch05:vqa}}.
\end{itemize}
\end{frame}

\begin{frame}\ft{QAOA as a special case of a VQA}
\onp{\footnotesize}
\ona{In terms of the variational loop of Fig.~\ref{fig:ch05:vqa}, QAOA fixes every ingredient other than the angles.}
\begin{description}\setlength\itemsep{0.15em}
    \item[Embedding] The problem enters only through the diagonal Hamiltonian $H_C$.
    \item[Ansatz] There are two generators, $H_C$ and $H_M$, used alternately, and hence $2p$ layers and $2p$ parameters, whatever the number of qubits. The cost unitary is a product of commuting rotations, one for each term of $H_C$, which for MaxCut means one $ZZ$ rotation for each edge; the mixer is a product of single-qubit $X$ rotations.
    \item[Measurement] Every qubit is measured in the computational basis, since $H_C$ is diagonal. Each measured copy yields a bit string $x$ and its objective $f(x)$, so that the same samples serve to estimate the energy and as candidate solutions; the output is the best bit string observed.
    \item[Optimizer] As for VQE, the optimizer is either derivative-free or based on gradients obtained by the parameter-shift rule (\Chap~\chref{C.PSR}).
\end{description}
\ona{Its row in Table~\ref{tab:ch05:fit} expresses it in the form of the loop. QAOA differs from VQE in two respects. Its ansatz is the discretised adiabatic evolution of Section~\ref{sec:ch05:adiabatic}, with $H_0 = -H_M$, whose ground state is $\ket+^{\otimes n}$, and $H_1 = H_C$, up to the sign of $\boldsymbol\beta$ \citep{farhi2014quantum}, so that Proposition~\ref{propn:ch05:adiabatic} applies, with the gap of the annealing path of \Chap~\chref{C.Annealing}. And its object is a bit string rather than an energy, so that what matters is how often good bit strings are sampled. Unlike VQAs in general, QAOA comes with a few guarantees: for MaxCut on 3-regular graphs, depth one achieves the ratio $0.692$ \citep{farhi2014quantum}, and warm-started QAOA inherits the ratio of the relaxation from which it starts \citep{egger2021warm}. Its performance at fixed depth, and as the depth grows, is studied in \Chap~\chref{C.QAOA}, and warm starts in \Chap~\chref{C.WSRounded}.}
\onp{Its ansatz is the discretised annealing of the adiabatic section, from $-H_M$ to $H_C$, and its output is a bit string. Its performance and convergence are studied in \Chap~\chref{C.QAOA}.}
\end{frame}

\section{A worked example with QAOA of depth one}\label{sec:ch05:worked}

\begin{frame}\ft{The instance and the ansatz}
\onp{\scriptsize}
\onp{\begin{columns}[T]\begin{column}{0.66\textwidth}}
\ona{To illustrate every step of the loop of Fig.~\ref{fig:ch05:vqa}, we carry out a few iterations on the smallest instance of interest, MaxCut on the 4-cycle with vertices $1,2,3,4$ and edges $12$, $23$, $34$, $41$, with QAOA of depth one, $p = 1$. In the notation of \eqref{eq:ch05:HC} and \eqref{eq:ch05:qaoa},}
\begin{gather}
    H_C = -\tfrac12\sum_{ij\in E}(\one - Z_iZ_j),\qquad H_M = \sum_{j=1}^4 X_j,\nonumber\\
    \ket{\psi(\beta,\gamma)} = e^{-i\beta H_M}e^{-i\gamma H_C}\ket{+}^{\otimes4}.
\end{gather}
\ona{The cost $F(\beta,\gamma) = \braket{\psi(\beta,\gamma)|H_C|\psi(\beta,\gamma)}$ is minus the expected cut of a measured bit string, and a short computation (Exercise~\ref{exe:ch05:c4}) gives}
\begin{equation}\label{eq:ch05:c4}
    F(\beta,\gamma) = -2 + \sin4\beta\,\sin2\gamma,
\end{equation}
\ona{with minimum $-3$ at $(\beta,\gamma) = (\pi/8,-\pi/4)$; at depth one, therefore, the expected cut is at most $3$ of the $4$ edges. The optimizer is not given \eqref{eq:ch05:c4}, and has only estimates obtained from measurements. All the numbers below were obtained with \texttt{code/ch05\_worked\_example.py}, which simulates the circuit gate by gate and samples its measurements.}
\begin{itemize}\setlength\itemsep{0.1em}
    \item two parameters $(\beta,\gamma)$, starting at $(0.1,-0.1)$;
    \item $M = 1000$ measured copies for each circuit;
    \item gradients obtained from shifted circuits, and steps $(\beta,\gamma)\leftarrow(\beta,\gamma) - 0.1\,\hat g$.
\end{itemize}
\onp{\end{column}\begin{column}{0.32\textwidth}}
\begin{center}
\begin{tikzpicture}[every node/.style={font=\scriptsize},lv/.style={circle,fill=black,inner sep=2pt},rv/.style={circle,draw,thick,fill=white,inner sep=1.8pt}]
  \node[lv,label={above left:$1$}] (v1) at (0,1.4) {};
  \node[rv,label={above right:$2$}] (v2) at (1.4,1.4) {};
  \node[lv,label={below right:$3$}] (v3) at (1.4,0) {};
  \node[rv,label={below left:$4$}] (v4) at (0,0) {};
  \foreach \u/\v in {v1/v2,v2/v3,v3/v4,v4/v1} \draw[green!45!black,thick] (\u) -- (\v);
  \node[note,anchor=north] at (0.7,-0.45) {the maximum cut $\{1,3\}\,|\,\{2,4\}$,\\given by the bit strings $1010$\\and $0101$, cuts all four edges};
\end{tikzpicture}
\end{center}
\onp{\end{column}\end{columns}}
\end{frame}

% CNOT (x, control y, target y) and measurement (x, y), for the circuit of the worked example
\mode<all>{\gdef\wecnot#1#2#3{\draw[thick] (#1,#2) -- (#1,#3); \fill (#1,#2) circle (2.2pt); \draw[thick,fill=white] (#1,#3) circle (0.14); \draw[thick] (#1,#3-0.14) -- (#1,#3+0.14); \draw[thick] (#1-0.14,#3) -- (#1+0.14,#3);}\gdef\wemeter#1#2{\draw[fill=white] (#1-0.26,#2-0.22) rectangle (#1+0.26,#2+0.22); \draw (#1-0.15,#2-0.08) arc (150:30:0.17); \draw[-{Stealth[length=2.5pt]}] (#1,#2-0.12) -- (#1+0.13,#2+0.14);}}
\begin{frame}\ft{The circuit}
\onp{\footnotesize}
\ona{The terms of $H_C$ commute, and so, up to a global phase, $e^{-i\gamma H_C} = \prod_{ij\in E}R_{ZZ}(\gamma)$, where $R_{ZZ}(\gamma) = e^{-i\gamma Z_iZ_j/2}$ acts on qubits $i$ and $j$. A $ZZ$ rotation is not a native gate on most hardware, and is built from controlled gates. Since conjugation by a CNOT maps $\one\otimes Z$ to $Z\otimes Z$,}
\begin{equation}\label{eq:ch05:rzz}
    R_{ZZ}(\gamma) = \mathrm{CNOT}\,\big(\one\otimes R_Z(\gamma)\big)\,\mathrm{CNOT} = e^{i\gamma/2}\big(R_Z(\gamma)\otimes R_Z(\gamma)\big)\,\mathrm{CP}(-2\gamma),
\end{equation}
\ona{where $R_Z(\gamma) = e^{-i\gamma Z/2}$ and $\mathrm{CP}(\varphi) = \mathrm{diag}(1,1,1,e^{i\varphi})$ is the controlled phase rotation. A $ZZ$ rotation is thus either two CNOTs around a $Z$ rotation of the target, or one controlled rotation and two single-qubit rotations (Exercise~\ref{exe:ch05:rzz}). The mixer is $e^{-i\beta H_M} = \prod_jR_X(2\beta)$, with $R_X(\theta) = e^{-i\theta X/2}$. The whole circuit, shown in Fig.~\ref{fig:ch05:circuit}, consists of four Hadamard gates, eight CNOTs, four rotations $R_Z(\gamma)$ and four rotations $R_X(2\beta)$, followed by a measurement of every qubit.}
\ona{\begin{figure}[htb]}\begin{center}
\begin{tikzpicture}[every node/.style={font=\scriptsize},g/.style={draw,fill=white,minimum height=0.48cm,minimum width=0.5cm,inner sep=1.5pt}]
  \foreach \q/\y in {1/0,2/-0.75,3/-1.5,4/-2.25} {\draw (0,\y) node[left] {$\ket0$} -- (8.75,\y); \node[g] at (0.5,\y) {$H$}; \node[g] at (7.9,\y) {$R_X(2\beta)$};}
  \wecnot{1.3}{0}{-0.75} \wecnot{1.3}{-1.5}{-2.25}
  \node[g] at (2.0,-0.75) {$R_Z(\gamma)$}; \node[g] at (2.0,-2.25) {$R_Z(\gamma)$};
  \wecnot{2.7}{0}{-0.75} \wecnot{2.7}{-1.5}{-2.25}
  \wecnot{3.5}{-0.75}{-1.5} \node[g] at (4.2,-1.5) {$R_Z(\gamma)$}; \wecnot{4.9}{-0.75}{-1.5}
  \wecnot{5.6}{-2.25}{0} \node[g] at (6.3,0) {$R_Z(\gamma)$}; \wecnot{7.0}{-2.25}{0}
  \foreach \y in {0,-0.75,-1.5,-2.25} {\wemeter{8.85}{\y}}
  \draw[decorate,decoration={brace,amplitude=4pt}] (1.1,0.35) -- (2.9,0.35) node[midway,above=4pt] {edges $12$, $34$};
  \draw[decorate,decoration={brace,amplitude=4pt}] (3.3,0.35) -- (5.1,0.35) node[midway,above=4pt] {edge $23$};
  \draw[decorate,decoration={brace,amplitude=4pt}] (5.4,0.35) -- (7.2,0.35) node[midway,above=4pt] {edge $41$};
  \draw[decorate,decoration={brace,amplitude=4pt}] (1.1,0.95) -- (7.2,0.95) node[midway,above=4pt] {$e^{-i\gamma H_C}$: one $R_{ZZ}(\gamma)$ per edge};
  \draw[decorate,decoration={brace,amplitude=4pt}] (7.45,0.35) -- (8.35,0.35) node[midway,above=4pt] {$e^{-i\beta H_M}$};
  \node[note,anchor=north] (n1) at (1.6,-2.75) {CNOT, $R_Z(\gamma)$, CNOT:\\one $R_{ZZ}(\gamma)$};
  \draw[notearrow] (n1.north) -- (2.0,-2.5);
  \node[note,anchor=north] (n2) at (5.4,-2.75) {edge $41$ spans qubits $2$ and $3$;\\on a line it needs SWAP gates};
  \draw[notearrow] (n2.north) to[bend right=10] (5.6,-1.85);
  \node[note,anchor=north] (n3) at (8.95,-2.75) {one bit string\\per copy};
  \draw[notearrow] (n3.north) -- (8.85,-2.48);
\end{tikzpicture}
\end{center}
\ona{\caption{The depth-one QAOA circuit for MaxCut on the 4-cycle. Each $ZZ$ rotation is compiled as in \eqref{eq:ch05:rzz} into two CNOTs around a $Z$ rotation of the target; the rotations of the edges $12$ and $34$ act on disjoint qubits and run in parallel.}\label{fig:ch05:circuit}\end{figure}}
\end{frame}

\begin{frame}\ft{Estimating the energy from measurements}
\onp{\footnotesize}
\ona{A measurement of all four qubits returns a single bit string $x$, with probability $|\braket{x|\psi(\beta,\gamma)}|^2$, and destroys the state. The energy is therefore estimated from $M$ copies of the state, each prepared afresh and measured once,}
\begin{equation}
    \hat F = -\frac1M\sum_{m=1}^M\mathrm{cut}(x_m), \qquad\text{with standard error}\quad \frac{\hat\sigma}{\sqrt M}.
\end{equation}
\ona{At the starting point $(0.1,-0.1)$, the first twelve of $M = 1000$ copies gave the bit strings $x_1x_2x_3x_4$ shown below, with their cuts. Of the $1000$ copies, $98$ cut no edge, $765$ cut two and $137$ cut all four, against the probabilities $0.106$, $0.749$ and $0.145$ (Fig.~\ref{fig:ch05:shots}). The estimate $\hat F = -(2\cdot765 + 4\cdot137)/1000 = -2.078\pm0.031$ lies within one standard error of the true expectation $-2.077$. Halving the error requires four times as many copies; the numbers of copies needed are discussed in \Chap~\chref{C.PerIteration}.}
\begin{center}\scriptsize
\texttt{0011}:2\quad\texttt{1010}:4\quad\texttt{0111}:2\quad\texttt{0101}:4\quad\texttt{0101}:4\quad\texttt{1100}:2\quad\texttt{1110}:2\quad\texttt{0010}:2\quad\texttt{1010}:4\quad\texttt{0100}:2\quad\texttt{1111}:0\quad\texttt{1110}:2\quad$\cdots$
\end{center}
\ona{\gdef\shH{0.36\linewidth}}\onp{\gdef\shH{0.3\linewidth}}
\ona{\begin{figure}[htb]}\begin{center}
\begin{tikzpicture}
\begin{axis}[ybar,bar width=14pt,width=0.42\linewidth,height=\shH,clip=false,axis lines=left,xmin=-1,xmax=5,ymin=0,ymax=0.85,xtick={0,2,4},ytick={0,0.25,0.5,0.75},xlabel={cut of the measured bit string},ylabel={fraction of copies},enlargelimits=false]
  \addplot[fill=blue!45,draw=blue!70!black] table[x=cut,y expr=\thisrow{count}/1000] {figures/ch05/worked_shots.dat};
  \addplot[only marks,mark=-,mark size=9pt,very thick,red!75!black] table[x=cut,y=prob] {figures/ch05/worked_shots.dat};
  \node[note,anchor=west] (s1) at (axis cs:5.4,0.72) {$765$ of the $1000$ copies\\cut two edges};
  \draw[notearrow] (s1.west) to[bend right=10] (axis cs:2.3,0.76);
  \node[note,anchor=west] (s2) at (axis cs:5.4,0.38) {red: the true\\probabilities};
  \draw[notearrow] (s2.west) to[bend left=10] (axis cs:4.35,0.15);
  \node[note,anchor=west] (s3) at (axis cs:5.4,0.06) {$\hat F = -2.078\pm0.031$; true\\expectation $-2.077$};
\end{axis}
\end{tikzpicture}
\end{center}
\ona{\caption{The cuts of the $1000$ bit strings measured at $(\beta,\gamma) = (0.1,-0.1)$ (blue), and their probabilities (red).}\label{fig:ch05:shots}\end{figure}}
\end{frame}

\mode<all>{\gdef\psrtab{\begin{tabular}{@{}lrrr@{\qquad}lrrr@{}}\toprule
gate & $\hat F_+$ & $\hat F_-$ & $\frac12(\hat F_+-\hat F_-)$ & gate & $\hat F_+$ & $\hat F_-$ & $\hat F_+-\hat F_-$ \\\midrule
$R_Z$, edge $12$ & $-1.832$ & $-2.260$ & $+0.214$ & $R_X$, qubit $1$ & $-2.168$ & $-1.942$ & $-0.226$ \\
$R_Z$, edge $23$ & $-1.832$ & $-2.222$ & $+0.195$ & $R_X$, qubit $2$ & $-2.098$ & $-1.906$ & $-0.192$ \\
$R_Z$, edge $34$ & $-1.824$ & $-2.212$ & $+0.194$ & $R_X$, qubit $3$ & $-2.148$ & $-1.952$ & $-0.196$ \\
$R_Z$, edge $41$ & $-1.838$ & $-2.228$ & $+0.195$ & $R_X$, qubit $4$ & $-2.200$ & $-1.948$ & $-0.252$ \\\midrule
\multicolumn{3}{@{}l}{$\hat g_\gamma$, the sum (true $+0.763$)} & $+0.798$ & \multicolumn{3}{l}{$\hat g_\beta$, the sum (true $-0.732$)} & $-0.866$ \\\bottomrule
\end{tabular}}}
\begin{frame}\ft{The gradient from shifted circuits}
\onp{\footnotesize}
\ona{The energy depends on the angle $\theta$ of a single gate $R_P(\theta) = e^{-i\theta P/2}$, with $P$ a Pauli string, as $a + b\cos\theta + c\sin\theta$, and hence, exactly, $\partial F/\partial\theta = \frac12[F(\theta+\frac\pi2) - F(\theta-\frac\pi2)]$; \Chap~\chref{C.PSR} proves this in general. Here $\gamma$ is the angle of four gates $R_Z(\gamma)$, one for each edge, and $2\beta$ that of four gates $R_X(2\beta)$, one for each qubit. Shifting one gate at a time and applying the chain rule gives}
\onp{For a single gate $R_P(\theta) = e^{-i\theta P/2}$, $\partial F/\partial\theta = \frac12[F(\theta+\frac\pi2) - F(\theta-\frac\pi2)]$ (\Chap~\chref{C.PSR}). Here $\gamma$ is the angle of four gates $R_Z(\gamma)$ and $2\beta$ that of four gates $R_X(2\beta)$, and the gates are shifted one at a time:}
\begin{equation}
    \frac{\partial F}{\partial\gamma} = \sum_{e\in E}\tfrac12\big[F_{e+} - F_{e-}\big], \qquad \frac{\partial F}{\partial\beta} = \sum_{j=1}^4 2\cdot\tfrac12\big[F_{j+} - F_{j-}\big],
\end{equation}
\ona{where $F_{e\pm}$ denotes the energy with only the $R_Z$ of edge $e$ shifted by $\pm\pi/2$, and $F_{j\pm}$ the energy with only the $R_X$ of qubit $j$ shifted by $\pm\pi/2$. This requires $16$ circuits, each estimated from $1000$ copies (Table~\ref{tab:ch05:psr}). By symmetry, the true contributions of all the edges are equal, to $0.191$, as are those of all the qubits, to $-0.183$; the estimates scatter about these values because of the shot noise.}
\onp{\begin{center}\scriptsize\psrtab\end{center}}
\ona{\begin{table}[htb]\centering\caption{The sixteen shifted circuits at the starting point $(\beta,\gamma) = (0.1,-0.1)$, each estimated from $1000$ copies, and the resulting estimates of the gradient.}\label{tab:ch05:psr}\footnotesize\psrtab\end{table}}
\end{frame}

\mode<all>{\gdef\itertab{\begin{tabular}{@{}rrrrrrr@{}}\toprule
$k$ & $\beta_k$ & $\gamma_k$ & $\hat F$ & $F$, true & $\hat g_\beta$ \ (true) & $\hat g_\gamma$ \ (true) \\\midrule
$0$ & $0.100$ & $-0.100$ & $-2.078\pm0.031$ & $-2.077$ & $-0.866$ \ ($-0.732$) & $+0.798$ \ ($+0.763$) \\
$1$ & $0.187$ & $-0.180$ & $-2.284\pm0.032$ & $-2.239$ & $-0.962$ \ ($-1.033$) & $+1.314$ \ ($+1.271$) \\
$2$ & $0.283$ & $-0.311$ & $-2.506\pm0.030$ & $-2.528$ & $-1.010$ \ ($-0.992$) & $+1.516$ \ ($+1.470$) \\
$3$ & $0.384$ & $-0.463$ & $-2.788\pm0.032$ & $-2.798$ & $-0.146$ \ ($-0.114$) & $+1.216$ \ ($+1.202$) \\
$4$ & $0.398$ & $-0.584$ & $-2.884\pm0.034$ & $-2.920$ & $+0.066$ \ ($+0.084$) & $+0.658$ \ ($+0.782$) \\
$5$ & $0.392$ & $-0.650$ & $-2.918\pm0.034$ & $-2.964$ & $-0.064$ \ ($-0.014$) & $+0.447$ \ ($+0.534$) \\
$6$ & $0.398$ & $-0.695$ & $-3.054\pm0.033$ & $-2.983$ & $+0.140$ \ ($+0.087$) & $+0.337$ \ ($+0.360$) \\\bottomrule
\end{tabular}}}
\begin{frame}\ft{Six iterations}
\onp{\footnotesize}
\ona{Table~\ref{tab:ch05:iter} shows six steps of gradient descent with step size $0.1$. The quantities with a hat are estimated from measured copies, and the true expectations and gradients are given only for comparison. Each iteration requires $17$ circuits, one for the energy and $16$ for the gradient, and so $17\,000$ copies.}
\onp{\begin{center}\scriptsize\itertab\end{center}}
\ona{\begin{table}[htb]\centering\caption{Gradient descent on the depth-one QAOA energy of the 4-cycle, from $(0.1,-0.1)$ with step size $0.1$. The estimates use $1000$ copies for each circuit; the true expectations and gradients are given for comparison. The optimum is $F = -3$ at $(\beta,\gamma) = (\pi/8,-\pi/4)\approx(0.393,-0.785)$.}\label{tab:ch05:iter}\footnotesize\itertab\end{table}}
\begin{itemize}\setlength\itemsep{0.15em}
    \item The angle $\beta$ settles near $\pi/8$ within three steps, while $\gamma$ approaches $-\pi/4$ more slowly, since the landscape is four times less curved in $\gamma$ than in $\beta$.
    \item At $k = 6$ the estimate $\hat F = -3.054$ lies below the minimum $-3$. The variational principle bounds the energy, not its estimate.
    \item The noise in an estimated component of the gradient is about $0.09$ for $\hat g_\beta$, which is a sum of four differences, and $0.04$ for $\hat g_\gamma$, which carries the factor $\frac12$. Near the optimum the true $g_\beta$ is smaller than this, and further progress requires more copies or a smaller step.
\end{itemize}
\end{frame}

\begin{frame}\ft{The iterates on the landscape}
\ona{Figure~\ref{fig:ch05:workedpath} shows the six steps on the landscape \eqref{eq:ch05:c4}. The starting point lies near the saddle point at the origin, where the gradient is small; the steps then lengthen, and shorten again near the minimum.}
\ona{\gdef\wpH{0.46\linewidth}\gdef\wpW{0.5\linewidth}\gdef\wpC{-0.2}}\onp{\gdef\wpH{0.36\linewidth}\gdef\wpW{0.4\linewidth}\gdef\wpC{-0.4}}
\ona{\begin{figure}[htb]}\begin{center}
\begin{tikzpicture}
\begin{axis}[name=wp,width=\wpW,height=\wpH,view={0}{90},xmin=0,xmax=0.785,ymin=-1.571,ymax=0,xtick={0,0.196,0.393,0.589,0.785},xticklabels={$0$,$\frac\pi{16}$,$\frac\pi8$,$\frac{3\pi}{16}$,$\frac\pi4$},ytick={-1.571,-1.178,-0.785,-0.393,0},yticklabels={$-\frac\pi2$,$-\frac{3\pi}8$,$-\frac\pi4$,$-\frac\pi8$,$0$},xlabel={$\beta$},ylabel={$\gamma$},colorbar horizontal,colorbar style={at={(0.5,\wpC)},anchor=north,height=0.18cm,font=\tiny},colormap/viridis,point meta min=-3,point meta max=-2,enlargelimits=false,axis on top]
  \addplot3[surf,shader=interp,domain=0:0.785,y domain=-1.571:0,samples=41] {-2+sin(deg(4*x))*sin(deg(2*y))};
  \addplot3[white,thick,mark=*,mark size=1.6pt] table[x=beta,y=gamma,z expr=0] {figures/ch05/worked_path.dat};
  \fill[red] (axis cs:0.3927,-0.7854,0) circle (2.5pt);
  \coordinate (st) at (axis cs:0.1,-0.1,0);
  \coordinate (mid) at (axis cs:0.283,-0.311,0);
  \coordinate (end) at (axis cs:0.398,-0.695,0);
  \coordinate (opt) at (axis cs:0.3927,-0.7854,0);
\end{axis}
  \node[note,anchor=west] (w1) at ([xshift=0.4cm,yshift=1.1cm]wp.east) {start $(0.1,-0.1)$, near\\the saddle at the origin};
  \draw[notearrow] (w1.west) to[bend right=10] (st);
  \node[note,anchor=west] (w2) at ([xshift=0.4cm,yshift=0.25cm]wp.east) {each step requires $17$ circuits\\and $17\,000$ copies};
  \draw[notearrow] (w2.west) to[bend right=10] (mid);
  \node[note,anchor=west] (w3) at ([xshift=0.4cm,yshift=-0.55cm]wp.east) {after six steps, $F = -2.983$};
  \draw[notearrow] (w3.west) to[bend left=10] (end);
  \node[note,anchor=west] (w4) at ([xshift=0.4cm,yshift=-1.2cm]wp.east) {the minimum $-3$,\\at $(\pi/8,-\pi/4)$};
  \draw[notearrow] (w4.west) to[bend left=15] (opt);
\end{tikzpicture}
\end{center}
\ona{\caption{The iterates of Table~\ref{tab:ch05:iter} (white) on the landscape $F(\beta,\gamma) = -2 + \sin4\beta\sin2\gamma$, with its minimum (red). The colour scale stops at $-2$, the value along the axes.}\label{fig:ch05:workedpath}\end{figure}}
\onp{\footnotesize Six noisy gradient steps, from near the saddle point to near the minimum.}
\end{frame}

\begin{frame}\ft{Observations on the worked example}\label{fr:ch05:workedlast}
\onp{\footnotesize}
\begin{itemize}\setlength\itemsep{0.25em}
    \item Six steps used $6\cdot17 = 102$ circuits and $102\,000$ copies. With the figures for hardware given in \Chap~\chref{C.PerIteration}, about $100\,\mu$s per copy and $6$~s of latency for each submitted circuit, this amounts to $10$~s of sampling and $10$ minutes of latency.
    \item The standard error of an energy is $0.03$, and that of a component of the gradient between $0.04$ and $0.09$; near the optimum the gradient can no longer be distinguished from the noise.
    \item At the last iterate a maximum cut is measured with probability $0.51$, in $526$ of $1000$ copies. The output of the algorithm is the best bit string observed, not the energy.
    \item At depth one no choice of angles gives an expected cut above $3$ of the $4$ edges; more layers can do better (\Chap~\chref{C.QAOA}).
    \item The landscape of this example is benign, with two parameters, a single minimum in the region of interest, and a starting point near a saddle. The theorems of \citet{bittel2021training}, later in this \chap{}, show how different the worst case is.
\end{itemize}
\ona{The instance is, of course, classically trivial, the 4-cycle being bipartite. The example serves only to exhibit the structure of a single run, which is the same at a hundred qubits.}
\end{frame}

\section{Training as a computational problem}\label{sec:ch05:formal}

\begin{frame}\ft{What a formalisation must specify}
\onp{\footnotesize}
\ona{Before one can prove that training is hard, one must say what training is. The loop of Fig.~\ref{fig:ch05:vqa} leaves a good deal unspecified, and \citet{bittel2021training} make it precise as follows (Fig.~\ref{fig:ch05:sketch}). From now on we follow them in writing $\boldsymbol\phi$ for the parameters.}
\begin{itemize}\setlength\itemsep{0.1em}
    \item The ansatz is a product $U_L(\phi_L)\cdots U_1(\phi_1)$ of \emph{layers} $U_i(\phi_i) = e^{-i\phi_iH_i}$, each generated by a given Hermitian matrix $H_i$ and controlled by a single real parameter, so that $\boldsymbol\phi\in\R^L$.
    \item The input consists of the $H_i$ and $O$, and the manner in which they are written down determines the size of an instance.
    \item The quantum computer is replaced by an oracle that returns $\braket{O(\boldsymbol\phi)}$, or by a classical computation where this is efficient.
    \item A solution consists of parameters whose energy is optimal up to an exponentially small error.
    \item The error is measured relative to the spectral width of $O$, and is divided into a part due to the ansatz and a part due to the optimizer.
\end{itemize}
\onp{\begin{center}\includegraphics[height=0.26\textheight]{figures/bittel/sketch.pdf}\end{center}\vspace{-0.6em}{\tiny After \citet{bittel2021training}; the classical optimization is shown in red.}}
\ona{\begin{figure}[htb]\centering\includegraphics[scale=0.86]{figures/bittel/sketch.pdf}\caption{A variational routine, after \citet{bittel2021training}. Parametrised layers $U_i(\phi_i)$ act on $\ket0^{\otimes N}$; a fixed unitary $U_{\rm M}$, measurements in the computational basis and post-processing estimate $\braket{O(\boldsymbol\phi)}$; and a classical optimizer proposes new parameters $\boldsymbol\phi_{\rm guess}$ until it returns $\boldsymbol\phi_{\rm final}$. The formalisation isolates the classical optimization (red).}\label{fig:ch05:sketch}\end{figure}}
\ona{The fourth point requires comment. The minimum of a real function is a real number which may have no finite representation, and the same is true of the parameters at which it is attained. We adopt the usual convention that a minimisation problem is solved if it is solved to exponential precision, that is, up to an error $2^{-\mathrm{poly}}$ in the size of the input. The hardness results below do not, therefore, derive from any difficulty in writing down the minimum.}
\end{frame}

\begin{frame}\ft{Training with a quantum computer: the oracular problem}
\onp{\footnotesize}
\begin{probl}[VQA minimisation, oracular formulation]\label{p:ch05:oracle}
\begin{description}\setlength\itemsep{0.1em}
    \item[Instance] Generators $H_1,\dots,H_L$ and an observable $O$ on $N$ qubits, each given as a real combination of polynomially many Pauli strings.
    \item[Oracle] For $\ket{\Psi(\boldsymbol\phi)} = U_L(\phi_L)\cdots U_1(\phi_1)\ket{0}^{\otimes N}$ with $U_i(\phi) = e^{-i\phi H_i}$, an oracle returns $\braket{O(\boldsymbol\phi)} = \braket{\Psi(\boldsymbol\phi)|O|\Psi(\boldsymbol\phi)}$ for any $\boldsymbol\phi$ up to any desired additive error $1/\mathrm{poly}$.
    \item[Task] Find $\boldsymbol\phi\in\R^L$ minimising $\braket{O(\boldsymbol\phi)}$.
\end{description}
\end{probl}
\ona{The oracle represents the quantum computer, idealised in two respects. It is deterministic, as if one post-selected on the runs that give a correct estimate, and it delivers whatever precision the algorithm requests, provided this is inverse-polynomial, which polynomially many shots suffice to achieve. The noise that this idealisation conceals is studied in \Chap~\chref{C.IterationComplexity}. An algorithm for Problem~\ref{p:ch05:oracle} is a classical algorithm that may call the oracle, each call counting as a single step. Whatever is hard for such an algorithm is therefore hard because of the classical optimization alone, and a better quantum computer would not help. The variational algorithms of Table~\ref{tab:ch05:fit} are instances of the problem.}
\begin{propn}[\citet{bittel2021training}]\label{propn:ch05:oracle}
Unless $\mathbf P = \mathbf{NP}$, no deterministic classical algorithm solves \ona{Problem~\ref{p:ch05:oracle}}\onp{the oracular problem} in polynomial time.
\end{propn}
\ona{This alone is not surprising. If the ansatz can reach every computational basis state and $O$ is a diagonal Ising Hamiltonian, then minimising over the ansatz includes finding the ground state of $O$, which is $\mathbf{NP}$-hard (\Chap~\chref{C.Inapprox}). The interest lies in how little the proof requires, and in the variants below, in which the ground-state problem is easy and training is nevertheless hard.}
\end{frame}

\begin{frame}\ft{The oracle, pictured}
\ona{Figure~\ref{fig:ch05:oracle} shows Problem~\ref{p:ch05:oracle} as a circuit inside a box. Everything inside the box is the idealised quantum computer, and everything outside it is the classical algorithm, whose running time the hardness results bound from below.}
\ona{\begin{figure}[htb]}\begin{center}
\begin{tikzpicture}[every node/.style={font=\scriptsize},layer/.style={draw,fill=blue!6,rounded corners=2pt,minimum width=1.05cm,minimum height=1.55cm,align=center}]
  \draw[rounded corners=8pt,dashed,thick,gray,fill=gray!4] (-0.75,0.55) rectangle (6.15,-1.95);
  \node[gray!60!black,anchor=south east] at (6.1,-1.9) {oracle $\mathcal O$};
  \foreach \q in {0,1,2} {\draw (0,-0.6*\q) node[left] {$\ket0$} -- (5.4,-0.6*\q);}
  \node[layer] (u1) at (0.9,-0.6) {$e^{-i\phi_1H_1}$};
  \node[layer] (u2) at (2.2,-0.6) {$e^{-i\phi_2H_2}$};
  \node[fill=gray!4] at (3.2,-0.6) {$\cdots$};
  \node[layer] (uL) at (4.2,-0.6) {$e^{-i\phi_LH_L}$};
  \node[draw,fill=green!10,rounded corners=2pt,minimum width=0.7cm,minimum height=1.55cm] (m) at (5.5,-0.6) {$O$};
  \node[draw,rounded corners,fill=red!8,minimum width=4.2cm,minimum height=0.7cm] (alg) at (2.7,-3.1) {classical algorithm};
  \draw[-{Latex},thick] (1.7,-2.75) -- node[left] {$\boldsymbol\phi$} (1.7,-1.95);
  \draw[-{Latex},thick] (3.7,-1.95) -- node[right] {$\braket{O(\boldsymbol\phi)}\pm\epsilon$} (3.7,-2.75);
  \node[note,anchor=west] (n1) at (6.45,0.45) {one real parameter per layer;\\a layer may act on all qubits};
  \draw[notearrow] (n1.west) to[bend right=15] (4.75,0.1);
  \node[note,anchor=west] (n2) at (6.45,-0.6) {the input: the $H_\ell$ and $O$,\\as Pauli strings with coefficients};
  \draw[notearrow] (n2.west) -- (5.86,-0.6);
  \node[note,anchor=west] (n3) at (6.45,-1.65) {deterministic, error $1/\mathrm{poly}$:\\shot noise idealised away};
  \draw[notearrow] (n3.west) -- (6.16,-1.5);
  \node[note,anchor=west] (n4) at (6.45,-3.1) {each call counts as one step;\\the hardness lies here};
  \draw[notearrow] (n4.west) -- (alg.east);
\end{tikzpicture}
\end{center}
\ona{\caption{Problem~\ref{p:ch05:oracle}. The oracle, shown as the dashed box, prepares $\ket{\Psi(\boldsymbol\phi)}$ by applying $L$ layers to $\ket0^{\otimes N}$ and returns the energy of $O$ to the requested precision; the classical algorithm chooses the parameters. Lower bounds on the running time of the classical algorithm hold however good the quantum computer may be.}\label{fig:ch05:oracle}\end{figure}}
\onp{\footnotesize Inside the box is an idealised quantum computer; outside it is the classical optimizer, whose cost the theorems bound.}
\end{frame}

\mode<all>{\gdef\classtab#1#2#3#4{{\setlength\tabcolsep{4pt}\begin{tabular}{@{}>{\raggedright\arraybackslash}p{#1\linewidth}>{\raggedright\arraybackslash}p{#2\linewidth}>{\raggedright\arraybackslash}p{#3\linewidth}>{\raggedright\arraybackslash}p{#4\linewidth}@{}}\toprule
 & $\mathbf{NP}$ & $\mathbf{QCMA}$ & $\mathbf{QMA}$ \\\midrule
certificate & a bit string $y$ & a bit string $y$ & a state $\ket\xi$ of $p(|x|)$ qubits \\
verifier & a deterministic classical algorithm & a quantum circuit & a quantum circuit \\
yes instance & some $y$ is accepted & some $y$ is accepted with probability $\ge\frac23$ & some $\ket\xi$ is accepted with probability $\ge\frac23$ \\
no instance & every $y$ is rejected & every $y$ is accepted with probability $\le\frac13$ & every $\ket\xi$ is accepted with probability $\le\frac13$ \\
complete problem & SAT; MaxCut, decision version & low energy of a $3$-local Hamiltonian in a state prepared by $k$ gates & low energy of a $2$-local Hamiltonian in any state \\
in this \chap{} & training in polynomial dimension & training with a quantum computer & the ground energy itself \\\bottomrule
\end{tabular}}}}
\begin{frame}\ft{The classes $\mathbf{NP}$, $\mathbf{QCMA}$ and $\mathbf{QMA}$}\label{fr:ch05:classes}
\onp{\footnotesize}
\ona{We shall place the training problems among three classes of decision problems, which differ in the kind of certificate that a prover supplies and in the kind of verifier that checks it. The class $\mathbf{NP}$ was recalled in \Chap~\chref{C.Inapprox}, and $\mathbf{QMA}$ was described informally in \Chap~\chref{C.EQA}; here we define both, together with the intermediate class $\mathbf{QCMA}$. We consider \emph{promise problems} $A = (A_{\rm yes},A_{\rm no})$, pairs of disjoint sets of instances, for which nothing is required of instances in neither set; this is natural when energies are to be compared with a threshold only up to a gap. A quantum verifier is a family of quantum circuits $V_x$, one for each instance $x$, of size polynomial in $|x|$ and generated by a classical algorithm in polynomial time; it acts on the certificate and on polynomially many ancillary qubits initialised to $\ket0$, and its output is the result of measuring one qubit.}
\begin{defn}\label{defn:ch05:classes}
Let $A = (A_{\rm yes},A_{\rm no})$ be a promise problem and $p$ a polynomial.
\begin{itemize}\setlength\itemsep{0.1em}
    \item[(i)] $A\in\mathbf{NP}$ if there is a deterministic classical algorithm $V$, running in polynomial time, such that for every $x\in A_{\rm yes}$ there is $y\in\{0,1\}^{p(|x|)}$ with $V(x,y) = 1$, and for every $x\in A_{\rm no}$, $V(x,y) = 0$ for all $y$.
    \item[(ii)] $A\in\mathbf{QCMA}$ if there is a quantum verifier $V_x$ such that for every $x\in A_{\rm yes}$ there is $y\in\{0,1\}^{p(|x|)}$ which $V_x$ accepts with probability at least $\frac23$, and for every $x\in A_{\rm no}$, $V_x$ accepts every $y$ with probability at most $\frac13$.
    \item[(iii)] $A\in\mathbf{QMA}$ if the same holds with the bit strings $y$ replaced by arbitrary states $\ket\xi$ of $p(|x|)$ qubits.
\end{itemize}
\end{defn}
\end{frame}

\begin{frame}\ft{The three classes compared}
\onp{\footnotesize}
\ona{Table~\ref{tab:ch05:classes} sets the three definitions side by side. Since a classical algorithm can be simulated by a quantum circuit, and a bit string is a computational basis state, $\mathbf{NP}\subseteq\mathbf{QCMA}\subseteq\mathbf{QMA}$; and $\mathbf{BQP}\subseteq\mathbf{QCMA}$, as a verifier may ignore its certificate. Neither containment is known to be strict. The constants $\frac23$ and $\frac13$ are arbitrary: any two thresholds separated by an inverse polynomial in $|x|$ give the same classes, since the verification may be repeated and the majority taken, with as many copies of the certificate as there are repetitions \citep{kitaev2002classical}. Deciding whether the ground energy of a $2$-local Hamiltonian is below $a$ or above $b$, with $b - a\ge1/\mathrm{poly}$, is $\mathbf{QMA}$-complete \citep{kitaev2002classical,kempe2006complexity}. A $\mathbf{QCMA}$ certificate, by contrast, is a classical description, such as a circuit that prepares a low-energy state. Indeed, deciding whether a $3$-local Hamiltonian has a state of energy below a given value that can be prepared by at most $k$ elementary gates is $\mathbf{QCMA}$-complete \citep{wocjan2003two}; the restriction on the preparation is precisely what distinguishes this problem from the $\mathbf{QMA}$-complete one, and it is the restriction that a variational ansatz imposes.}
\onp{\begin{center}\scriptsize\classtab{0.13}{0.22}{0.27}{0.27}\end{center}
$\mathbf{NP}\subseteq\mathbf{QCMA}\subseteq\mathbf{QMA}$ and $\mathbf{BQP}\subseteq\mathbf{QCMA}$; neither containment is known to be strict. The constants $\frac23$ and $\frac13$ may be replaced by any two thresholds separated by an inverse polynomial.}
\ona{\begin{table}[htb]\centering\caption{The classes $\mathbf{NP}$, $\mathbf{QCMA}$ and $\mathbf{QMA}$ of Definition~\ref{defn:ch05:classes}, for an instance $x$ and certificates of length $p(|x|)$, with the problems of this \chap{} that belong to them. The verifiers run in time polynomial in $|x|$, and the training problems are the decision versions of Problems~\ref{p:ch05:vqa} and \ref{p:ch05:oracle}.}\label{tab:ch05:classes}\footnotesize\classtab{0.15}{0.23}{0.27}{0.27}\end{table}}
\end{frame}

\begin{frame}\ft{Training without a quantum computer}
\onp{\footnotesize}
\ona{If the Hilbert space has polynomial dimension $n$, that is, if there are only $\mathcal O(\log n)$ qubits, then the state $\ket{\Psi(\boldsymbol\phi)}$ and its energy can be computed classically in polynomial time, and the oracle is not needed. It is then convenient to work in the computational basis.}
\begin{probl}[VQA minimisation]\label{p:ch05:vqa}
\begin{description}\setlength\itemsep{0.1em}
    \item[Instance] An initial state $\ket{\Psi_0}\in\C^n$, generators $H_1,\dots,H_L\in\Herm(\C^n)$ and an observable $O\in\Herm(\C^n)$, given as vectors and matrices.
    \item[Task] Find $\boldsymbol\phi\in\R^L$ minimising $\braket{O(\boldsymbol\phi)}$ for $\ket{\Psi(\boldsymbol\phi)} = U_L(\phi_L)\cdots U_1(\phi_1)\ket{\Psi_0}$.
\end{description}
\end{probl}
\ona{Since a matrix on $\lceil\log_2n\rceil$ qubits has at most $\mathcal O(n^2)$ Pauli coefficients, this is a special case of Problem~\ref{p:ch05:oracle}, with an oracle that can be computed classically. The \emph{decision version} asks, for a given threshold $a\in\R$, whether some $\boldsymbol\phi$ achieves $\braket{O(\boldsymbol\phi)}\le a$. With the convention on precision, it lies in $\mathbf{NP}$, the parameters being a certificate that can be checked in polynomial time. The decision version of Problem~\ref{p:ch05:oracle}, with a polynomial-time quantum computer in place of the oracle and an inverse-polynomial gap about $a$, lies in $\mathbf{QCMA}$: the parameters are a classical certificate, and the quantum computer checks it \citep{bittel2021training,bittelQCMA}. This allows the two sources of difficulty to be distinguished.}
\begin{itemize}\setlength\itemsep{0.1em}
    \item The ansatz. Deciding the ground energy of a local Hamiltonian is $\mathbf{QMA}$-complete \citep{kitaev2002classical,kempe2006complexity}, the certificate being a quantum state. If a polynomial-size ansatz always reached the ground energy, its parameters would be classical certificates; hence, if $\mathbf{QCMA}\ne\mathbf{QMA}$, every such ansatz fails on some Hamiltonians, whatever the optimizer.
    \item The optimizer. Even when the ground-state problem is easy and the ground energy is within reach of the ansatz, finding the best parameters is $\mathbf{NP}$-hard, as we show below.
\end{itemize}
\end{frame}

\begin{frame}\ft{Where the training problems live}
\ona{Figure~\ref{fig:ch05:classes} places the decision versions of the training problems among the classes of Definition~\ref{defn:ch05:classes}. The classes are distinguished by the kind of certificate: classical parameters checked classically, classical parameters checked on a quantum computer, or a quantum state.}
\ona{\begin{figure}[htb]}\begin{center}
\begin{tikzpicture}[every node/.style={font=\scriptsize,align=center}]
  \draw[rounded corners=12pt,thick,fill=blue!5] (-4.4,-2.25) rectangle (2.4,2.25);
  \node[anchor=north west,align=left] at (-4.3,2.2) {$\mathbf{QMA}$: quantum certificates\\local Hamiltonians: $\mathbf{QMA}$-complete};
  \draw[rounded corners=10pt,thick,fill=green!7] (-4.05,-1.95) rectangle (2.05,0.95);
  \node[anchor=north west,align=left] at (-3.95,0.9) {$\mathbf{QCMA}$: classical certificate, quantum check\\training on qubits, with a quantum computer};
  \draw[rounded corners=8pt,thick,fill=orange!10] (-3.7,-1.7) rectangle (0.7,-0.3);
  \node at (-1.5,-1.0) {$\mathbf{NP}$\\training in polynomial dimension:\\$\mathbf{NP}$-complete (step 3)};
  \node[note,anchor=west] (nq) at (2.75,1.7) {certificate: the ground\\state itself, a quantum state};
  \draw[notearrow] (nq.west) -- (2.25,1.55);
  \node[note,anchor=west] (nc) at (2.75,0.45) {certificate: the parameters $\boldsymbol\phi$,\\checked on a quantum computer};
  \draw[notearrow] (nc.west) -- (1.9,0.3);
  \node[note,anchor=west] (nn) at (2.75,-1.0) {certificate: $\boldsymbol\phi$, checked\\classically in polynomial time};
  \draw[notearrow] (nn.west) to[bend left=10] (0.6,-1.0);
  \node[note,anchor=north] (nm) at (-1.0,-2.5) {if $\mathbf{QCMA}\ne\mathbf{QMA}$, some ground states lie beyond every polynomial-size ansatz,\\and the model mismatch cannot always be removed, whatever the optimizer};
  \draw[notearrow] ([xshift=-2.6cm]nm.north) to[bend left=20] (-4.22,-1.6);
\end{tikzpicture}
\end{center}
\ona{\caption{The decision versions of the training problems. In polynomial dimension the parameters are a certificate that can be checked classically, and the problem is $\mathbf{NP}$-complete. With a quantum computer in place of the oracle, the parameters are a classical certificate checked on a quantum computer, as in $\mathbf{QCMA}$. Deciding the ground energy of a local Hamiltonian is $\mathbf{QMA}$-complete, its certificate being a ground state \citep{kitaev2002classical,kempe2006complexity}. The containments $\mathbf{NP}\subseteq\mathbf{QCMA}\subseteq\mathbf{QMA}$ are not known to be strict.}\label{fig:ch05:classes}\end{figure}}
\onp{\footnotesize Certificates: $\boldsymbol\phi$ checked classically ($\mathbf{NP}$), $\boldsymbol\phi$ checked on a quantum computer ($\mathbf{QCMA}$), a quantum state ($\mathbf{QMA}$).}
\end{frame}

\begin{frame}\ft{QAOA as a training problem}
\onp{\footnotesize}
\begin{probl}[QAOA minimisation]\label{p:ch05:qaoa}
\begin{description}\setlength\itemsep{0.1em}
    \item[Instance] Hamiltonians $H_b,H_c\in\Herm(\C^n)$, and the number of layers $L$, written in unary.
    \item[Task] For $\ket{\Psi(\boldsymbol\beta,\boldsymbol\gamma)} = U_b(\beta_L)U_c(\gamma_L)\cdots U_b(\beta_1)U_c(\gamma_1)\ket{\Psi_0}$, where $\ket{\Psi_0}$ is the ground state of $H_b$, $U_b(\beta) = e^{-i\beta H_b}$ and $U_c(\gamma) = e^{-i\gamma H_c}$, find $\boldsymbol\beta,\boldsymbol\gamma\in\R^L$ minimising $\braket{H_c(\boldsymbol\beta,\boldsymbol\gamma)} = \braket{\Psi(\boldsymbol\beta,\boldsymbol\gamma)|H_c|\Psi(\boldsymbol\beta,\boldsymbol\gamma)}$.
\end{description}
\end{probl}
\ona{This problem differs from Problem~\ref{p:ch05:vqa} in three respects: there are only two generators, used alternately, as in the discretised annealing of Section~\ref{sec:ch05:adiabatic}; the observable is one of them; and the initial state is not free, but is the ground state of $H_b$. QAOA for MaxCut is the case $H_b = -\sum_jX_j$, whose ground state is $\ket+^{\otimes n}$, and $H_c = H_C$, the diagonal cost Hamiltonian of Section~\ref{sec:ch05:qaoa}. The hardness proofs below use small matrices given explicitly, and therefore apply to any formulation that contains them.}
\begin{itemize}\setlength\itemsep{0.15em}
    \item The number of layers is written in unary because, in binary, an input of $\mathcal O(\log L)$ bits could ask for $L$ layers, and the evaluation of a single energy would take time exponential in the size of the input. In unary, the circuit is no longer than its description.
    \item If the eigenvalues of $H_b$ and $H_c$ are integers, the energy is $2\pi$-periodic in every angle, and one may restrict attention to $\boldsymbol\beta,\boldsymbol\gamma\in[0,2\pi)^L$, as for MaxCut.
\end{itemize}
\end{frame}

\begin{frame}\ft{Model mismatch and optimization error}
\onp{\footnotesize}
\ona{Let $\braket{O}_a$ be the energy at the parameters returned by an algorithm, $\braket{O}_{\min} = \min_{\boldsymbol\phi}\braket{O(\boldsymbol\phi)}$ the least energy attainable in the ansatz, and $w(O) = \lambda_{\max}(O) - \lambda_{\min}(O)$ the \emph{spectral width} of $O$. The \emph{approximation error} of a run is the sum of two parts (Fig.~\ref{fig:ch05:errors}),}
\begin{equation}\label{eq:ch05:split}
    \delta = \frac{\braket{O}_a - \lambda_{\min}(O)}{w(O)}
    = \underbrace{\frac{\braket{O}_{\min} - \lambda_{\min}(O)}{w(O)}}_{\delta_m\text{: model mismatch}}
    + \underbrace{\frac{\braket{O}_a - \braket{O}_{\min}}{w(O)}}_{\delta_o\text{: optimization error}}.
\end{equation}
\ona{\begin{figure}[htb]}\begin{center}
\begin{tikzpicture}[every node/.style={font=\scriptsize}]
  \draw[thick,-{Latex}] (0,0) -- (10.4,0) node[right] {energy};
  \foreach \x/\l in {0.3/{$\lambda_{\min}(O)$},3.3/{$\braket{O}_{\min}$},6.1/{$\braket{O}_a$},9.7/{$\lambda_{\max}(O)$}} {\draw[thick] (\x,0.12) -- (\x,-0.12) node[below] {\l};}
  \draw[decorate,decoration={brace,amplitude=5pt}] (0.3,0.22) -- (3.3,0.22) node[midway,above=6pt] {$\delta_m\,w(O)$};
  \draw[decorate,decoration={brace,amplitude=5pt}] (3.3,0.22) -- (6.1,0.22) node[midway,above=6pt] {$\delta_o\,w(O)$};
  \draw[decorate,decoration={brace,amplitude=5pt,mirror}] (0.3,-0.62) -- (9.7,-0.62) node[midway,below=6pt] {spectral width $w(O)$};
  \node[note,align=center] at (1.8,1.15) {the ansatz\\can do no better};
  \node[note,align=center] at (4.7,1.15) {left by the\\optimizer};
  \node[note,anchor=west] (nb) at (6.7,0.75) {Bittel and Kliesch bound this\\part from below, with $\delta_m = 0$};
  \draw[notearrow] (nb.west) to[bend right=20] (5.4,0.55);
\end{tikzpicture}
\end{center}
\ona{\caption{The two parts of the error of a variational algorithm, both in units of the spectral width. The model mismatch is a property of the ansatz; the optimization error is the part left by the optimizer, and is the subject of the hardness results.}\label{fig:ch05:errors}\end{figure}}
\begin{defn}[Optimization error, \citet{bittel2021training}]\label{defn:ch05:opterr}
The \emph{optimization error} of an algorithm is the smallest $\Delta\in[0,1]$ such that $\delta_o\le\Delta$ on every instance considered.
\end{defn}
\ona{Three remarks are in order. First, the normalisation places $\delta$, $\delta_m$ and $\delta_o$ in $[0,1]$ and makes them invariant under $O\mapsto aO + b\one$ with $a>0$, so that neither the unit of energy nor a constant shift matters. Second, $\Delta = 1$ is trivial, since it is achieved by any parameters whatever; a theorem that no polynomial-time algorithm has $\Delta<1$ therefore asserts that nothing at all can be guaranteed. Third, like an approximation ratio (\Chap~\chref{C.Inapprox}), $\Delta$ is a worst case over all instances in a class, here a class of VQAs. When $\lambda_{\max}(O) = 0$ and the ansatz contains a ground state, as in step 2 below, an output within a ratio $\alpha$ of the optimum has $\delta_o = 1-\alpha$.}
\onp{$\Delta = 1$ is trivial; $\Delta$ is to training what $1-\alpha$ is to approximation.}
\end{frame}

\section{Hardness of training: Bittel and Kliesch}\label{sec:ch05:hardness}

\begin{frame}\ft{Training a variational algorithm is $\mathbf{NP}$-hard}
\onp{\footnotesize}
\ona{Throughout this section, $A\in\{0,1\}^{d\times d}$ denotes the adjacency matrix of a simple, undirected, unweighted graph with $d$ vertices and $|E|\ge1$ edges, so that $A$ is symmetric with zero diagonal and $\sum_{ij}A_{ij} = 2|E|$, and $\mathrm{MC} = \mathrm{MaxCut}(A)$. In the terms just introduced, the main results of \citet{bittel2021training} are as follows.}
\begin{thm}[\citet{bittel2021training}]\label{thm:ch05:bittel}
The following are $\mathbf{NP}$-hard:
\begin{enumerate}\setlength\itemsep{0.05em}
    \item[(a)] \ona{Problem~\ref{p:ch05:oracle}}\onp{the oracular problem}, with single-qubit rotations on $d$ qubits;
    \item[(b)] \ona{Problem~\ref{p:ch05:vqa}}\onp{VQA minimisation} in dimension $2d$, that is, on $\mathcal O(\log d)$ qubits, with $d$ layers or with one, its decision version being $\mathbf{NP}$-complete;
    \item[(c)] \ona{Problem~\ref{p:ch05:qaoa}}\onp{QAOA minimisation} with one layer;
    \item[(d)] \ona{Problem~\ref{p:ch05:qaoa}}\onp{QAOA minimisation} with $d$ layers, angles in $[0,2\pi)$, and $\|H_b\|= 3$, $\|H_c\| = 1$;
    \item[(e)] the analogue of \ona{Problem~\ref{p:ch05:vqa}}\onp{VQA minimisation} for free fermions.
\end{enumerate}
Moreover, unless $\mathbf P = \mathbf{NP}$, no polynomial-time algorithm has optimization error $\Delta<1$ for (a), nor $\Delta < (1-\alpha_{\max})/2$ for (d), where $\alpha_{\max}$ is the best approximation ratio for MaxCut achievable in polynomial time.
\end{thm}
\ona{In (b), (c) and (e) all the energies can be computed classically in polynomial time, so that the hardness lies in the optimization itself, and is not inherited from a hard ground-state problem or a hard simulation. The proof is a chain of reductions from MaxCut (Fig.~\ref{fig:ch05:chain}). We give each of them in full, except for the construction behind (d), of which we explain only the idea.}
\end{frame}

\begin{frame}\ft{The chain of reductions}
\ona{\begin{figure}[htb]}\begin{center}
\begin{tikzpicture}[every node/.style={font=\scriptsize},box/.style={draw,rounded corners,fill=blue!6,minimum width=2.7cm,minimum height=0.9cm,align=center},lab/.style={font=\tiny,align=center,inner sep=1pt}]
  \node[box,fill=gray!12] (mc) at (0,0) {MaxCut\\$\max_v\mathrm{cut}(v)$};
  \node[box] (cmc) at (4.2,0) {continuous MaxCut\\$\min_{\boldsymbol\phi}\mu$ \quad(step 1)};
  \node[box] (orc) at (8.4,0) {oracle: $d$ qubits,\\$d$ layers \quad(step 2)};
  \node[box] (log) at (4.2,-1.75) {dimension $2d$,\\$d$ layers \quad(step 3)};
  \node[box] (qd) at (8.4,-1.75) {QAOA, $d$ layers,\\$\|H_b\| = 3$ \quad(step 6)};
  \node[box] (one) at (4.2,-3.5) {dimension $2d$,\\one layer \quad(step 4)};
  \node[box] (ff) at (8.4,-3.5) {free fermions,\\$2d$ modes \quad(step 7)};
  \node[box] (q1) at (0,-3.5) {QAOA,\\one layer \quad(step 5)};
  \draw[-{Latex},thick] (mc) -- node[lab,above=2pt] {$v_i = \cos\phi_i$;\\rounding\\loses nothing} (cmc);
  \draw[-{Latex},thick] (cmc) -- node[lab,above=2pt] {$R_Y(\phi_i)$\\on qubit $i$} (orc);
  \draw[-{Latex},thick] (cmc) -- node[lab,left=3pt] {phases $e^{\mp i\phi_j}$;\\$O$ a Laplacian} (log);
  \draw[-{Latex},thick] (cmc) -- node[lab,sloped,above] {$2d+1$ copies} (qd);
  \draw[-{Latex},thick] (log) -- node[lab,left=3pt] {a clock:\\$E_j = 2\pi/m^j$} (one);
  \draw[-{Latex},thick] (log) -- node[lab,sloped,above] {one fermion} (ff);
  \draw[-{Latex},thick] (one) -- node[lab,above=2pt] {additional level,\\coupling $\tau$} (q1);
  \node[note,anchor=north] at ([yshift=-2pt]orc.south) {no $\Delta<1$, by boosting};
  \node[note,anchor=north] at ([yshift=-2pt]qd.south) {no $\Delta<(1-\alpha_{\max})/2$};
\end{tikzpicture}
\end{center}
\ona{\caption{The reductions of \citet{bittel2021training}. Each arrow is a polynomial-time reduction that preserves the objective exactly or to any desired precision; the labels indicate the device used in each step, and the steps are numbered as in the text.}\label{fig:ch05:chain}\end{figure}}
\onp{\footnotesize Each step preserves the objective exactly or to any desired precision.}
\end{frame}

\begin{frame}\ft{Step 1: from MaxCut to cosines}
\onp{\scriptsize}
\begin{probl}[Continuous MaxCut]\label{p:ch05:cmc}
Given $A$, find $\boldsymbol\phi\in[0,2\pi)^d$ minimising
\begin{equation}\label{eq:ch05:mu}
    \mu(\boldsymbol\phi) \;=\; \tfrac14\sum_{i,j}A_{ij}\,\big(\cos\phi_i\cos\phi_j - 1\big).
\end{equation}
\end{probl}
\begin{lem}[\citet{bittel2021training}]\label{lem:ch05:cmc}
$\min_{\boldsymbol\phi}\mu = -\mathrm{MC}$, and from any $\boldsymbol\phi$ one obtains in polynomial time a cut of value at least $-\mu(\boldsymbol\phi)$. Hence \ona{Problem~\ref{p:ch05:cmc}}\onp{continuous MaxCut} is $\mathbf{NP}$-hard, and unless $\mathbf P = \mathbf{NP}$ no polynomial-time algorithm approximates it within a ratio better than $\alpha_{\max}$.
\end{lem}
\ona{\emph{Proof.} Fix all the coordinates except $\phi_i$. Since $A$ is symmetric with zero diagonal,}
\begin{equation}
    \mu(\boldsymbol\phi) \;=\; \cos\phi_i\,\Big[\tfrac12\sum_{j\ne i}A_{ij}\cos\phi_j\Big] \;+\; C_i,
\end{equation}
\ona{where $C_i$ does not depend on $\phi_i$. As a function of $\phi_i$, therefore, $\mu$ is affine in $\cos\phi_i \in [-1,1]$, and attains its least value at $\phi_i = 0$ or at $\phi_i = \pi$. Moving the coordinates one at a time to the better of $0$ and $\pi$ never increases $\mu$, and after $d$ moves $\boldsymbol\phi \in \{0,\pi\}^d$. There $v_i = \cos\phi_i \in \{\pm1\}$, and}
\begin{equation}
    \mu \;=\; \tfrac14\sum_{i,j}A_{ij}(v_iv_j - 1) \;=\; -\tfrac12\sum_{i,j}A_{ij}\,\mathbb 1[v_i \ne v_j] \;=\; -\,\mathrm{cut}(v). \qquad\square
\end{equation}
\ona{The reduction preserves values exactly, and continuous MaxCut is therefore exactly as hard to approximate as MaxCut. The rounding loses nothing: the cut read off from a continuous point is never worse than the point itself.}
\end{frame}

\begin{frame}\ft{Rounding without loss}
\ona{The proof of Lemma~\ref{lem:ch05:cmc} rests on a single observation, illustrated in Fig.~\ref{fig:ch05:affine}: along each coordinate, $\mu$ is an affine function of $\cos\phi_i$, and an affine function on an interval attains its minimum at an endpoint. The two lines correspond to the two moves of the rounding shown in Fig.~\ref{fig:ch05:landscape}, for the 4-cycle with $\phi_3 = 0$ and $\phi_4 = \pi$, starting from $(\phi_1,\phi_2) = (4\pi/3,4\pi/3)$.}
\ona{\gdef\afH{0.42\linewidth}}\onp{\gdef\afH{0.36\linewidth}}
\ona{\begin{figure}[htb]}\begin{center}
\begin{tikzpicture}
\begin{axis}[width=0.48\linewidth,height=\afH,clip=false,axis lines=left,xmin=-1.15,xmax=1.15,ymin=-4.3,ymax=-1.7,xtick={-1,-0.5,0,0.5,1},ytick={-4,-3.5,-3,-2.5,-2},xlabel={$\cos\phi_i$},ylabel={$\mu$},domain=-1:1,samples=2,no markers,thick]
  \addplot[blue!70!black] {-0.75*x-2.75};
  \addplot[red!70!black] {x-3};
  \draw[-{Latex},very thick,blue!70!black,opacity=0.5] (axis cs:-0.5,-2.375) -- (axis cs:0.96,-3.47);
  \draw[-{Latex},very thick,red!70!black,opacity=0.5] (axis cs:-0.5,-3.5) -- (axis cs:-0.96,-3.96);
  \fill[blue!70!black] (axis cs:-0.5,-2.375) circle (2pt);
  \fill[blue!70!black] (axis cs:1,-3.5) circle (2pt);
  \fill[red!70!black] (axis cs:-0.5,-3.5) circle (2pt);
  \fill[red!70!black] (axis cs:-1,-4) circle (2.5pt);
  \node[note,anchor=south] (n4) at (axis cs:-0.3,-1.6) {affine in $\cos\phi_i$, and so\\least at an endpoint};
  \draw[notearrow] (n4.south) to[bend right=10] (axis cs:-0.78,-2.17);
  \node[note,anchor=west] (n1) at (axis cs:1.3,-2.5) {first $\phi_1$, with $\cos\phi_2 = -\frac12$\\fixed: from $\cos\phi_1 = -\frac12$ to $1$};
  \draw[notearrow] (n1.west) to[bend right=15] (axis cs:0.35,-3.01);
  \node[note,anchor=west] (n2) at (axis cs:1.3,-3.45) {then $\phi_2$, with $\phi_1 = 0$\\fixed: to $\cos\phi_2 = -1$};
  \draw[notearrow] (n2.west) to[bend left=12] (axis cs:-0.2,-3.2);
  \node[note,anchor=west] (n3) at (axis cs:1.3,-4.15) {$\mu = -4$, the maximum cut};
  \draw[notearrow] (n3.west) to[bend left=8] (axis cs:-0.94,-4.04);
\end{axis}
\end{tikzpicture}
\end{center}
\ona{\caption{The rounding of Lemma~\ref{lem:ch05:cmc}. In blue, $\mu$ as a function of $\cos\phi_1$, with $\cos\phi_2 = -\frac12$; in red, as a function of $\cos\phi_2$, with $\phi_1 = 0$. Each is a straight line, so that moving the coordinate to the better endpoint, $\phi_i\in\{0,\pi\}$, never increases $\mu$; two moves reach $-4$, the maximum cut of the 4-cycle.}\label{fig:ch05:affine}\end{figure}}
\onp{\footnotesize Along each coordinate $\mu$ is a straight line in $\cos\phi_i$, and the better endpoint is never worse.}
\end{frame}

\begin{frame}\ft{The landscape of continuous MaxCut}
\ona{Figure~\ref{fig:ch05:landscape} shows $\mu$ for the 4-cycle $v_1v_2v_3v_4$ with $\phi_3 = 0$ and $\phi_4 = \pi$ fixed, in which case
\begin{equation*}
    \mu(\phi_1,\phi_2) = \tfrac12\big(\cos\phi_1\cos\phi_2 + \cos\phi_2 - \cos\phi_1 - 5\big).
\end{equation*}}
\ona{\gdef\lsH{0.48\linewidth}\gdef\lsW{0.52\linewidth}\gdef\lsC{-0.2}}\onp{\gdef\lsH{0.36\linewidth}\gdef\lsW{0.44\linewidth}\gdef\lsC{-0.3}}
\ona{\begin{figure}[htb]}\begin{center}
\begin{tikzpicture}
\begin{axis}[name=ls,width=\lsW,height=\lsH,view={0}{90},xmin=0,xmax=360,ymin=0,ymax=360,xtick={0,90,180,270,360},xticklabels={$0$,$\frac\pi2$,$\pi$,$\frac{3\pi}2$,$2\pi$},ytick={0,90,180,270,360},yticklabels={$0$,$\frac\pi2$,$\pi$,$\frac{3\pi}2$,$2\pi$},xlabel={$\phi_1$},ylabel={$\phi_2$},colorbar horizontal,colorbar style={at={(0.5,\lsC)},anchor=north,height=0.18cm,font=\tiny},colormap/viridis,point meta min=-4,point meta max=-2,enlargelimits=false,axis on top]
  \addplot3[surf,shader=interp,domain=0:360,y domain=0:360,samples=41] {0.5*(cos(x)*cos(y)+cos(y)-cos(x)-5)};
  \draw[-{Latex},very thick,white] (axis cs:240,240,0) -- (axis cs:356,240,0);
  \draw[-{Latex},very thick,white] (axis cs:356,240,0) -- (axis cs:356,184,0);
  \fill[white] (axis cs:240,240,0) circle (2.5pt);
  \node[white,font=\tiny,above] at (axis cs:240,240,0) {start};
  \fill[red] (axis cs:0,180,0) circle (3pt);
  \fill[red] (axis cs:360,180,0) circle (3pt);
  \coordinate (a1) at (axis cs:298,240,0);
  \coordinate (a2) at (axis cs:356,212,0);
  \coordinate (mn) at (axis cs:360,180,0);
\end{axis}
  \node[note,anchor=west] (n1) at ([xshift=0.4cm,yshift=0.95cm]ls.east) {first $\phi_1$,\\with $\phi_2$ fixed};
  \draw[notearrow] (n1.west) to[bend right=10] (a1);
  \node[note,anchor=west] (n2) at ([xshift=0.4cm,yshift=0.2cm]ls.east) {then $\phi_2$, with\\$\phi_1 = 2\pi\equiv0$};
  \draw[notearrow] (n2.west) to[bend right=10] (a2);
  \node[note,anchor=west] (n3) at ([xshift=0.4cm,yshift=-0.7cm]ls.east) {the minimum $-4$,\\the maximum cut};
  \draw[notearrow] (n3.west) to[bend left=10] (mn);
\end{tikzpicture}
\end{center}
\ona{\caption{Continuous MaxCut on the 4-cycle with $\phi_3 = 0$ and $\phi_4 = \pi$ fixed. The minimum $-4$ (red) is attained at $(\phi_1,\phi_2) = (0,\pi)$, which corresponds to the maximum cut of the 4-cycle. The white arrows show the rounding of the proof: first $\phi_1$ is moved to the better of $0$ and $\pi$, here $2\pi\equiv0$, with $\phi_2$ fixed, and then $\phi_2$; neither move increases $\mu$.}\label{fig:ch05:landscape}\end{figure}}
\onp{\footnotesize The minimum $-4 = -\mathrm{MaxCut}(C_4)$ is at $(0,\pi)$; the arrows show the rounding, one coordinate at a time.}
\end{frame}

\begin{frame}\ft{Local minima are the local optima of MaxCut}
\onp{\footnotesize}
\ona{The proof of Lemma~\ref{lem:ch05:cmc} tells us more about the landscape than the location of its minimum. The derivatives of $\mu$ are}
\begin{equation}
    \frac{\partial\mu}{\partial\phi_i} = -\tfrac12\sin\phi_i\sum_{j}A_{ij}\cos\phi_j,
    \qquad
    \frac{\partial^2\mu}{\partial\phi_i\,\partial\phi_k} = -\tfrac12\,\delta_{ik}\cos\phi_i\sum_jA_{ij}\cos\phi_j + \tfrac12A_{ik}\sin\phi_i\sin\phi_k.
\end{equation}
\ona{At a point of $\{0,\pi\}^d$, with $v_i = \cos\phi_i\in\{\pm1\}$, the gradient vanishes and the Hessian is diagonal, with entries $-\frac12v_i\sum_jA_{ij}v_j$, that is, half the number of neighbours of $i$ on the other side of the cut less half the number on its own side. The entry is positive exactly when moving $i$ to the other side would decrease the cut. It follows that:}
\begin{itemize}\setlength\itemsep{0.15em}
    \item the points of $\{0,\pi\}^d$ that are strict local minima of $\mu$ are the cuts made worse by every single-vertex move, which are the fixed points of greedy local search;
    \item such a cut contains at least half of the edges at every vertex, and hence at least $|E|/2\ge\mathrm{MC}/2$ edges, and the bound $1/2$ is tight \ona{(Exercise~\ref{exe:ch05:persistent})}\onp{(exercise)};
    \item a gradient method, or any other method that converges to a local minimum, such as Newton's method, a quasi-Newton algorithm such as Broyden--Fletcher--Goldfarb--Shanno (BFGS), or the natural gradient, may, if started in the basin of such a minimum, return a cut of half the optimum; its ratio is then $1/2$, against $0.878$ for Goemans--Williamson.
\end{itemize}
\ona{\citet{bittel2021training} call these minima \emph{persistent}. They are not artefacts of a poor parametrisation, to be removed by a better ansatz, but the local optima of the combinatorial problem itself, embedded in the landscape. Which of them a local method finds depends on where it starts (Figs.~\ref{fig:ch05:persistent} and~\ref{fig:ch05:localminima}), and the choice of starting point is the subject of \Chaps~\chref{C.WSRounded}--\chref{C.WSContinuous}.}
\end{frame}

\begin{frame}\ft{A persistent local minimum}
\ona{Figure~\ref{fig:ch05:persistent} shows the graph of Exercise~\ref{exe:ch05:persistent} for $k = 2$, which has eight vertices and twelve edges. In the cut on the left, every vertex has two neighbours across the cut and one on its own side, so that every single-vertex move makes the cut worse, and the point of $\{0,\pi\}^8$ that encodes the cut is a strict local minimum of $\mu$. The graph is nevertheless bipartite, and the cut on the right contains all twelve edges. As $k$ grows, the ratio $k/(2k-1)$ attained by the local minimum tends to $1/2$.}
\ona{\begin{figure}[htb]}\begin{center}
\begin{tikzpicture}[every node/.style={font=\scriptsize},lv/.style={circle,fill=black,inner sep=1.7pt},rv/.style={circle,draw,thick,fill=white,inner sep=1.5pt},ce/.style={green!45!black,thick},ue/.style={red!70!black,thick,dashed}]
  \fill[blue!7,rounded corners=6pt] (-0.35,0.8) rectangle (3.35,1.8);
  \fill[orange!12,rounded corners=6pt] (-0.35,-1.0) rectangle (3.35,0.0);
  \node[anchor=east] at (-0.4,1.3) {side $S$};
  \node[anchor=east] at (-0.4,-0.5) {side $T$};
  \node[lv] (a1) at (0,1.3) {}; \node[lv] (a2) at (0.8,1.3) {}; \node[rv] (b1) at (2.2,1.3) {}; \node[rv] (b2) at (3.0,1.3) {};
  \node[lv] (a3) at (0,-0.5) {}; \node[lv] (a4) at (0.8,-0.5) {}; \node[rv] (b3) at (2.2,-0.5) {}; \node[rv] (b4) at (3.0,-0.5) {};
  \foreach \u in {a1,a2} \foreach \v in {b3,b4} \draw[ce] (\u) -- (\v);
  \foreach \u in {a3,a4} \foreach \v in {b1,b2} \draw[ce] (\u) -- (\v);
  \draw[ue] (a2) -- (b1); \draw[ue] (a1) to[bend left=25] (b2);
  \draw[ue] (a4) -- (b3); \draw[ue] (a3) to[bend right=25] (b4);
  \node at (1.5,-1.4) {(a) a strict local minimum: $8$ of $12$ edges};
  \begin{scope}[xshift=5.9cm]
    \fill[gray!10,rounded corners=6pt] (-0.35,-0.85) rectangle (0.35,1.95);
    \fill[gray!10,rounded corners=6pt] (1.85,-0.85) rectangle (2.55,1.95);
    \node[lv] (c1) at (0,1.6) {}; \node[lv] (c2) at (0,0.9) {}; \node[lv] (c3) at (0,0.2) {}; \node[lv] (c4) at (0,-0.5) {};
    \node[rv] (d1) at (2.2,1.6) {}; \node[rv] (d2) at (2.2,0.9) {}; \node[rv] (d3) at (2.2,0.2) {}; \node[rv] (d4) at (2.2,-0.5) {};
    \foreach \u/\v in {c1/d2,c2/d1,c3/d4,c4/d3,c1/d3,c1/d4,c2/d3,c2/d4,c3/d1,c3/d2,c4/d1,c4/d2} \draw[ce] (\u) -- (\v);
    \node at (1.1,-1.4) {(b) the bipartition: all $12$ edges};
  \end{scope}
  \node[note,anchor=south] (nv) at (1.5,2.2) {each vertex has two neighbours across and one on its\\own side, so every single-vertex move loses an edge};
  \draw[notearrow] (nv.south) to[bend right=15] (a2);
  \node[note,anchor=west] (nu) at (3.6,0.4) {dashed: the\\four uncut edges};
  \draw[notearrow] (nu.west) to[bend left=15] (3.2,-0.75);
\end{tikzpicture}
\end{center}
\ona{\caption{The graph of Exercise~\ref{exe:ch05:persistent} for $k = 2$, with the vertices of $L_1\cup L_2$ filled and those of $R_1\cup R_2$ hollow; $L_1$ and $R_1$ form the top row of (a), and $L_2$ and $R_2$ the bottom row. (a) The cut $(L_1\cup R_1, L_2\cup R_2)$ is optimal under single-vertex moves, and is therefore a strict local minimum of $\mu$, but it contains only $8$ edges. (b) The same graph, drawn as the bipartite graph that it is; the cut $(L,R)$ contains all $12$ edges.}\label{fig:ch05:persistent}\end{figure}}
\onp{\footnotesize Every single move loses an edge, yet the cut contains only $8$ of $12$ edges; in general the ratio is $k/(2k-1)\to\frac12$.}
\end{frame}

\begin{frame}\ft{The quality of typical local minima}
\ona{The bound $1/2$ is a worst case. Figure~\ref{fig:ch05:localminima} shows a more typical case, a random 3-regular graph on $20$ vertices, with $30$ edges and maximum cut $28$. Gradient descent on $\mu$ from 1000 random points, with rounding at the end, and greedy single-vertex search from 1000 random cuts end at the same values, as they should, since the local minima of $\mu$ are the local optima of the greedy search. About a third of the runs of either method reach the maximum cut; the remainder stop at persistent local minima, some of which contain only $22$ edges. The data were produced by \texttt{code/ch05\_figures.py}.}
\ona{\gdef\lmH{0.4\linewidth}}\onp{\gdef\lmH{0.34\linewidth}}
\ona{\begin{figure}[htb]}\begin{center}
\begin{tikzpicture}
\begin{axis}[ybar,bar width=4.5pt,width=0.56\linewidth,height=\lmH,clip=false,axis lines=left,xmin=21.3,xmax=28.7,ymin=0,ymax=0.4,xtick={22,...,28},ytick={0,0.1,0.2,0.3,0.4},yticklabel style={/pgf/number format/fixed},xlabel={cut value at the end of a run (maximum $28$)},ylabel={fraction of runs},enlargelimits=false]
  \addplot[fill=blue!55,draw=blue!70!black] [restrict x to domain=21.5:28.5] table[x expr=\thisrow{ratio}*28,y=gd] {figures/ch05/local_minima.dat};
  \addplot[fill=orange!60,draw=orange!80!black] [restrict x to domain=21.5:28.5] table[x expr=\thisrow{ratio}*28,y=greedy] {figures/ch05/local_minima.dat};
  \draw[dashed,thick,green!40!black] (axis cs:24.6,0) -- (axis cs:24.6,0.4);
  \node[note,anchor=west] (n4) at (axis cs:24.75,0.37) {GW guarantee,\\$0.878\cdot28$};
  \node[note,anchor=west] (n1) at (axis cs:29.0,0.37) {$35\%$ of the runs\\reach the maximum cut};
  \draw[notearrow] (n1.west) to[bend right=10] (axis cs:28.1,0.355);
  \node[note,anchor=west] (n2) at (axis cs:29.0,0.2) {blue: gradient descent on $\mu$;\\orange: greedy search};
  \draw[notearrow] (n2.west) to[bend left=10] (axis cs:27.1,0.15);
  \node[note,anchor=south] (n3) at (axis cs:22.6,0.1) {persistent local\\minima, as low\\as $22$ edges};
  \draw[notearrow] (n3.south) -- (axis cs:22,0.025);
\end{axis}
\end{tikzpicture}
\end{center}
\ona{\caption{The values at which local methods stop on continuous MaxCut, for a random 3-regular graph on $20$ vertices with maximum cut $28$: the fraction of 1000 runs of gradient descent on $\mu$ (blue) and of greedy single-vertex search (orange) that end at each cut value. The two distributions agree. The dashed line marks the expected value guaranteed by Goemans--Williamson on every instance.}\label{fig:ch05:localminima}\end{figure}}
\onp{\footnotesize Gradient descent on $\mu$ and greedy search stop at the same persistent local minima, and a third of the runs reach the optimum.}
\end{frame}

\begin{frame}\ft{Step 2: a circuit that computes $\mu$}
\onp{\footnotesize}
\ona{\emph{Proof of Proposition~\ref{propn:ch05:oracle}.} We reduce Problem~\ref{p:ch05:cmc} to Problem~\ref{p:ch05:oracle}. Take $N = d$ qubits and $L = d$ layers with generators $H_i = \frac12Y_i$, so that $U_i(\phi_i) = R_Y(\phi_i)$ acts on qubit $i$, together with the Ising observable}
\begin{equation}
    O \;=\; \tfrac14\sum_{i,j}A_{ij}\big(Z_iZ_j - \one\big).
\end{equation}
\ona{The state is the product state $\bigotimes_i\big(\cos\tfrac{\phi_i}2\ket0 + \sin\tfrac{\phi_i}2\ket1\big)$, for which $\braket{Z_i} = \cos\phi_i$ and $\braket{Z_iZ_j} = \cos\phi_i\cos\phi_j$ when $i\ne j$, and hence}
\begin{equation}
    \braket{O(\boldsymbol\phi)} \;=\; \mu(\boldsymbol\phi) \qquad\text{exactly.}
\end{equation}
\ona{The instance has $d$ generators and $\mathcal O(d^2)$ Pauli strings, and the oracle can be simulated, since $\mu(\boldsymbol\phi)$ is a sum of $\mathcal O(d^2)$ products of cosines. A polynomial-time algorithm for Problem~\ref{p:ch05:oracle}, with the oracle replaced by this computation and followed by the rounding of Lemma~\ref{lem:ch05:cmc}, would therefore solve MaxCut in polynomial time. \qed}
\begin{itemize}\setlength\itemsep{0.15em}
    \item The circuit creates no entanglement, and the ground state of $O$, a computational basis state, lies in the ansatz, so that $\braket{O}_{\min} = \lambda_{\min}(O) = -\mathrm{MC}$ and $\delta_m = 0$. The hardness lies entirely in the classical optimization.
    \item The eigenvalues of $O$ are the numbers $-\mathrm{cut}(z)$, $z\in\{0,1\}^d$, so that $w(O) = \mathrm{MC}$, and a training algorithm with optimization error $\Delta$ is a MaxCut algorithm with ratio $1-\Delta$. Unless $\mathbf P = \mathbf{NP}$, therefore, $\Delta\ge1-\alpha_{\max}$.
\end{itemize}
\end{frame}

\begin{frame}\ft{Boosting, and the absence of any guarantee}
\onp{\footnotesize}
\begin{cor}[\citet{bittel2021training}]\label{cor:ch05:boost}
Unless $\mathbf P = \mathbf{NP}$, no polynomial-time algorithm for \ona{Problem~\ref{p:ch05:oracle}}\onp{the oracular problem} has optimization error $\Delta<1$.
\end{cor}
\ona{\emph{Proof.} Take $k$ copies of the circuit of step 2, on $kd$ qubits, all controlled by the same parameters, $\tilde U(\boldsymbol\phi) = U(\boldsymbol\phi)^{\otimes k}$, and the observable}
\begin{equation}
    \tilde O = (-1)^{k-1}O^{\otimes k}, \qquad \braket{\tilde O(\boldsymbol\phi)} = (-1)^{k-1}\mu(\boldsymbol\phi)^k = -|\mu(\boldsymbol\phi)|^k,
\end{equation}
\ona{where we have used $\mu\le0$; the sign ensures that the problem remains a minimisation. For constant $k$, $\tilde O$ has $\mathcal O(d^{2k})$ Pauli strings. Its eigenvalues are $(-1)^{k-1}$ times products of $k$ eigenvalues of $O$, which lie in $[-\mathrm{MC},0]$; they therefore lie in $[-\mathrm{MC}^k,0]$, and $w(\tilde O) = \mathrm{MC}^k = -\braket{\tilde O}_{\min}$. Suppose that an algorithm guarantees $\Delta\le\rho<1$. On this instance it returns some $\boldsymbol\phi$ with}
\begin{equation}
    \mathrm{MC}^k - |\mu(\boldsymbol\phi)|^k \;\le\; \rho\,\mathrm{MC}^k, \qquad\text{that is,}\qquad |\mu(\boldsymbol\phi)| \;\ge\; (1-\rho)^{1/k}\,\mathrm{MC}.
\end{equation}
\ona{Choose $k$ so large that $(1-\rho)^{1/k} > 16/17$; for $\rho = 0.9$, $k = 38$ suffices. Rounding $\boldsymbol\phi$ as in Lemma~\ref{lem:ch05:cmc} then gives, in polynomial time, a cut within the factor $16/17$ of the optimum, and this is $\mathbf{NP}$-hard \citep{Hstad2001}. \qed}
\ona{The $k$-th power stretches the scale: a relative error $\epsilon$ in $|\mu|$ becomes a relative error $1-(1-\epsilon)^k\approx k\epsilon$ in $|\mu|^k$. For general VQAs, therefore, the worst-case guarantee of a training algorithm can be made as poor as one pleases.}
\onp{With $k$ copies, an error $\rho$ gives a cut of ratio $(1-\rho)^{1/k}\to1$, and to exceed $16/17$ is $\mathbf{NP}$-hard.}
\end{frame}

\begin{frame}\ft{The number of copies required}
\ona{Figure~\ref{fig:ch05:boost} shows the ratio $(1-\rho)^{1/k}$ that an algorithm with $\Delta\le\rho$ would guarantee on $k$ copies. Every curve eventually crosses the threshold $16/17$ of H\accent23 astad. A weaker guarantee merely requires more copies, but a fixed number of them for each $\rho$, so that the instance remains of polynomial size.}
\ona{\gdef\bsH{0.4\linewidth}}\onp{\gdef\bsH{0.34\linewidth}}
\ona{\begin{figure}[htb]}\begin{center}
\begin{tikzpicture}
\begin{axis}[width=0.56\linewidth,height=\bsH,clip=false,axis lines=left,xmin=0,xmax=100,ymin=0.3,ymax=1.0,xlabel={number of copies $k$},ylabel={guaranteed cut ratio},no markers,thick,samples=150]
  \addplot[blue!70!black,domain=1:100] {0.5^(1/x)};
  \addplot[red!70!black,domain=2:100] {0.1^(1/x)};
  \addplot[orange!80!black,domain=4:100] {0.01^(1/x)};
  \addplot[gray,dashed,domain=0:100] {16/17};
  \fill[blue!70!black] (axis cs:12,0.9439) circle (1.8pt);
  \fill[red!70!black] (axis cs:38,0.9412) circle (1.8pt);
  \fill[orange!80!black] (axis cs:76,0.9412) circle (1.8pt);
  \node[font=\scriptsize,anchor=south east] at (axis cs:99,0.31) {{\color{blue!70!black}$\rho = 0.5$}\quad{\color{red!70!black}$\rho = 0.9$}\quad{\color{orange!80!black}$\rho = 0.99$}};
  \node[note,anchor=north west] (n1) at (axis cs:40,0.8) {for $\rho = 0.9$, $38$ copies give a ratio\\above $16/17$, which is $\mathbf{NP}$-hard};
  \draw[notearrow] (n1.north west) to[bend left=15] (axis cs:37.5,0.93);
  \node[note,anchor=north west] (n3) at (axis cs:60,0.55) {for $\rho = 0.99$, $76$ copies};
  \draw[notearrow] (n3.north) to[bend right=15] (axis cs:76,0.93);
  \node[note,anchor=west] (n2) at (axis cs:103,0.9412) {$16/17$, the threshold\\of H\accent23 astad};
  \draw[notearrow] (n2.west) -- (axis cs:100.5,0.9412);
\end{axis}
\end{tikzpicture}
\end{center}
\ona{\caption{Boosting. An algorithm that guarantees optimization error $\rho$ on the $k$-fold product of step 2 yields a cut within the ratio $(1-\rho)^{1/k}$ of the optimum. For every $\rho<1$ this ratio exceeds $16/17$, beyond which approximating MaxCut is $\mathbf{NP}$-hard \citep{Hstad2001}, at a finite $k$, namely at $12$, $38$ and $76$ copies for $\rho = 0.5$, $0.9$ and $0.99$.}\label{fig:ch05:boost}\end{figure}}
\onp{\footnotesize For every $\rho<1$ the curve crosses $16/17$ at a finite $k$.}
\end{frame}

\begin{frame}\ft{Step 3: logarithmically many qubits}
\onp{\scriptsize}
\ona{The circuit of step 2 acts on $d$ qubits, and it is only its special structure that makes its energy easy to compute. To show that the hardness does not depend on this, \citet{bittel2021training} encode Problem~\ref{p:ch05:cmc} in a Hilbert space of dimension $2d$, in which every quantity is given by a small matrix computation.}
\begin{thm}[\citet{bittel2021training}]\label{thm:ch05:log}
\ona{Problem~\ref{p:ch05:vqa}}\onp{VQA minimisation} is $\mathbf{NP}$-hard already in dimension $n = 2d$ with $L = d$ layers, and its decision version is $\mathbf{NP}$-complete.
\end{thm}
\ona{\emph{Proof.} Take the uniform initial state $\ket{\Psi_0} = \ket{+_{2d}} = (2d)^{-1/2}\sum_{k=1}^{2d}\ket k$ and the diagonal generators $H_j = \ket{2j-1}\!\bra{2j-1} - \ket{2j}\!\bra{2j}$, which impose opposite phases on the $j$-th pair of basis states,}
\begin{equation}
    \ket{\Psi(\boldsymbol\phi)} = \frac1{\sqrt{2d}}\sum_{j=1}^d\big(e^{-i\phi_j}\ket{2j-1} + e^{i\phi_j}\ket{2j}\big).
\end{equation}
\ona{For the observable, double every vertex. Let $W = A\otimes J_2$, where $J_2 = \begin{psmallmatrix}1&1\\1&1\end{psmallmatrix}$, be the adjacency matrix of the graph on $2d$ vertices in which both copies of $i$ are joined to both copies of $j$ whenever $ij$ is an edge, let $L_W = \mathrm{diag}(W\mathbf 1) - W$ be its Laplacian, and let}
\begin{equation}
    O \;=\; -\tfrac d8\,L_W \;=\; \tfrac d8\,A\otimes J_2 \;-\; \tfrac d4\,\mathrm{diag}(A\mathbf 1)\otimes\one_2.
\end{equation}
\ona{This is the observable of \citet{bittel2021training}, who write it as $O' = \frac d8A\otimes J_2$ with the diagonal replaced by minus the column sums. A Laplacian satisfies $x^\dagger L_Wx = \frac12\sum_{k,l}W_{kl}|x_k-x_l|^2$. The amplitudes of $\ket{\Psi(\boldsymbol\phi)}$ are $(2d)^{-1/2}e^{-is\phi_j}$ with $s = \pm1$, and for an edge $ij$ the four pairs of signs give $\sum_{s,s'}|e^{-is\phi_i} - e^{-is'\phi_j}|^2 = 8 - 4\cos(\phi_i-\phi_j) - 4\cos(\phi_i+\phi_j)$, whence}
\begin{equation}
    \braket{O(\boldsymbol\phi)} = \tfrac18\sum_{i,j}A_{ij}\big(\cos(\phi_i-\phi_j) + \cos(\phi_i+\phi_j) - 2\big) = \mu(\boldsymbol\phi).
\end{equation}
\ona{This reduces Problem~\ref{p:ch05:cmc} to Problem~\ref{p:ch05:vqa}. For the decision version, a parameter vector is a certificate that can be checked in polynomial time. \qed}
\onp{The observable $O = -\frac d8L_W$ is a negative multiple of the Laplacian of the doubled graph, and $\braket{O(\boldsymbol\phi)} = \mu(\boldsymbol\phi)$.}
\end{frame}

\begin{frame}\ft{The doubled graph}
\ona{Figure~\ref{fig:ch05:doubled} shows the construction of step 3 for the triangle. All the amplitudes of $\ket{\Psi(\boldsymbol\phi)}$ have the same modulus, so that the state is described by phasors, one for each vertex of the doubled graph, the two copies of vertex $j$ carrying conjugate phases. The energy is a negative multiple of the Laplacian form; it therefore decreases when the phasors at the ends of an edge point in different directions, and vanishes when all the phasors agree.}
\ona{\begin{figure}[htb]}\begin{center}
\begin{tikzpicture}[every node/.style={font=\scriptsize},pv/.style={circle,draw,thick,fill=white,minimum size=0.6cm,inner sep=0pt},ph/.style={-{Stealth[length=3.5pt]},thick,blue!70!black}]
  \foreach \n/\x/\y/\a/\p in {1/0/1.4/30/above,2/-1.2/-0.7/150/left,3/1.2/-0.7/260/right} {
    \node[pv,label={[font=\tiny]\p:$\n$}] (v\n) at (\x,\y) {};
    \draw[ph] (v\n.center) -- ++(\a:0.22);
  }
  \draw[thick] (v1) -- (v2); \draw[thick] (v2) -- (v3); \draw[thick] (v1) -- (v3);
  \node at (0,-1.55) {(a) the triangle, with angles $\phi_j$};
  \begin{scope}[xshift=5.0cm]
    \foreach \n/\x/\y/\a in {1p/0.4/1.4/30,1m/-0.4/1.4/-30,2p/-1.0/-0.95/150,2m/-1.8/-0.5/-150,3p/1.0/-0.95/260,3m/1.8/-0.5/-260} {
      \node[pv] (w\n) at (\x,\y) {};
      \draw[ph] (w\n.center) -- ++(\a:0.22);
    }
    \foreach \i/\j in {1/2,2/3,1/3} \foreach \s in {p,m} \foreach \t in {p,m} \draw[gray!70!black] (w\i\s) -- (w\j\t);
    \node[font=\tiny] at (-0.4,1.85) {$1^-$}; \node[font=\tiny] at (0.4,1.85) {$1^+$};
    \node[font=\tiny] at (-2.2,-0.5) {$2^-$}; \node[font=\tiny] at (-1.0,-1.38) {$2^+$};
    \node[font=\tiny] at (2.2,-0.5) {$3^-$}; \node[font=\tiny] at (1.0,-1.38) {$3^+$};
    \node at (0,-1.85) {(b) the doubled graph, $W = A\otimes J_2$};
  \end{scope}
  \node[note,anchor=west] (np) at (6.9,1.7) {the copies $j^\pm$ of vertex $j$ have\\amplitudes $(2d)^{-1/2}e^{\pm i\phi_j}$};
  \draw[notearrow] (np.west) -- (5.7,1.45);
  \node[note,anchor=west] (ne) at (7.3,0.35) {an edge $ij$ becomes\\four edges};
  \draw[notearrow] (ne.west) to[bend right=10] (5.62,0.35);
  \node[note,anchor=north] (nE) at (4.4,-2.25) {the energy $-\frac d{16}\sum_{kl}W_{kl}|x_k-x_l|^2$ is lowered by phasors that point apart across an edge;\\equal phasors give energy $0$, since the uniform state lies in the kernel of $O$};
\end{tikzpicture}
\end{center}
\ona{\caption{Step 3 for the triangle, $d = 3$. (a) The graph, with the angle $\phi_j$ at each vertex shown as a phasor. (b) The state $\ket{\Psi(\boldsymbol\phi)}$ on the doubled graph, with $2d = 6$ basis states; each vertex $j$ has two copies $j^\pm$ with conjugate phases, and each edge of the graph joins both copies of one end to both copies of the other. The observable is $O = -\frac d8L_W$.}\label{fig:ch05:doubled}\end{figure}}
\onp{\footnotesize Pairs of conjugate phasors on the doubled graph; the energy of $O = -\frac d8L_W$ is lowered by phasors that point apart.}
\end{frame}

\begin{frame}\ft{An easy ground-state problem and hard training}
\onp{\footnotesize}
\ona{The construction of step 3 separates the two parts of the error \eqref{eq:ch05:split} sharply. Since $O = -\frac d8L_W$, and a Laplacian is positive semidefinite with the constant vector in its kernel, $\lambda_{\max}(O) = 0$ and $\lambda_{\min}(O) = -\frac d8\lambda_{\max}(L_W)$. The vectors $v\otimes(1,1)$ and $v\otimes(1,-1)$ show that the spectrum of $L_W$ consists of twice the eigenvalues of the Laplacian $L_A = \mathrm{diag}(A\mathbf 1) - A$ of the graph together with twice its degrees; since $\lambda_{\max}(L_A)\ge e_i^\top L_Ae_i = \deg(i)$, the largest of these is $2\lambda_{\max}(L_A)$. Hence}
\begin{equation}
    \lambda_{\min}(O) \;=\; -\tfrac d4\,\lambda_{\max}(L_A) \;\le\; -\mathrm{MC} \;=\; \braket{O}_{\min},
\end{equation}
\ona{where the inequality is the eigenvalue bound on MaxCut of \citet{mohar1990eigenvalues}: for $v\in\{\pm1\}^d$, $\mathrm{cut}(v) = \frac14v^\top L_Av\le\frac d4\lambda_{\max}(L_A)$.}
\begin{itemize}\setlength\itemsep{0.15em}
    \item The ground-state problem for $O$ amounts to computing the largest eigenvalue of a $d\times d$ matrix, which can be done in polynomial time; its value is an upper bound on MaxCut, which is usually not attained.
    \item The ansatz reaches only states whose amplitudes have equal moduli and opposite phases in each pair. Its least energy is $-\mathrm{MC}$, and the model mismatch $\delta_m$ is therefore the gap between the eigenvalue bound and the maximum cut, relative to the bound.
    \item The optimization within the ansatz is $\mathbf{NP}$-hard. Restricting the states has made the problem harder, not easier.
\end{itemize}
\end{frame}

\begin{frame}\ft{The example of the 5-cycle}
\ona{Figure~\ref{fig:ch05:c5} shows the spectrum of the observable of step 3 for the 5-cycle, $d = 5$. The smallest eigenvalue is the eigenvalue bound $-\frac54\lambda_{\max}(L_{C_5}) = -\frac54\big(2 + 2\cos\frac\pi5\big)\approx-4.52$, whereas the states of the ansatz reach only $-\mathrm{MaxCut}(C_5) = -4$. The model mismatch is thus $0.52/4.52\approx0.12$ of the spectral width, and the remainder of the error is due to the optimizer.}
\ona{\begin{figure}[htb]}\begin{center}
\begin{tikzpicture}[x=1.65cm,every node/.style={font=\scriptsize}]
  \fill[green!14] (-4,0) rectangle (0,0.3);
  \draw[thick,-{Latex}] (-4.85,0) -- (0.55,0) node[right] {energy};
  \foreach \x in {-4,-3,-2,-1,0} \draw (\x,0.06) -- (\x,-0.06) node[below] {$\x$};
  \foreach \e/\k in {-4.5225/2,-2.5/5,-1.7275/2,0/1} \foreach \j in {1,...,\k} \fill[blue!70!black] (\e,0.42+0.17*\j) circle (1.7pt);
  \node[red!70!black,fill,rectangle,rotate=45,inner sep=2pt] at (-4,0) {};
  \draw[decorate,decoration={brace,amplitude=4pt,mirror}] (-4.5225,-0.42) -- (-4,-0.42) node[midway,below=5pt] {mismatch};
  \node[note,anchor=south west] (n1) at (-4.85,1.45) {the ground energy $-\frac54\lambda_{\max}(L_{C_5})\approx-4.52$, an eigenvalue,\\is the bound of Mohar and Poljak, computable in polynomial time};
  \draw[notearrow] ([xshift=0.6cm]n1.south west) to[bend right=10] (-4.52,0.84);
  \node[note,anchor=west] (n2) at (-3.3,-0.75) {the best state of the ansatz has energy $-\mathrm{MaxCut}(C_5) = -4$,\\which is $\mathbf{NP}$-hard to find for general graphs};
  \draw[notearrow] (n2.west) to[bend left=15] (-3.95,-0.08);
  \node[note,anchor=east] (n3) at (-2.68,0.95) {the spectrum of $O$,\\with multiplicities};
  \draw[notearrow] (n3.east) -- (-2.56,0.95);
  \node[note,anchor=west] (n4) at (-1.6,0.95) {in green, energies\\in the ansatz};
  \draw[notearrow] (n4.south) to[bend left=10] (-0.95,0.24);
  \node[note,anchor=west] (n5) at (0.3,0.62) {the uniform state, at $0$};
  \draw[notearrow] (n5.west) -- (0.06,0.59);
\end{tikzpicture}
\end{center}
\ona{\caption{The observable of step 3 for the 5-cycle. The dots are its eigenvalues $-\frac58\cdot2\lambda$, where $\lambda$ runs over the eigenvalues $0$, $2-2\cos\frac{2\pi}5$, $2+2\cos\frac\pi5$ of the Laplacian of $C_5$ and over its degrees, stacked by multiplicity. In green, the interval $[-4,0]$ of energies of states in the ansatz. The ground energy is easy to compute but cannot be reached, and the least energy that can be reached is hard to find.}\label{fig:ch05:c5}\end{figure}}
\onp{\footnotesize The ground energy $-4.52$ is an eigenvalue and easy to compute; the best energy in the ansatz, $-4$, is $\mathbf{NP}$-hard to find in general.}
\end{frame}

\begin{frame}\ft{Step 4: a clock with $d$ hands}
\onp{\footnotesize}
\ona{With $d$ layers, each parameter sets one phase. A single layer has a single parameter, a time $t$, and $e^{-itH}$ with $H = \sum_jE_j(\ket{2j-1}\!\bra{2j-1} - \ket{2j}\!\bra{2j})$ sets all the phases at once, to $\phi_j = E_jt$. The question is whether, for energies with short descriptions, the curve $t\mapsto\boldsymbol Et$ comes close to every point of the torus $[0,2\pi)^d$.}
\begin{defn}[\citet{bittel2021training}]
Energies $E_1,\dots,E_d$ form an \emph{$\epsilon$-approximate ergodic spectrum} if for every $\boldsymbol\phi\in[0,2\pi)^d$ there is a $t\ge0$ with $|\phi_j - E_jt|_{2\pi}\le\epsilon$ for all $j$, where $|x|_{2\pi} = \min_{k\in\mathbb Z}|x-2\pi k|$.
\end{defn}
\ona{By Kronecker's theorem, energies that are linearly independent over the rationals form an $\epsilon$-approximate ergodic spectrum for every $\epsilon>0$, but the times required may be very large, and such energies cannot be written down exactly. A clock serves better.}
\begin{lem}[\citet{bittel2021training}]\label{lem:ch05:clock}
For an integer $m\ge2$, the energies $E_j = 2\pi/m^j$, $j = 1,\dots,d$, form a $4\pi/m$-approximate ergodic spectrum, with integer times $0\le t<m^d$.
\end{lem}
\ona{\emph{Proof.} Write the target angles in base $m$. Let $s_j = \lfloor m\phi_j/2\pi\rfloor\in\{0,\dots,m-1\}$, and let $t = \sum_{k=1}^ds_km^{k-1}$ be the number with digits $s_d\cdots s_2s_1$. Then}
\begin{equation}
    E_jt \;=\; 2\pi\sum_{k>j}s_km^{k-1-j} \;+\; \frac{2\pi s_j}m \;+\; 2\pi\sum_{k<j}s_km^{k-1-j}.
\end{equation}
\ona{The first sum is a multiple of $2\pi$. By the choice of $s_j$, the middle term is within $2\pi/m$ of $\phi_j$. The last sum is at most $2\pi(m-1)\sum_{k<j}m^{k-1-j}<2\pi/m$. Hence $|\phi_j - E_jt|_{2\pi}<4\pi/m$. \qed}
\ona{Hand $j$ turns once in $m^j$ units of time, $m$ times more slowly than hand $j-1$. For $m = 10$ it points, apart from a small contribution from the lower digits, at the $j$-th digit of $t$ from the right, as on an odometer (Exercise~\ref{exe:ch05:clock}).}
\onp{The digits of $t$ in base $m$ set the hands; the error is less than $2\pi/m$ from rounding plus $2\pi/m$ from the lower digits.}
\end{frame}

\begin{frame}\ft{An example of the clock}
\ona{In Fig.~\ref{fig:ch05:clock} the hands of a clock with $m = 10$ and $d = 3$ are set to the angles $\boldsymbol\phi = 2\pi\,(0.62, 0.15, 0.48)$. The digits $s_j = \lfloor10\phi_j/2\pi\rfloor$ are $6$, $1$ and $4$, so that $t = 416$. Hand $1$ stands at $0.6$ of a turn, within a tenth of a turn of its target; hands $2$ and $3$ stand at their digits, advanced by the lower digits by $0.06$ and $0.016$ of a turn respectively.}
\ona{\begin{figure}[htb]}\begin{center}
\begin{tikzpicture}[every node/.style={font=\scriptsize}]
  \foreach \j/\x/\tgt/\hand/\per/\dig in {1/0/223.2/216/10/6,2/3.4/54/57.6/100/1,3/6.8/172.8/149.76/1000/4} {
    \begin{scope}[xshift=\x cm]
      \draw[thick,fill=gray!5] (0,0) circle (1.0);
      \foreach \k in {0,...,9} \draw (36*\k:0.86) -- (36*\k:1.0);
      \node[font=\tiny] at (0:1.17) {$0$};
      \draw[dashed,gray!60!black,thick] (0,0) -- (\tgt:0.98);
      \draw[annot,very thick] (\hand:0.5) arc[start angle=\hand,end angle=\tgt,radius=0.5];
      \draw[-{Stealth[length=5pt]},very thick,blue!70!black] (0,0) -- (\hand:0.8);
      \fill (0,0) circle (1.5pt);
      \node[align=center] at (0,-1.45) {hand $\j$, digit $s_\j = \dig$\\a turn in $\per$ ticks};
    \end{scope}
  }
  \node[note,anchor=south] (n1) at (-0.2,1.3) {rounding $\phi_1$ to a tenth\\of a turn errs by $0.02$};
  \draw[notearrow] (n1.south) to[bend right=25] (-0.4,-0.31);
  \node[note,anchor=south] (n2) at (3.4,1.3) {dashed, the target $\phi_j$;\\arrow, the hand at $t = 416$};
  \draw[notearrow] (n2.south) to[bend left=15] (3.92,0.74);
  \node[note,anchor=south] (n3) at (7.0,1.3) {the lower digits advance the\\hand by $0.016$ of a turn};
  \draw[notearrow] (n3.south) to[bend left=20] (6.32,0.17);
\end{tikzpicture}
\end{center}
\ona{\caption{The clock of Lemma~\ref{lem:ch05:clock} for $m = 10$ and $d = 3$, at $t = 416$, set to the targets $\boldsymbol\phi = 2\pi\,(0.62,0.15,0.48)$, shown dashed. Hand $j$ turns once in $10^j$ ticks and stands at $E_jt = 2\pi t/10^j$ mod $2\pi$. The teal arcs show the errors, $0.02$, $0.01$ and $0.064$ of a turn, all below the bound $4\pi/m$, which is a fifth of a turn.}\label{fig:ch05:clock}\end{figure}}
\onp{\footnotesize At $t = 416$ the digits $4,1,6$ set hands $3,2,1$, with errors below a fifth of a turn.}
\end{frame}

\begin{frame}\ft{The reduction with one layer}
\onp{\footnotesize}
\begin{thm}[\citet{bittel2021training}]\label{thm:ch05:onelayer}
\ona{Problem~\ref{p:ch05:vqa}}\onp{VQA minimisation} is $\mathbf{NP}$-hard with a single layer, $L = 1$.
\end{thm}
\ona{\emph{Proof.} Keep $\ket{\Psi_0}$ and $O$ from step 3, and use the single generator $H = \sum_jE_j(\ket{2j-1}\!\bra{2j-1} - \ket{2j}\!\bra{2j})$ with the energies of Lemma~\ref{lem:ch05:clock}, so that $\braket{O(t)} = \mu(\boldsymbol Et)$. It remains to determine how fine the clock must be. Since $|\partial\mu/\partial\phi_j|\le\frac12\deg(j)$, and the degrees sum to $2|E|$,}
\begin{equation}
    |\mu(\boldsymbol\phi) - \mu(\boldsymbol\phi')| \;\le\; |E|\,\max_j|\phi_j - \phi'_j|_{2\pi}.
\end{equation}
\ona{Taking for $\boldsymbol\phi$ a maximum cut, and for $t$ the time given by the lemma, we obtain $\min_t\braket{O(t)}\le-\mathrm{MC} + 4\pi|E|/m$. If an algorithm returns a $t$ within $\eta$ of the minimum, rounding $\boldsymbol Et$ as in Lemma~\ref{lem:ch05:cmc} gives a cut of value greater than $\mathrm{MC} - 4\pi|E|/m - \eta$. When $m>8\pi|E|$ and $\eta<1/2$, this is a maximum cut, since cut values are integers. The energies then require $\mathcal O(d\log|E|)$ bits; after a rescaling of $t$ by $2\pi$, they are the rationals $m^{-j}$. \qed}
\begin{cor}[\citet{bittel2021training}]\label{cor:ch05:time}
For a time-independent Hamiltonian on logarithmically many qubits, minimising $\braket{O(t)}$ over the time $t\ge0$ of the evolution is $\mathbf{NP}$-hard.
\end{cor}
\ona{A single parameter does not make the landscape simple. As a function of $t$, $\mu(\boldsymbol Et)$ is a sum of cosines whose frequencies span $d$ orders of magnitude, and it has exponentially many local minima (Fig.~\ref{fig:ch05:onelayer}).}
\end{frame}

\begin{frame}\ft{The curve of a single parameter on the torus}
\ona{Figure~\ref{fig:ch05:torus} shows the curve $t\mapsto(E_1t,E_2t)$ for $m = 5$ on the landscape of a single edge, $\mu(\phi_1,\phi_2) = \frac12(\cos\phi_1\cos\phi_2 - 1)$, whose minimum $-1$ is attained where one angle is $0$ and the other $\pi$. The curve winds around the torus in $m$ parallel strands. Its best point, at $t\approx10.1$, has $\mu\approx-0.91$, and rounding it gives the cut.}
\ona{\gdef\tsH{0.46\linewidth}}\onp{\gdef\tsH{0.36\linewidth}}
\ona{\begin{figure}[htb]}\begin{center}
\begin{tikzpicture}
\begin{axis}[width=0.5\linewidth,height=\tsH,view={0}{90},xmin=0,xmax=360,ymin=0,ymax=360,xtick={0,90,180,270,360},xticklabels={$0$,$\frac\pi2$,$\pi$,$\frac{3\pi}2$,$2\pi$},ytick={0,90,180,270,360},yticklabels={$0$,$\frac\pi2$,$\pi$,$\frac{3\pi}2$,$2\pi$},xlabel={$\phi_1 = E_1t$},ylabel={$\phi_2 = E_2t$},colorbar,colorbar style={width=0.25cm,font=\tiny},colormap/viridis,point meta min=-1,point meta max=0,enlargelimits=false,axis on top]
  \addplot3[surf,shader=interp,domain=0:360,y domain=0:360,samples=41] {0.5*(cos(x)*cos(y)-1)};
  \draw[white,thick] (axis cs:0,0,0) -- (axis cs:360,72,0);
  \draw[white,thick] (axis cs:0,72,0) -- (axis cs:360,144,0);
  \draw[white,thick] (axis cs:0,144,0) -- (axis cs:360,216,0);
  \draw[white,thick] (axis cs:0,216,0) -- (axis cs:360,288,0);
  \draw[white,thick] (axis cs:0,288,0) -- (axis cs:360,360,0);
  \fill[red] (axis cs:0,180,0) circle (3pt);
  \fill[red] (axis cs:360,180,0) circle (3pt);
  \fill[red] (axis cs:180,0,0) circle (3pt);
  \fill[red] (axis cs:180,360,0) circle (3pt);
  \fill[white] (axis cs:7.8,145.6,0) circle (2.5pt);
  \node[white,font=\tiny,anchor=west] at (axis cs:12,138,0) {best on the curve};
\end{axis}
\end{tikzpicture}
\end{center}
\ona{\caption{Continuous MaxCut for a single edge, $\mu(\phi_1,\phi_2) = \frac12(\cos\phi_1\cos\phi_2-1)$, with its minima $-1$ in red, and, in white, the curve $t\mapsto(E_1t,E_2t)$ mod $2\pi$ of a single layer with the clock energies $E_j = 2\pi/5^j$, for $0\le t<25$. The best point on the curve has $\mu\approx-0.91$, and rounding it as in Lemma~\ref{lem:ch05:cmc} gives $(0,\pi)$, the cut.}\label{fig:ch05:torus}\end{figure}}
\onp{\footnotesize For one edge and $m = 5$, the curve of a single parameter (white) passes near every point; its best point, with $\mu\approx-0.91$, rounds to the cut.}
\end{frame}

\begin{frame}\ft{One parameter and many local minima}
\ona{Figure~\ref{fig:ch05:onelayer} shows $\mu(\boldsymbol Et)$ for the path on three vertices, whose maximum cut is $2$, with the clock energies $E_j = 2\pi/6^j$. Over one turn of the slowest hand, $0\le t<216$, the curve has $42$ local minima, most of them far from $-2$. Its two best points, at $t\approx18$ and $t\approx198$, have $\mu\approx-1.93$ and round to the maximum cut. The data were produced by \texttt{code/ch05\_figures.py}.}
\ona{\gdef\olH{0.36\linewidth}\gdef\olW{0.7\linewidth}}\onp{\gdef\olH{0.3\linewidth}\gdef\olW{0.62\linewidth}}
\ona{\begin{figure}[htb]}\begin{center}
\begin{tikzpicture}
\begin{axis}[width=\olW,height=\olH,clip=false,axis lines=left,xmin=0,xmax=216,ymin=-2.2,ymax=0.1,xtick={0,36,72,108,144,180,216},ytick={-2,-1.5,-1,-0.5,0},xlabel={$t$},ylabel={$\mu(\boldsymbol Et)$},no markers]
  \addplot[blue!70!black,thin] table[x=t,y=mu] {figures/ch05/onelayer_curve.dat};
  \addplot[gray,dashed,thick,domain=0:216,samples=2] {-2};
  \fill[red!70!black] (axis cs:198.0,-1.933) circle (2pt);
  \fill[red!70!black] (axis cs:18.0,-1.933) circle (2pt);
  \node[note,anchor=south west] (n1) at (axis cs:0,0.25) {the fastest hand turns\\once in $6$ units of $t$};
  \draw[notearrow] (n1.south) -- (axis cs:36,-0.22);
  \node[note,anchor=south west] (n2) at (axis cs:112,0.25) {$42$ local minima, most\\of them far from $-2$};
  \draw[notearrow] (n2.south) to[bend left=10] (axis cs:96,-0.99);
  \node[note,anchor=west] (n3) at (axis cs:226,-1.45) {the best value, $-1.93$,\\rounds to the maximum cut};
  \draw[notearrow] (n3.west) to[bend left=10] (axis cs:199.5,-1.91);
  \node[note,anchor=west] (n4) at (axis cs:226,-2.1) {$-\mathrm{MaxCut}(P_3) = -2$};
  \draw[notearrow] (n4.west) -- (axis cs:217,-2.0);
\end{axis}
\end{tikzpicture}
\end{center}
\ona{\caption{Continuous MaxCut along the single-layer curve $t\mapsto\boldsymbol Et$ for the path on three vertices, with $E_j = 2\pi/6^j$, over $0\le t<6^3$. The curve is symmetric about $t = 108$. Its best points, in red, round to the maximum cut; the dashed line is $-\mathrm{MaxCut} = -2$.}\label{fig:ch05:onelayer}\end{figure}}
\onp{\footnotesize A single parameter, and $42$ local minima, since the frequencies of the clock span $d$ orders of magnitude.}
\end{frame}

\begin{frame}\ft{Step 5: QAOA with one layer}
\onp{\scriptsize}
\begin{thm}[\citet{bittel2021training}]\label{thm:ch05:qaoa1}
\ona{Problem~\ref{p:ch05:qaoa}}\onp{QAOA minimisation} is $\mathbf{NP}$-hard with a single layer, $L = 1$.
\end{thm}
\ona{\emph{Proof.} QAOA differs from step 4 in that the initial state must be the ground state of $H_b$, and the observable must be $H_c$. Both requirements are met by adding a single level. On $\C^{2d+1}$, let}
\begin{equation}
    H_b = \mathrm{diag}(E_1,-E_1,\dots,E_d,-E_d,-1), \qquad H_c = O\oplus0 + \tau\big(\ket{+_{2d}}\!\bra{2d+1} + \ket{2d+1}\!\bra{+_{2d}}\big),
\end{equation}
\ona{with $O$ from step 3, the energies of Lemma~\ref{lem:ch05:clock} with $m\ge7$, so that $|E_j|<1$, and a coupling $\tau>0$. The ground state of $H_b$ is the additional level $\ket{2d+1}$. Since $O\ket{+_{2d}} = 0$, the uniform vector lying in the kernel of a Laplacian, $H_c$ acts on the span of $\ket{2d+1}$ and $\ket{+_{2d}}$ as $\tau$ times a Pauli $X$, and}
\begin{equation}
    U_c(\gamma)\ket{2d+1} = \cos(\tau\gamma)\ket{2d+1} - i\sin(\tau\gamma)\ket{+_{2d}}.
\end{equation}
\ona{The mixer $U_b(\beta)$ then imposes the phases $e^{\mp iE_j\beta}$ on the pairs, as does the single layer of step 4, and the phase $e^{i\beta}$ on the additional level. A short computation (Exercise~\ref{exe:ch05:qaoa}) gives}
\begin{equation}
    \braket{H_c(\beta,\gamma)} = \sin^2(\tau\gamma)\,\mu(\boldsymbol E\beta) + \tau\sin(2\tau\gamma)\,g(\beta), \qquad g(\beta) = -\frac{\sin\beta}d\sum_{j=1}^d\cos(E_j\beta),
\end{equation}
\ona{where $|g|\le1$. Since $\mu\le0$, every $(\beta,\gamma)$ has $\braket{H_c}\ge\mu(\boldsymbol E\beta) - \tau$, while $\gamma = \pi/(2\tau)$ gives $\braket{H_c} = \mu(\boldsymbol E\beta)$ exactly. Angles within $\eta$ of the optimum therefore yield a $\beta$ with $\mu(\boldsymbol E\beta)\le-\mathrm{MC} + 4\pi|E|/m + \tau + \eta$, and with $\tau = 1/4$, $\eta<1/4$ and $m>8\pi|E|$ the rounding of step 4 recovers a maximum cut. \qed}
\onp{A single additional level, the ground state of $H_b$, is weakly coupled to $\ket{+_{2d}}$, and $\gamma = \pi/2\tau$ carries the state into the construction of step 4.}
\end{frame}

\begin{frame}\ft{The construction for QAOA with one layer}
\ona{Figure~\ref{fig:ch05:qaoa1} shows the two ingredients of step 5. On the left is the spectrum of $H_b$, consisting of the pairs $\pm E_j$ of the clock, all inside $(-1,1)$, and below them the additional level, which is the ground state and hence the initial state. On the right is the plane spanned by the additional level and $\ket{+_{2d}}$, which $H_c$ leaves invariant. In this plane $U_c(\gamma)$ rotates the initial state by the angle $\tau\gamma$, and a quarter turn, $\gamma = \pi/(2\tau)$, moves all the amplitude into the uniform state on the pairs, on which $U_b(\beta)$ then acts as the single layer of step 4.}
\ona{\begin{figure}[htb]}\begin{center}
\begin{tikzpicture}[every node/.style={font=\scriptsize}]
  \foreach \e/\l in {0.85/1,0.45/2,0.18/3} {
    \draw[thick,blue!70!black] (0,1.8*\e) -- (1.1,1.8*\e) node[right] {$+E_\l$};
    \draw[thick,blue!70!black] (0,-1.8*\e) -- (1.1,-1.8*\e) node[right] {$-E_\l$};
  }
  \draw[very thick,red!70!black] (0,-1.8) -- (1.1,-1.8) node[right] {$-1$};
  \draw[decorate,decoration={brace,amplitude=4pt}] (-0.15,-1.6) -- (-0.15,1.6) node[midway,left=6pt,align=center] {$\ket{+_{2d}}$: uniform\\over the pairs};
  \draw[{Latex}-{Latex},thick,orange!80!black] (-0.05,-1.8) to[bend left=45] node[left] {$\tau$} (-1.0,-0.45);
  \node[note,anchor=north] (na) at (0.6,-2.0) {the ground state of $H_b$,\\the additional level $\ket{2d+1}$};
  \draw[notearrow] (na.north) -- (0.6,-1.83);
  \node[note,anchor=south] (nb) at (0.6,1.75) {the clock energies,\\all inside $(-1,1)$};
  \draw[notearrow] (nb.south) -- (0.6,1.56);
  \begin{scope}[shift={(4.4,-1.8)}]
    \draw[-{Latex}] (0,0) -- (3.1,0) node[below] {$\ket{2d+1}$};
    \draw[-{Latex}] (0,0) -- (0,3.1) node[left] {$-i\ket{+_{2d}}$};
    \draw[gray,dashed] (2.6,0) arc[start angle=0,end angle=90,radius=2.6];
    \draw[annot,thick] (0.8,0) arc[start angle=0,end angle=35,radius=0.8];
    \node[annot] at (17:1.08) {$\tau\gamma$};
    \draw[-{Latex},very thick,blue!70!black] (0,0) -- (35:2.6);
    \node[blue!70!black,anchor=south west] at (35:2.6) {$U_c(\gamma)\ket{2d+1}$};
    \fill[red!70!black] (0,2.6) circle (2pt);
    \node[note,anchor=west] (nq) at (0.25,3.05) {at $\gamma = \pi/2\tau$, a quarter turn,\\all the amplitude is in $\ket{+_{2d}}$};
    \draw[notearrow] (nq.west) to[bend right=20] (0.07,2.66);
    \node[note,anchor=north] (nu) at (1.3,-0.5) {then $U_b(\beta)$ imprints $e^{\mp iE_j\beta}$ on the pairs};
  \end{scope}
\end{tikzpicture}
\end{center}
\ona{\caption{Step 5. On the left, the spectrum of $H_b$, not to scale, with the additional level $\ket{2d+1}$ at $-1$ as its ground state; $H_c$ couples it with strength $\tau$ to the uniform state on the $2d$ levels of the pairs, which $O$ annihilates. On the right, the span of $\ket{2d+1}$ and $\ket{+_{2d}}$, on which $H_c$ acts as $\tau X$ and $U_c(\gamma)$ as a rotation by $\tau\gamma$.}\label{fig:ch05:qaoa1}\end{figure}}
\onp{\footnotesize A single additional level is the ground state of $H_b$, and a quarter turn of $U_c$ carries it to $\ket{+_{2d}}$.}
\end{frame}

\begin{frame}\ft{Step 6: QAOA with bounded energies}
\onp{\scriptsize}
\ona{The energies $2\pi/m^j$ of step 5 span $d$ orders of magnitude, and a device implementing them would require $d\log m$ bits of precision. With $d$ layers, small integer spectra suffice.}
\begin{thm}[\citet{bittel2021training}]\label{thm:ch05:qaoad}
\ona{Problem~\ref{p:ch05:qaoa}}\onp{QAOA minimisation} is $\mathbf{NP}$-hard with $L = d$ layers and angles in $[0,2\pi)$, even for $H_b$ and $H_c$ with integer eigenvalues, $\|H_b\| = 3$ and $\|H_c\| = 1$, and with $\braket{H_c}_{\min} = 1 - 2\,\mathrm{MC}/|E|$.
\end{thm}
\ona{\emph{Idea of the construction.} The Hilbert space is the direct sum of $2d+1$ copies of $\C^d\otimes\C^d\otimes\C^2\otimes\C^2$, with basis states $\ket{i,j,a,b}$ labelled by an ordered pair of vertices and two bits. The ground state of $H_b$, which is the initial state, is the uniform superposition over the edges $ij$ and the bits, in the first copy. Each half-layer moves amplitude one copy further, $U_c(\gamma_k)$ with amplitude $\sin(\gamma_k/2)$ and $U_b(\beta_k)$ with amplitude $\sin\beta_k$; the latter also attaches the phase $e^{ia(\beta_k+\pi/2)}$ to the components with $i = k$, and $e^{(-1)^bia(\beta_k+\pi/2)}$ to those with $j = k$. After all $d$ layers, the last copy carries the phases of all the $\beta_k$, and there $H_c$ measures their interference. Amplitude left behind in the other copies can only add energy, since $H_c$ is positive semidefinite there. The result is}
\begin{equation}
    \braket{H_c(\boldsymbol\beta,\boldsymbol\gamma)} \;\ge\; \prod_{k=1}^d\sin^2\beta_k\,\sin^2\frac{\gamma_k}2\;\cdot\;\frac1{2|E|}\sum_{i,j}A_{ij}\sin\beta_i\sin\beta_j,
\end{equation}
\ona{with equality for $\gamma_k = \pi$ and $\boldsymbol\beta\in\{\pi/2,3\pi/2\}^d$. Since $\sum_{ij}A_{ij}\sin\beta_i\sin\beta_j = 4\mu(\boldsymbol\beta-\pi/2) + 2|E|$, the minimum is $1 - 2\,\mathrm{MC}/|E|$. From any angles, a cut of value at least $\frac{|E|}2\big(1-\braket{H_c}\big)$ is obtained by rounding $\boldsymbol\beta-\pi/2$ as in Lemma~\ref{lem:ch05:cmc}, or, if that is worse, by greedy local search, which cuts at least $|E|/2$ edges.}
\onp{There are $2d+1$ copies of a register; each half-layer moves amplitude one copy further and imposes the phase of $\beta_k$ on vertex $k$, and the energy is read in the last copy.}
\end{frame}

\begin{frame}\ft{The copies of the register}
\ona{Figure~\ref{fig:ch05:copies} shows the construction of step 6 as a relay. The register $\mathcal K = \C^d\otimes\C^d\otimes\C^2\otimes\C^2$ is copied $2d+1$ times, and $H_c$ and $H_b$ couple consecutive copies alternately, $H_c$ joining copies $2k-1$ and $2k$, and $H_b$ copies $2k$ and $2k+1$. Each half-layer passes amplitude one copy further, and layer $k$ also attaches the phase of $\beta_k$ to the terms that involve vertex $k$. Only in the last copy does $H_c$ measure anything that depends on the phases.}
\ona{\begin{figure}[htb]}\begin{center}
\begin{tikzpicture}[every node/.style={font=\scriptsize},cp/.style={draw,thick,rounded corners=3pt,fill=blue!6,minimum width=1.05cm,minimum height=0.8cm,align=center}]
  \node[cp] (k1) at (0,0) {$\mathcal K_1$};
  \node[cp] (k2) at (2.3,0) {$\mathcal K_2$};
  \node[cp] (k3) at (4.6,0) {$\mathcal K_3$};
  \node (dots) at (6.35,0) {$\cdots$};
  \node[cp,fill=green!12] (kl) at (8.1,0) {$\mathcal K_{2d+1}$};
  \draw[-{Latex},thick] (k1) -- node[above] {$U_c(\gamma_1)$} node[below] {$\sin\frac{\gamma_1}2$} (k2);
  \draw[-{Latex},thick] (k2) -- node[above] {$U_b(\beta_1)$} node[below] {$\sin\beta_1$} (k3);
  \draw[-{Latex},thick] (k3) -- node[above] {$U_c(\gamma_2)$} (dots);
  \draw[-{Latex},thick] (dots) -- node[above] {$U_b(\beta_d)$} node[below] {$\sin\beta_d$} (kl);
  \node[note,anchor=south] (n1) at (0,1.15) {initially uniform over\\edges $ij$ and bits $a,b$};
  \draw[notearrow] (n1.south) -- (k1.north);
  \node[note,anchor=south] (n2) at (3.45,1.15) {phase $e^{\pm ia(\beta_1+\pi/2)}$ on\\the terms with $i = 1$ or $j = 1$};
  \draw[notearrow] (n2.south) to[bend left=15] (3.95,0.03);
  \node[note,anchor=south] (n4) at (7.6,1.15) {only here does $H_c$ read\\the phases, by interference};
  \draw[notearrow] (n4.south) to[bend left=10] (kl.north);
  \node[note,anchor=north] (n3) at (2.6,-1.0) {amplitude left in the earlier copies can\\only add energy, since $H_c\succeq0$ there};
  \draw[notearrow] (n3.north) to[bend right=10] (k3.south);
  \node[note,anchor=north] (n5) at (7.6,-1.0) {there $\braket{H_c}\propto\sum_{ij}A_{ij}\sin\beta_i\sin\beta_j$,\\an affine image of $\mu(\boldsymbol\beta-\frac\pi2)$};
  \draw[notearrow] (n5.north) to[bend right=10] (kl.south);
\end{tikzpicture}
\end{center}
\ona{\caption{Step 6 as a relay through $2d+1$ copies of a register. The initial state, the ground state of $H_b$, lies in the first copy. Each of $U_c(\gamma_k)$ and $U_b(\beta_k)$ moves amplitude one copy further, by the factors shown, and $U_b(\beta_k)$ also attaches the phase of $\beta_k$ to the terms that involve vertex $k$. The energy depends on the phases only through the last copy, and is least when all the amplitude arrives there, at $\gamma_k = \pi$ and $\beta_k\in\{\pi/2,3\pi/2\}$.}\label{fig:ch05:copies}\end{figure}}
\onp{\footnotesize Amplitude is relayed through $2d+1$ copies, layer $k$ writes $\beta_k$ onto vertex $k$, and the last copy is read.}
\end{frame}

\begin{frame}\ft{A lower bound on the optimization error of QAOA}
\onp{\footnotesize}
\begin{cor}[\citet{bittel2021training}]\label{cor:ch05:qaoa}
Unless $\mathbf P = \mathbf{NP}$, every polynomial-time algorithm for \ona{Problem~\ref{p:ch05:qaoa}}\onp{QAOA minimisation}, and hence for \ona{Problem~\ref{p:ch05:vqa}}\onp{VQA minimisation}, has optimization error $\Delta\ge(1-\alpha_{\max})/2$.
\end{cor}
\ona{\emph{Proof.} On the instances of Theorem~\ref{thm:ch05:qaoad}, $w(H_c) = 2$ and $\braket{H_c}_{\min} = 1-2\,\mathrm{MC}/|E|$. An algorithm with optimization error $\Delta$ returns angles with $\braket{H_c}_a\le1-2\,\mathrm{MC}/|E| + 2\Delta$, from which a cut can be read off of value at least}
\begin{equation}
    \frac{|E|}2\big(1-\braket{H_c}_a\big) \;\ge\; \mathrm{MC} - \Delta|E| \;\ge\; (1-2\Delta)\,\mathrm{MC},
\end{equation}
\ona{where we have used $\mathrm{MC}\ge|E|/2$. This is a MaxCut algorithm with ratio $1-2\Delta$, and so $1-2\Delta\le\alpha_{\max}$. \qed}
\begin{itemize}\setlength\itemsep{0.2em}
    \item With the bound $\alpha_{\max}\le16/17$ of H\accent23 astad \citep{Hstad2001}, $\Delta\ge1/34\approx0.029$. Under the Unique Games Conjecture, $\alpha_{\max} = \alpha_{\rm GW}$ and $\Delta\ge0.061$.
    \item A local minimum at a locally optimal cut of value $\mathrm{MC}/2$ has $\braket{H_c}_a - \braket{H_c}_{\min} = \mathrm{MC}/|E|\ge1/2$, so that gradient methods, and other methods that may stop at local minima, have $\Delta\ge1/4$, without any assumption on $\mathbf P$ and $\mathbf{NP}$.
    \item The boosting of \ona{Corollary~\ref{cor:ch05:boost}}\onp{the oracular problem} does not apply, since a tensor power of $H_c$ is not the cost Hamiltonian of a QAOA circuit with mixer $H_b$; whether $\Delta$ can be pushed towards $1$ for QAOA is open.
\end{itemize}
\end{frame}

\begin{frame}\ft{Step 7: free fermions}
\onp{\scriptsize}
\ona{Free fermions are a second family of quantum systems that can be simulated classically, this time with an exponentially large Hilbert space. Consider $n$ fermionic modes with annihilation operators $c_1,\dots,c_n$, where $\{c_i^\dagger,c_j\} = \delta_{ij}$ and $\{c_i,c_j\} = 0$. A \emph{quadratic} operator $H = \sum_{ij}h_{ij}c_i^\dagger c_j$ is described by its Hermitian \emph{coefficient matrix} $h$, and all the relevant computations reduce to computations with $n\times n$ matrices.}
\begin{itemize}\setlength\itemsep{0.1em}
    \item The evolved operator $e^{iHt}Oe^{-iHt}$ is quadratic, with coefficient matrix $e^{iht}oe^{-iht}$.
    \item The ground state $\rho[H_0]$ of a quadratic $H_0$ fills the modes of the negative eigenvalues of $h_0$, and its \emph{correlation matrix} $\Gamma_{ij} = \Tr(c_j^\dagger c_i\,\rho)$ is the projector onto their span.
    \item Expectations are given by $\Tr(O\rho) = \sum_{ij}o_{ij}\Gamma_{ji} = \Tr(o\,\Gamma)$.
\end{itemize}
\begin{probl}[VQA minimisation, free fermions]\label{p:ch05:ff}
Given coefficient matrices $h_0,h_1,\dots,h_L,o\in\Herm(\C^n)$, find $\boldsymbol\phi\in\R^L$ minimising $\Tr\big(O\,U(\boldsymbol\phi)\,\rho[H_0]\,U(\boldsymbol\phi)^\dagger\big)$, where $U(\boldsymbol\phi) = e^{-i\phi_LH_L}\cdots e^{-i\phi_1H_1}$.
\end{probl}
\begin{thm}[\citet{bittel2021training}]\label{thm:ch05:ff}
\ona{Problem~\ref{p:ch05:ff}}\onp{The free-fermion problem} is $\mathbf{NP}$-hard, even when the initial state $\rho[H_0]$ is pure.
\end{thm}
\ona{\emph{Proof.} Take $n = 2d$ modes, $h_0 = \one - \frac2nJ$, where $J$ is the all-ones matrix, $h_j = \ket{2j-1}\!\bra{2j-1} - \ket{2j}\!\bra{2j}$ as in step 3, and $o = O$, the matrix of step 3. The spectrum of $h_0$ is $(-1,1,\dots,1)$, with eigenvector $\ket{+_{2d}}$ for $-1$, so that the ground state is pure: it is a single fermion in the mode $n^{-1/2}\sum_jc_j^\dagger$, with $\Gamma = \ket{+_{2d}}\!\bra{+_{2d}}$. The generators commute, and}
\begin{equation}
    \braket{O(\boldsymbol\phi)} = \Tr\big(o(\boldsymbol\phi)\,\Gamma\big) = \braket{+_{2d}|\,e^{i\sum_j\phi_jh_j}\,o\,e^{-i\sum_j\phi_jh_j}\,|+_{2d}} = \mu(\boldsymbol\phi),
\end{equation}
\ona{by the computation of step 3, a single fermion in $2d$ modes being a $2d$-level system. \qed{} (The published version has $h_0 = \one - J/n$, whose spectrum $(0,1,\dots,1)$ would give a degenerate ground state; with the factor $2$ one obtains the spectrum stated there.)}
\onp{A single fermion in $2d$ modes is a $2d$-level system, and the free-fermion problem therefore contains step 3.}
\end{frame}

\begin{frame}\ft{One fermion in $2d$ modes}
\ona{Figure~\ref{fig:ch05:fermion} shows why step 7 is step 3 in another guise. The Fock space of $2d$ modes has dimension $2^{2d}$, but the initial state has a single fermion, quadratic evolutions preserve the number of fermions, and the states with one fermion span a space of dimension $2d$, on which a quadratic operator acts by its coefficient matrix.}
\ona{\begin{figure}[htb]}\begin{center}
\begin{tikzpicture}[every node/.style={font=\scriptsize},md/.style={circle,draw,thick,fill=white,minimum size=0.56cm,inner sep=0pt},ph/.style={-{Stealth[length=3pt]},thick,blue!70!black}]
  \fill[blue!9,rounded corners=14pt] (-0.6,-0.5) rectangle (6.1,0.5);
  \foreach \j/\a in {1/40,2/115,3/250} {
    \draw[dashed,gray,rounded corners=4pt] ({2.2*(\j-1)-0.42},-0.4) rectangle ({2.2*(\j-1)+1.52},0.4);
    \node[md] (ma\j) at ({2.2*(\j-1)},0) {};
    \node[md] (mb\j) at ({2.2*(\j-1)+1.1},0) {};
    \draw[ph] (ma\j.center) -- ++(-\a:0.22);
    \draw[ph] (mb\j.center) -- ++(\a:0.22);
  }
  \foreach \k/\x in {1/0,2/1.1,3/2.2,4/3.3,5/4.4,6/5.5} \node[font=\tiny] at (\x,-0.68) {$c_\k$};
  \node[note,anchor=south] (n1) at (2.75,0.85) {a single fermion, spread evenly over all $2d$ modes, is the ground state\\of $h_0 = \one - \frac2nJ$, with correlation matrix $\Gamma = \ket{+_{2d}}\!\bra{+_{2d}}$};
  \draw[notearrow] (n1.south) -- (2.75,0.52);
  \node[note,anchor=north] (n2) at (2.75,-1.0) {$h_j$ acts only on the modes $2j-1$ and $2j$,\\and imposes the phases $e^{\mp i\phi_j}$};
  \draw[notearrow] (n2.north) -- (2.75,-0.42);
  \node[note,anchor=west] (n3) at (6.4,0) {the Fock space has $2^{2d}$\\dimensions, of which a single\\fermion uses only $2d$};
\end{tikzpicture}
\end{center}
\ona{\caption{Step 7 for $d = 3$. The six modes are grouped into the pairs on which the generators $h_j$ act. The shaded region represents the single fermion of the initial state, in the uniform mode $n^{-1/2}\sum_kc_k^\dagger$, and the phasors represent the phases $e^{\mp i\phi_j}$ that the evolution gives to its amplitudes, as in Fig.~\ref{fig:ch05:doubled}.}\label{fig:ch05:fermion}\end{figure}}
\onp{\footnotesize A single fermion in $2d$ modes is a $2d$-level system, and step 7 therefore contains step 3.}
\end{frame}

\mode<all>{\gdef\sumtab#1#2#3#4{\begin{tabular}{@{}>{\raggedright\arraybackslash}p{#1\linewidth}>{\raggedright\arraybackslash}p{#2\linewidth}>{\raggedright\arraybackslash}p{#3\linewidth}>{\raggedright\arraybackslash}p{#4\linewidth}@{}}\toprule
setting & size & layers & unless $\mathbf P = \mathbf{NP}$, no polynomial-time algorithm achieves \\\midrule
oracle, rotations (step 2) & $d$ qubits & $d$ & $\Delta<1$ \\
dimension $2d$ (steps 3, 4) & $\mathcal O(\log d)$ qubits & $d$, or $1$ & the optimum \\
QAOA (step 5) & $\mathcal O(\log d)$ qubits & $1$ & the optimum \\
QAOA, bounded energies (step 6) & $\mathcal O(\log d)$ qubits & $d$ & $\Delta<\frac12(1-\alpha_{\max})$, that is, $\frac1{34}$, or $0.061$ under the UGC \\
free fermions (step 7) & $2d$ modes & $d$ & the optimum \\\midrule
local methods on step 6 & & & $\Delta<\frac14$, unconditionally \\\bottomrule
\end{tabular}}}
\begin{frame}\ft{A summary of the hardness results}
\ona{Table~\ref{tab:ch05:summary} collects what the reductions exclude. In every row the energy can be computed efficiently, by a quantum computer in the first row and classically in the others, so that the obstruction lies in the search over the parameters.}
\onp{\begin{center}\scriptsize\sumtab{0.25}{0.15}{0.09}{0.38}\end{center}}
\ona{\begin{table}[htb]\centering\caption{The hardness results of \citet{bittel2021training}. ``The optimum'' means that it is $\mathbf{NP}$-hard to find parameters that are optimal up to exponential precision; $\Delta$ is the optimization error of Definition~\ref{defn:ch05:opterr}.}\label{tab:ch05:summary}\footnotesize\sumtab{0.25}{0.15}{0.09}{0.38}\end{table}}
\end{frame}

\begin{frame}\ft{The excluded guarantees on a common scale}
\ona{Figure~\ref{fig:ch05:deltaline} shows, on a common scale, the optimization errors that cannot be guaranteed in each setting of Table~\ref{tab:ch05:summary}. For QAOA, no guarantee below $1/34$ is possible in polynomial time unless $\mathbf P = \mathbf{NP}$, and none below $0.061$ if the Unique Games Conjecture holds as well. Methods that may stop at a local minimum cannot guarantee less than $1/4$. For general VQAs with an oracle, nothing below the trivial value $1$ can be guaranteed at all.}
\ona{\begin{figure}[htb]}\begin{center}
\begin{tikzpicture}[x=6.5cm,every node/.style={font=\scriptsize}]
  \foreach \y/\l in {2.0/{QAOA (step 6)},1.0/{local methods},0/{oracle (step 2)}} {\draw[gray!50] (0,\y) -- (1,\y); \node[anchor=east] at (-0.02,\y) {\l};}
  \fill[red!55] (0,1.84) rectangle (1/34,2.16);
  \fill[orange!50] (1/34,1.84) rectangle (0.061,2.16);
  \fill[gray!45] (0,0.84) rectangle (0.25,1.16);
  \fill[red!55] (0,-0.16) rectangle (0.995,0.16);
  \draw[thick,-{Latex}] (0,-0.5) -- (1.07,-0.5) node[right] {$\Delta$};
  \foreach \x in {0,0.25,0.5,0.75,1} \draw (\x,-0.45) -- (\x,-0.55) node[below] {$\x$};
  \node[note,anchor=west] (n1) at (0.13,2.3) {in red, below $\frac1{34}$, excluded unless $\mathbf P = \mathbf{NP}$ by the $\frac{16}{17}$ of H\accent23 astad};
  \draw[notearrow] (n1.west) to[bend right=10] (0.016,2.14);
  \node[note,anchor=west] (n2) at (0.13,1.7) {in orange, below $0.061$, excluded under the Unique Games Conjecture};
  \draw[notearrow] (n2.west) to[bend left=10] (0.046,1.86);
  \node[note,anchor=west] (n3) at (0.29,1.0) {the local minima include all local optima of single-vertex\\search, so that $\Delta\ge\frac14$, without any assumption};
  \draw[notearrow] (n3.west) -- (0.255,1.0);
  \node[note,anchor=west] (n4) at (1.02,0.0) {nothing below $1$,\\the trivial value};
  \draw[notearrow] (n4.west) -- (0.997,0.0);
\end{tikzpicture}
\end{center}
\ona{\caption{The optimization errors $\Delta$ that cannot be guaranteed, shaded, in three settings: QAOA with bounded energies (Corollary~\ref{cor:ch05:qaoa}), methods that may stop at local minima on the same instances, and general VQAs with an oracle (Corollary~\ref{cor:ch05:boost}). The unshaded values are not excluded by these results, and whether they can be achieved is open.}\label{fig:ch05:deltaline}\end{figure}}
\onp{\footnotesize Excluded are $\Delta<\frac1{34}$, or $0.061$ under the UGC, for QAOA; $\Delta<\frac14$ for local methods; and $\Delta<1$ with an oracle.}
\end{frame}

\begin{frame}\ft{Training, if the Unique Games Theorem holds}
\onp{\scriptsize}
\ona{The bound of Corollary~\ref{cor:ch05:qaoa} is stated in terms of $\alpha_{\max}$, the best ratio achievable for MaxCut in polynomial time, and it is therefore sharpened if the Unique Games Theorem claimed in \Chap~\chref{C.Inapprox} is correct. Two cases must be distinguished, according to the computer on which the training is carried out.}
\begin{itemize}\setlength\itemsep{0.3em}
    \item Classical training. Goemans--Williamson achieves $\alpha_{\rm GW}$, and if the theorem holds and $\mathbf P\ne\mathbf{NP}$, no polynomial-time algorithm achieves more on every instance, so that $\alpha_{\max} = \alpha_{\rm GW}$. Every polynomial-time classical training procedure for the QAOA instances of \ona{Theorem~\ref{thm:ch05:qaoad}}\onp{step 6} then has optimization error
    \begin{equation*}
        \Delta \;\ge\; \tfrac12(1-\alpha_{\rm GW}) \;\approx\; 0.061,
    \end{equation*}
    with no assumption beyond $\mathbf P\ne\mathbf{NP}$. Without the theorem, the unconditional bound is $1/34 \approx 0.029$, from the $16/17$ of H\accent23 astad.
    \item Training with a quantum computer in the loop. An optimizer that queries a quantum device, or runs any other quantum subroutine, is a polynomial-time quantum algorithm. If it guaranteed $\Delta < (1-\alpha_{\rm GW})/2$, the reduction would turn it into a polynomial-time quantum algorithm approximating MaxCut beyond $\alpha_{\rm GW}$. By the theorem, such an algorithm would decide 3-SAT, so that the same bound holds unless $\mathbf{NP}\subseteq\mathbf{BQP}$. In any case the quantum device does not help with the hard instances, which live on $\mathcal O(\log d)$ qubits and can be simulated classically; the difficulty lies in the search over the parameters.
    \item Some results are unaffected. The oracle setting already excludes any guarantee $\Delta<1$, assuming only $\mathbf P\ne\mathbf{NP}$, or $\mathbf{NP}\not\subseteq\mathbf{BQP}$ for quantum optimizers. Local methods cannot guarantee $\Delta<1/4$ in any case. The $\mathbf{QCMA}$-hardness of approximating the optimal depth \citep{bittelQCMA} does not involve unique games.
\end{itemize}
\ona{If the theorem holds, then, an error of about $6\%$ of the spectral width is the best that any efficient training procedure, classical or quantum, can guarantee on all instances of QAOA, unless $\mathbf{NP}\subseteq\mathbf{BQP}$. As before, this says nothing about typical instances.}
\end{frame}

\begin{frame}\ft{The scope of the hardness results}
\begin{itemize}\setlength\itemsep{0.3em}
    \item In the worst case, unless $\mathbf P = \mathbf{NP}$, no polynomial-time training procedure finds the optimal parameters, nor even parameters within a fixed fraction of the spectral width of the optimum. The difficulty lies in the classical optimization, and not in the quantum hardware or in the estimation of expectation values.
    \item The results do not imply that good parameters are hard to find for typical instances, and \Chap~\chref{C.PerIteration} finds that often they are not.
    \item Nor do they imply that local minima are useless: a warm start which inherits a classical guarantee retains it, whatever the optimizer (\Chaps~\chref{C.WSRounded}--\chref{C.WSContinuous}).
\end{itemize}
\ona{Variational methods are in good company. The optimization in the Hartree--Fock method is $\mathbf{NP}$-hard \citep{schuch2009computational}, as are global variations in the density matrix renormalisation group \citep{eisert2006computational}, and both are in daily use. Later work has extended these results for VQAs. Approximating the depth that an ansatz requires to reach a target energy is $\mathbf{QCMA}$-hard \citep{bittelQCMA}; in an idealised digitised setting the training problem is even undecidable \citep{korpas2025undecidable}; and for a constant number of angles there are tractable regimes \citep{sakos2026global}.}
\onp{\footnotesize Also $\mathbf{NP}$-hard, and in daily use, are Hartree--Fock \citep{schuch2009computational} and DMRG \citep{eisert2006computational}. Later results: approximating the depth is $\mathbf{QCMA}$-hard \citep{bittelQCMA}, training is undecidable in an idealised setting \citep{korpas2025undecidable}, and there are tractable regimes for a constant number of angles \citep{sakos2026global}.}
\end{frame}

\section{A second obstacle: barren plateaus}\label{sec:ch05:barren}

\begin{frame}\ft{Barren plateaus}\label{fr:ch05:barren}
\ona{Persistent local minima are one obstacle to training. A second, of a quite different kind, is the \emph{barren plateau} \citep{Clean_2018_BarrenPlateaus}: for sufficiently random and deep circuits that do not reflect the structure of the problem, the gradient of $C$ has mean zero and a variance that decays exponentially in the number of qubits,}
\begin{equation}
    \mathrm{Var}\Big[\frac{\partial C}{\partial\vartheta_i}\Big] \in \mathcal O(2^{-n}),
\end{equation}
\ona{so that exponentially many shots are needed to detect any gradient at all. Several causes have since been identified, among them circuits close to unitary 2-designs, excessive entanglement combined with measurements of part of the system, global cost functions, and noise; they have been unified in terms of the dimension of the dynamical Lie algebra of the circuit. The landscapes of \citet{bittel2021training} have no barren plateaus, since they live on few qubits; their minima are hard to find for combinatorial reasons (Fig.~\ref{fig:ch05:obstacles}).}
\begin{itemize}\setlength\itemsep{0.2em}
    \item Ans\"atze with few parameters that reflect the structure of the problem, such as QAOA and UCC, are less affected than hardware-efficient ones.
    \item An initialisation that places the state near the optimum avoids the plateau; how to find one is the question that warm starts address (\Chaps~\chref{C.WSRounded}--\chref{C.WSContinuous}).
\end{itemize}
\end{frame}

\begin{frame}\ft{The two obstacles compared}
\ona{Figure~\ref{fig:ch05:obstacles} contrasts the two obstacles schematically. On a barren plateau the landscape is flat to exponential precision, and the optimizer has no signal to follow. In the landscapes of \citet{bittel2021training} there is a signal everywhere, but it leads to the nearest persistent local minimum.}
\ona{\gdef\obW{0.46\linewidth}\gdef\obH{0.38\linewidth}}\onp{\gdef\obW{0.46\linewidth}\gdef\obH{0.4\linewidth}}
\ona{\begin{figure}[htb]}\begin{center}
\begin{tikzpicture}
\begin{axis}[name=bp,width=\obW,height=\obH,clip=false,axis lines=left,xtick=\empty,ytick=\empty,xlabel={parameter},ylabel={cost},xmin=0,xmax=1,ymin=-1.15,ymax=2.4,samples=500,domain=0:1,no markers,thick,title={\scriptsize (a) a barren plateau}]
  \addplot[blue!70!black] {-exp(-((x-0.62)/0.025)^2)+0.006*sin(deg(37*x))};
  \node[note,anchor=north west] (p1) at (axis cs:0.02,2.3) {a gradient of order $2^{-n}$\\gives no signal to follow};
  \draw[notearrow] (p1.south) to[bend right=10] (axis cs:0.3,0.04);
  \node[note,anchor=north east] (p2) at (axis cs:0.99,1.25) {the minimum, in\\a narrow gorge};
  \draw[notearrow] (p2.south) to[bend left=15] (axis cs:0.635,-0.95);
\end{axis}
\begin{axis}[at={(bp.east)},anchor=west,xshift=1.2cm,width=\obW,height=\obH,clip=false,axis lines=left,xtick=\empty,ytick=\empty,xlabel={parameter},ylabel={cost},xmin=0,xmax=1,ymin=-1.15,ymax=2.4,samples=500,domain=0:1,no markers,thick,title={\scriptsize (b) persistent local minima}]
  \addplot[red!70!black] {-0.6*cos(deg(2*pi*(x-0.62)))-0.3*cos(deg(10*pi*x+1))-0.15*cos(deg(26*pi*x+2))};
  \fill[orange!80!black] (axis cs:0.365,-0.418) circle (2pt);
  \fill[green!40!black] (axis cs:0.5865,-0.979) circle (2pt);
  \node[note,anchor=north west] (q2) at (axis cs:0.02,2.3) {a signal everywhere, but it\\leads to the nearest local\\minimum, a local optimum};
  \draw[notearrow] (q2.south) to[bend right=15] (axis cs:0.36,-0.38);
  \node[note,anchor=north east] (q3) at (axis cs:0.99,1.3) {the global\\minimum};
  \draw[notearrow] (q3.south) to[bend left=10] (axis cs:0.6,-0.95);
\end{axis}
\end{tikzpicture}
\end{center}
\ona{\caption{The two obstacles to training, schematically. (a) On a barren plateau the gradient is exponentially small almost everywhere, and the minimum lies in a narrow gorge. (b) In the landscapes of \citet{bittel2021training} the gradient is informative, but local methods stop at the nearest of many persistent local minima.}\label{fig:ch05:obstacles}\end{figure}}
\onp{\footnotesize (a) There is no signal. (b) There is a signal, but it leads to the wrong place.}
\end{frame}

\begin{frame}\ft{Questions for the rest of the course}
\begin{enumerate}\setlength\itemsep{0.3em}
    \item Which problems can be encoded, and with which ans\"atze? Combinatorial optimization is treated in \Chap~\chref{C.QAOA}, linear systems in \Chap~\chref{C.VQLS}, and portfolios in \Chap~\chref{C.Portfolio}.
    \item How are gradients computed? See \Chap~\chref{C.PSR}.
    \item How many iterations are needed, given noisy and biased measurements? See \Chap~\chref{C.IterationComplexity}.
    \item What does a single iteration cost? See \Chap~\chref{C.PerIteration}.
    \item Where should one start, given that the landscape has persistent local minima? See \Chaps~\chref{C.WSRounded}--\chref{C.WSContinuous}.
\end{enumerate}
\end{frame}

\ona{\section*{Exercises}}

\begin{frame}\ft{Exercises}
\onp{\scriptsize}
\begin{exe}[A first VQA]
Consider the problem of finding the eigenstate of lowest energy of the spin-chain Hamiltonian
\begin{align*}
    H = J \sum_{i=1}^{n-1} Z_iZ_{i+1} + h \sum_{i=1}^n X_i,
\end{align*}
with $J<0$, the \emph{transverse-field Ising model}. Apply the variational principle to the product ansatz
\begin{align}
    \ket{\psi(\theta,\phi)} = \big(\cos(\theta)\ket{0} + e^{i \phi}\sin(\theta)\ket{1}\big)^{\otimes n}.
\end{align}
Then carry out a small experiment on at least three qubits, with small $h$, and compare the energy obtained with the exact ground energy, which may be found by means of the Jordan--Wigner transformation. This is a first variational quantum algorithm.
\end{exe}
\end{frame}

\begin{frame}\ft{Exercises, continued}
\onp{\footnotesize}
\begin{exe}[The error measure]
Show that the quantities $\delta$, $\delta_m$ and $\delta_o$ of \eqref{eq:ch05:split} lie in $[0,1]$ and are invariant under $O\mapsto aO+b\one$ for $a>0$. Compute the spectral widths of the Ising observable of step 2 and of the observable of step 3, for the 4-cycle and for the triangle.
\end{exe}
\begin{exe}[Continuous MaxCut]
For the triangle, write out $\mu(\phi_1,\phi_2,\phi_3)$ from \eqref{eq:ch05:mu}, find its minimum, and show that the set of minimisers is infinite. Does the rounding of the proof still reach a maximum cut?
\end{exe}
\begin{exe}[Persistent local minima]\label{exe:ch05:persistent}
Let $V = L_1\cup L_2\cup R_1\cup R_2$ with four parts of $k$ vertices each, and join every vertex of $L_1\cup L_2$ to every vertex of $R_1\cup R_2$, except for a perfect matching between $L_1$ and $R_1$ and one between $L_2$ and $R_2$. Show that the cut $(L_1\cup R_1, L_2\cup R_2)$ corresponds to a strict local minimum of $\mu$, and that its ratio to the maximum cut is $k/(2k-1)\to1/2$.
\end{exe}
\end{frame}

\begin{frame}\ft{Exercises, continued}
\onp{\footnotesize}
\begin{exe}[Boosting]
Verify that the eigenvalues of $(-1)^{k-1}O^{\otimes k}$ lie in $[-\mathrm{MC}^k,0]$. Find the smallest $k$ for which an algorithm with $\Delta\le\rho$ would exceed the ratio $16/17$ of H\accent23 astad, for $\rho = 0.9$ and for $\rho = 0.99$.
\end{exe}
\begin{exe}[A clock]\label{exe:ch05:clock}
Let $m = 10$, $d = 3$ and $\boldsymbol\phi = 2\pi\,(0.37, 0.81, 0.25)$. Compute the digits $s_j$ and the time $t$ of \ona{Lemma~\ref{lem:ch05:clock}}\onp{the clock lemma}, and the errors $|\phi_j - E_jt|_{2\pi}$, and compare them with $4\pi/m$. Explain why the error of hand $j$ depends only on $\phi_j$ and on the digits $s_k$ with $k<j$.
\end{exe}
\end{frame}

\begin{frame}\ft{Exercises, continued}
\onp{\footnotesize}
\begin{exe}[QAOA with one layer]\label{exe:ch05:qaoa}
In the construction of step 5, show that $O\ket{+_{2d}} = 0$, compute $U_b(\beta)U_c(\gamma)\ket{2d+1}$, and derive the formula for $\braket{H_c(\beta,\gamma)}$. Show that $|g(\beta)|\le1$.
\end{exe}
\begin{exe}[Free fermions]
Show that the quadratic Hamiltonian with $h_0 = \one - \frac2nJ$ has a unique ground state, with one fermion and correlation matrix $J/n$, and that with $h_0 = \one - \frac1nJ$ the ground state is degenerate. In the latter case, find the correlation matrix of the maximally mixed state on the ground space.
\end{exe}
\end{frame}

\begin{frame}\ft{Exercises, continued}
\onp{\footnotesize}
\begin{exe}[The worked example]\label{exe:ch05:c4}
For the 4-cycle, show that $F(\beta,\gamma) = -2 + \sin4\beta\,\sin2\gamma$. (By symmetry, $F = -2 + 2\braket{Z_1Z_2}$; conjugate $Z_1Z_2$ first by $e^{-i\beta H_M}$ and then by $e^{-i\gamma H_C}$, and take the expectation in $\ket+^{\otimes4}$.)
\end{exe}
\begin{exe}[Two ways to build a $ZZ$ rotation]\label{exe:ch05:rzz}
Verify both identities in \eqref{eq:ch05:rzz} by multiplying out the $4\times4$ matrices. Then construct $e^{-i\theta Z_1Z_2Z_3/2}$ from CNOTs and a single $R_Z$, and count the CNOTs.
\end{exe}
\end{frame}

\begin{frame}\ft{Exercises, continued}
\onp{\footnotesize}
\begin{exe}[Trotter error]
For $A = aX$ and $B = bZ$, expand $e^{-i(A+B)} - e^{-iA}e^{-iB}$ to second order in $a$ and $b$, and compare its norm with the bound $\frac12\|[A,B]\| = |ab|$ of \ona{Lemma~\ref{lem:ch05:trotter}}\onp{the Trotter lemma}.
\end{exe}
\begin{exe}[H$_2$ in one layer]
On the span of $\ket{\rm HF}$ and the doubly excited determinant, write $F$ and $H$ as $2\times2$ matrices, and show, on the Bloch sphere, that $e^{-i\beta F}e^{-i\gamma H}\ket{\rm HF}$ is the ground state of $H$ for suitable $\beta$ and $\gamma$. Explain why $e^{-i\gamma H}\ket{\rm HF}$ alone is never better than $\ket{\rm HF}$.
\end{exe}
\end{frame}

\begin{frame}[allowframebreaks]\ft{References}
\biblio
\end{frame}

\end{document}
